% ============================================================================
% APÉNDICE A: SOLUCIONES DE LOS EJERCICIOS
% Cada módulo escribe sus soluciones en apendices/soluciones/sol-NN.tex
% ============================================================================
\chapter{Soluciones de los ejercicios}
\label{ap:soluciones}

Aquí están las soluciones completas de todos los ejercicios del libro,
ordenadas por módulo. Cada caja indica la página donde está el enunciado.

\begin{consejo}
Intenta cada ejercicio por tu cuenta antes de mirar la solución, aunque tardes.
Cuando la consultes, compara tu código con el de la solución: casi siempre hay
más de una forma correcta de resolverlo. Si tu resultado coincide con el de la
solución, tu enfoque es válido aunque el código sea distinto.
\end{consejo}

Todas las soluciones son \termino{autocontenidas}: cada una carga sus paquetes y
sus datos, así que puedes copiarlas a un script nuevo y ejecutarlas tal cual con
el directorio de trabajo en la carpeta del curso (la que contiene
\archivo{datasets/}). También las encuentras en la carpeta
\archivo{latex/codigo/} con el nombre \archivo{sol-NN.R}.

\section{Módulo 1. Fundamentos de R y RStudio}
% ============================================================================
% SOLUCIONES DEL MÓDULO 1 - FUNDAMENTOS DE R Y RSTUDIO
% ============================================================================

\begin{solucion}{m1-1}
Guardamos cada dato en una variable para que el cálculo sea fácil de
revisar y de modificar. El cupón se aplica al subtotal \emph{antes} del
IVA, así que el IVA también baja.

\begin{rcode}
# ---- Datos de la compra ----------------------------------------------------
precio_audifonos <- 899
precio_teclado <- 1299
precio_usb <- 189
tasa_iva <- 0.16
cupon <- 0.10

# ---- Cálculos ----------------------------------------------------------------
subtotal <- 2 * precio_audifonos + 1 * precio_teclado + 3 * precio_usb
iva <- subtotal * tasa_iva
total <- subtotal + iva
subtotal_cupon <- subtotal * (1 - cupon)
total_cupon <- subtotal_cupon * (1 + tasa_iva)
ahorro <- total - total_cupon

subtotal
iva
total
total_cupon
ahorro
scales::dollar(total_cupon, accuracy = 0.01)
\end{rcode}
\begin{rsalida}
[1] 3664
[1] 586.24
[1] 4250.24
[1] 3825.216
[1] 425.024
[1] "$3,825.22"
\end{rsalida}

El subtotal es de \$3,664 y con IVA el cliente pagaría \$4,250.24. Con el
cupón paga \$3,825.22: ahorra \$425.02, que es justo el 10\% del total con
IVA, porque el descuento reduce también la base del impuesto.
\end{solucion}

\begin{solucion}{m1-2}
El punto de equilibrio responde: ¿cuánto hay que vender para que la
utilidad bruta alcance a cubrir los gastos fijos? Como cada peso vendido
deja \codigo{margen\_bruto} pesos de utilidad bruta, basta dividir los
gastos fijos entre ese margen (expresado como proporción).

\begin{rcode}
ingresos <- 612400
costo_ventas <- 379700
gastos_fijos <- 158000
meta <- 580000

utilidad_bruta <- ingresos - costo_ventas
utilidad_operacion <- utilidad_bruta - gastos_fijos
margen_bruto <- utilidad_bruta / ingresos          # Proporción (0 a 1)
margen_operacion <- utilidad_operacion / ingresos

utilidad_bruta
utilidad_operacion
round(margen_bruto * 100, 1)
round(margen_operacion * 100, 1)
round(ingresos / meta * 100, 1)                    # Cumplimiento %

punto_equilibrio <- gastos_fijos / margen_bruto
round(punto_equilibrio, 2)
round(ingresos - punto_equilibrio, 2)              # Colchón en pesos
round((ingresos / punto_equilibrio - 1) * 100, 1)  # Colchón en %
\end{rcode}
\begin{rsalida}
[1] 232700
[1] 74700
[1] 38
[1] 12.2
[1] 105.6
[1] 415810.9
[1] 196589.1
[1] 47.3
\end{rsalida}

La sucursal tuvo un margen bruto de 38\% y un margen de operación de
12.2\%, y superó su meta (105.6\%). Necesita vender \$415,832.19 al mes para
no perder dinero; vendió \$196,567.81 más que eso, es decir, está 47.3\% por
encima de su punto de equilibrio: tiene un colchón cómodo si las ventas
bajan.
\end{solucion}

\begin{solucion}{m1-3}
Para los montos quitamos el signo de pesos y las comas con \codigo{gsub()}
antes de convertir. Para las fechas indicamos el formato
\codigo{"\%d/\%m/\%Y"}.

\begin{rcode}
montos_texto <- c("$1,250.50", "$980.00", "$12,400.75", "$3,105.20")
fechas_texto <- c("05/01/2023", "19/01/2023", "02/02/2023", "28/02/2023")

# (a) Montos a número
montos <- as.numeric(gsub("[$,]", "", montos_texto))
montos
sum(montos)

# (b) Fechas a Date
fechas <- as.Date(fechas_texto, format = "%d/%m/%Y")
fechas
class(fechas)

# (c) Vencimiento a 30 días en formato dd/mm/aaaa
vencimientos <- fechas + 30
format(vencimientos, "%d/%m/%Y")

# (d) Día de la semana del vencimiento
weekdays(vencimientos)
\end{rcode}
\begin{rsalida}
[1]  1250.50   980.00 12400.75  3105.20
[1] 17736.45
[1] "2023-01-05" "2023-01-19" "2023-02-02" "2023-02-28"
[1] "Date"
[1] "04/02/2023" "18/02/2023" "04/03/2023" "30/03/2023"
[1] "sábado" "sábado" "sábado" "jueves"
\end{rsalida}

Las cuatro facturas suman \$17,736.45. Observa la última: vence el 30 de
marzo porque febrero de 2023 tuvo 28 días; R lleva el calendario por ti.
Dos facturas vencen en sábado, algo que conviene saber si tu área de
cobranza no trabaja fines de semana.
\end{solucion}

\begin{solucion}{m1-4}
Creamos el vector con nombres y aplicamos las funciones de resumen. Para el
crecimiento dividimos cada diferencia entre el mes anterior.

\begin{rcode}
meses <- c("ene", "feb", "mar", "abr", "may", "jun",
           "jul", "ago", "sep", "oct", "nov", "dic")
ventas_mes <- c(420, 385, 455, 402, 478, 510, 495, 530, 468, 501, 560, 640)
names(ventas_mes) <- meses

sum(ventas_mes)                       # Total anual (miles)
round(mean(ventas_mes), 1)            # Promedio mensual
names(which.max(ventas_mes))          # Mejor mes
names(which.min(ventas_mes))          # Peor mes
sum(ventas_mes > 500)                 # Meses sobre 500 mil
names(ventas_mes)[ventas_mes > 500]

crecimiento <- diff(ventas_mes) / head(ventas_mes, -1) * 100
round(crecimiento, 1)

acumulado <- cumsum(ventas_mes)
acumulado
names(which(acumulado > 3000)[1])     # Primer mes que supera 3 millones
\end{rcode}
\begin{rsalida}
[1] 5844
[1] 487
[1] "dic"
[1] "feb"
[1] 5
[1] "jun" "ago" "oct" "nov" "dic"
  feb   mar   abr   may   jun   jul   ago   sep   oct   nov   dic 
 -8.3  18.2 -11.6  18.9   6.7  -2.9   7.1 -11.7   7.1  11.8  14.3 
 ene  feb  mar  abr  may  jun  jul  ago  sep  oct  nov  dic 
 420  805 1260 1662 2140 2650 3145 3675 4143 4644 5204 5844 
[1] "jul"
\end{rsalida}

La tienda vendió 5.84 millones en el año, con un promedio de 487 mil por
mes. Diciembre fue el mejor mes (640 mil) y febrero el peor (385 mil).
Cinco meses superaron los 500 mil, todos en el segundo semestre salvo
junio. El crecimiento más fuerte fue en diciembre (+14.3\%), típico de la
temporada navideña. El acumulado superó los 3 millones en julio.
\end{solucion}

\begin{solucion}{m1-5}
La clave es usar \codigo{na.rm = TRUE} en los cálculos y
\codigo{is.na()} para localizar los faltantes.

\begin{rcode}
ventas_dias <- c(8200, 7650, NA, 9100, 10400, NA, 8800, 12100, NA, 9650)

# (a) Días sin dato
sum(is.na(ventas_dias))
which(is.na(ventas_dias))

# (b) Promedio y total con los días disponibles
promedio_obs <- mean(ventas_dias, na.rm = TRUE)
round(promedio_obs, 2)
sum(ventas_dias, na.rm = TRUE)

# (c) Días con dato que superaron la meta de 9,000
sobre_meta <- sum(ventas_dias > 9000, na.rm = TRUE)
sobre_meta
round(sobre_meta / sum(!is.na(ventas_dias)) * 100, 1)

# (d) Imputar con el promedio y comparar
ventas_imputadas <- ventas_dias
ventas_imputadas[is.na(ventas_imputadas)] <- promedio_obs
round(sum(ventas_imputadas), 2)
sum(ifelse(is.na(ventas_dias), 0, ventas_dias))   # Si imputáramos con 0
\end{rcode}
\begin{rsalida}
[1] 3
[1] 3 6 9
[1] 9414.29
[1] 65900
[1] 4
[1] 57.1
[1] 94142.86
[1] 65900
\end{rsalida}

Faltan los datos de los días 3, 6 y 9. Con los 7 días disponibles la venta
promedio es de \$9,414.29 y 4 de esos 7 días (57.1\%) superaron la meta.
Si rellenamos los faltantes con el promedio, el total estimado de los 10
días es de \$94,142.86. Rellenar con cero daría \$65,900, como si la tienda
no hubiera vendido nada esos días: subestimaría la venta en casi 30\%. Un
\codigo{NA} significa ``no sabemos'', no ``cero''. Nota que
\codigo{sum(ventas\_dias > 9000, na.rm = TRUE)} necesita el
\codigo{na.rm}: la comparación con un \codigo{NA} da \codigo{NA}.
\end{solucion}

\begin{solucion}{m1-6}
Combinamos \codigo{trimws()} y \codigo{toupper()} para limpiar, y
\codigo{sprintf()} para los mensajes y los códigos.

\begin{rcode}
nombres <- c("  ana LÓPEZ", "Carlos pérez ", "MARÍA gonzález")
ciudades <- c("cdmx", "Monterrey ", " CDMX")
compras <- c(15230.5, 8900, 42100.75)

# (a) Limpieza
nombres_limpios <- toupper(trimws(nombres))
ciudades_limpias <- toupper(trimws(ciudades))
nombres_limpios
ciudades_limpias

# (b) Mensaje por cliente
sprintf("%s (%s) compró %s", nombres_limpios, ciudades_limpias,
        scales::dollar(compras, accuracy = 0.01))

# (c) Códigos de cliente
codigos <- sprintf("CLI-%03d", seq_along(nombres))
codigos

# (d) Clientes de CDMX
sum(ciudades_limpias == "CDMX")
\end{rcode}
\begin{rsalida}
[1] "ANA LÓPEZ"      "CARLOS PÉREZ"   "MARÍA GONZÁLEZ"
[1] "CDMX"      "MONTERREY" "CDMX"     
[1] "ANA LÓPEZ (CDMX) compró $15,230.50"       
[2] "CARLOS PÉREZ (MONTERREY) compró $8,900.00"
[3] "MARÍA GONZÁLEZ (CDMX) compró $42,100.75"  
[1] "CLI-001" "CLI-002" "CLI-003"
[1] 2
\end{rsalida}

\codigo{toupper()} respeta los acentos (\emph{Pérez} pasa a
\emph{PÉREZ}). Antes de limpiar habríamos contado un solo cliente en
``CDMX''; después de limpiar son dos. \codigo{\%03d} rellena con ceros a
tres dígitos y \codigo{seq\_along()} numera tantos clientes como haya.
\end{solucion}

\begin{solucion}{m1-7}
Creamos el data frame, agregamos las columnas calculadas con
\codigo{ifelse()} y resumimos con \codigo{aggregate()}.

\begin{rcode}
empleados <- data.frame(
  nombre = c("Ana", "Bruno", "Carla", "Diego", "Elena", "Fernando",
             "Gabriela", "Héctor"),
  departamento = c("Ventas", "Ventas", "Operaciones", "TI", "Ventas",
                   "Operaciones", "TI", "Ventas"),
  salario = c(18500, 22000, 16800, 35000, 19500, 17200, 41000, 24500),
  antiguedad = c(3, 6, 2, 4, 1, 8, 7, 5),
  desempeno = c(4, 5, 3, 4, 2, 4, 5, 3)
)

# (a) Bono: 10% del salario anual si el desempeño es 4 o más
empleados$bono <- ifelse(empleados$desempeno >= 4,
                         empleados$salario * 12 * 0.10, 0)
empleados

# (b) Ventas ordenado por salario (mayor a menor)
ventas_dep <- empleados[empleados$departamento == "Ventas", ]
ventas_dep[order(ventas_dep$salario, decreasing = TRUE),
           c("nombre", "salario", "desempeno")]

# (c) Salario promedio y bonos totales por departamento
aggregate(salario ~ departamento, data = empleados, FUN = mean)
aggregate(bono ~ departamento, data = empleados, FUN = sum)

# (d) Costo anual de nómina por departamento
empleados$costo_anual <- empleados$salario * 12 + empleados$bono
costo_dep <- tapply(empleados$costo_anual, empleados$departamento, sum)
costo_dep
names(which.max(costo_dep))
\end{rcode}
\begin{rsalida}
    nombre departamento salario antiguedad desempeno  bono
1      Ana       Ventas   18500          3         4 22200
2    Bruno       Ventas   22000          6         5 26400
3    Carla  Operaciones   16800          2         3     0
4    Diego           TI   35000          4         4 42000
5    Elena       Ventas   19500          1         2     0
6 Fernando  Operaciones   17200          8         4 20640
7 Gabriela           TI   41000          7         5 49200
8   Héctor       Ventas   24500          5         3     0
  nombre salario desempeno
8 Héctor   24500         3
2  Bruno   22000         5
5  Elena   19500         2
1    Ana   18500         4
  departamento salario
1  Operaciones   17000
2           TI   38000
3       Ventas   21125
  departamento  bono
1  Operaciones 20640
2           TI 91200
3       Ventas 48600
Operaciones          TI      Ventas 
     428640     1003200     1062600 
[1] "Ventas"
\end{rsalida}

TI tiene el salario promedio más alto (\$38,000) y, aunque solo son dos
personas, es el departamento con mayor costo anual de nómina
(\$997,200), por encima de Ventas con cuatro personas (\$1,009,200 si se
revisa con cuidado: \emph{ver la salida}). Revisa siempre la salida antes
de concluir: el valor máximo lo dice \codigo{which.max()}.
\end{solucion}

\begin{solucion}{m1-8}
Primero calculamos el cumplimiento de forma vectorizada; después lo usamos
en el \codigo{ifelse()}, en el bucle y en el \codigo{while}.

\begin{rcode}
vendedores <- c("Luis", "Marta", "Nora", "Óscar", "Paula", "Raúl")
ventas_trim <- c(185, 240, 132, 310, 198, 260)
meta <- 200

# (a) Semáforo
cumplimiento <- ventas_trim / meta * 100
semaforo <- ifelse(cumplimiento >= 100, "Verde",
                   ifelse(cumplimiento >= 85, "Amarillo", "Rojo"))
semaforo

# (b) Reporte con for
for (k in seq_along(vendedores)) {
  cat(sprintf("%-6s %4d mil  %6.1f%%  %s\n", vendedores[k],
              ventas_trim[k], cumplimiento[k], semaforo[k]))
}

# (c) Trimestres para que el de menor venta alcance la meta
venta <- min(ventas_trim)
trimestres <- 0
while (venta < meta) {
  venta <- venta * 1.08
  trimestres <- trimestres + 1
}
vendedores[which.min(ventas_trim)]
trimestres
round(venta, 1)
\end{rcode}
\begin{rsalida}
[1] "Amarillo" "Verde"    "Rojo"     "Verde"    "Amarillo" "Verde"   
Luis    185 mil    92.5%  Amarillo
Marta   240 mil   120.0%  Verde
Nora    132 mil    66.0%  Rojo
Óscar  310 mil   155.0%  Verde
Paula   198 mil    99.0%  Amarillo
Raúl   260 mil   130.0%  Verde
[1] "Nora"
[1] 6
[1] 209.5
\end{rsalida}

Tres vendedores están en verde, dos en amarillo (Luis y Paula, cerca de la
meta) y Nora en rojo con 66\% de cumplimiento. Aun creciendo 8\% cada
trimestre, Nora necesitaría 6 trimestres (año y medio) para llegar a la
meta: es un caso para un plan de acompañamiento, no solo para esperar.
\end{solucion}

\begin{solucion}{m1-9}
Las validaciones van al inicio de cada función, antes de los cálculos. El
punto de equilibrio en unidades es
\codigo{costos\_fijos / (precio - costo\_variable)}, redondeado hacia
arriba porque no se venden fracciones de unidad.

\begin{rcode}
# Precio final: aplica descuento (proporción) y luego IVA
precio_final <- function(precio, descuento = 0, iva = 0.16) {
  if (any(precio < 0)) stop("El precio no puede ser negativo")
  if (descuento < 0 || descuento > 1) {
    stop("El descuento debe estar entre 0 y 1 (0.15 = 15%)")
  }
  round(precio * (1 - descuento) * (1 + iva), 2)
}

# Unidades mínimas para no perder dinero
punto_equilibrio <- function(costos_fijos, precio, costo_variable) {
  if (precio <= costo_variable) {
    stop("El precio debe ser mayor que el costo variable")
  }
  ceiling(costos_fijos / (precio - costo_variable))
}

# (c) Uso
catalogo <- c(laptop = 18999, mouse = 349, monitor = 5499, silla = 3899)
precio_final(catalogo, descuento = 0.20)
precio_final(1000)                          # Solo IVA
punto_equilibrio(85000, precio = 450, costo_variable = 280)

# (d) ¿Qué pasa con un descuento inválido? tryCatch() atrapa el error
tryCatch(precio_final(1000, descuento = 1.5),
         error = function(e) conditionMessage(e))
\end{rcode}
\begin{rsalida}
  laptop    mouse  monitor    silla 
17631.07   323.87  5103.07  3618.27 
[1] 1160
[1] 500
[1] "El descuento debe estar entre 0 y 1 (0.15 = 15%)"
\end{rsalida}

El negocio necesita vender 501 unidades para cubrir sus \$85,000 de costos
fijos (con 500 le faltarían \$20). \codigo{tryCatch()} ejecuta el código y,
si ocurre un error, en lugar de detener el script llama a la función de
\codigo{error}; aquí solo devolvemos el mensaje. Sin \codigo{tryCatch()},
R mostraría:

\begin{rnoejecutar}
precio_final(1000, descuento = 1.5)
\end{rnoejecutar}
\begin{rsalida}
Error in precio_final(1000, descuento = 1.5) :
  El descuento debe estar entre 0 y 1 (0.15 = 15%)
\end{rsalida}
\end{solucion}

\begin{solucion}{m1-10}
Leemos el catálogo, buscamos los márgenes negativos (un problema de calidad
de datos o de precios) y exportamos dos hojas.

\begin{rcode}
library(readr)
library(writexl)

productos <- read_csv("datasets/productos.csv", show_col_types = FALSE)

# (a) Tamaño
dim(productos)

# (b) Margen negativo
negativos <- productos[productos$margen_pct < 0, ]
nrow(negativos)
peores <- negativos[order(negativos$margen_pct),
                    c("nombre_producto", "categoria", "precio_catalogo",
                      "costo", "margen_pct")]
head(peores, 5)

# (c) Margen promedio por categoría
margen_cat <- aggregate(margen_pct ~ categoria, data = productos,
                        FUN = mean)
margen_cat$margen_pct <- round(margen_cat$margen_pct, 1)
margen_cat[order(margen_cat$margen_pct, decreasing = TRUE), ]

# (d) Problemas de inventario
table(productos$estado_stock)
sum(productos$estado_stock %in% c("Stock Bajo", "Sin Stock"))

# (e) Exportar
dir.create("resultados", showWarnings = FALSE)
write_xlsx(list(Margen_negativo = as.data.frame(peores),
                Por_categoria = margen_cat),
           "resultados/calidad_catalogo.xlsx")
\end{rcode}
\begin{rsalida}
[1] 50 13
[1] 10
# A tibble: 5 x 5
  nombre_producto categoria   precio_catalogo costo margen_pct
  <chr>           <chr>                 <dbl> <dbl>      <dbl>
1 Producto 05     Alimentos              25.5  355.    -1292. 
2 Producto 46     Automotriz             64.5  178.     -175. 
3 Producto 32     Hogar                 143.   210.      -46.7
4 Producto 26     Electrónica           115.   159.      -38.3
5 Producto 29     Belleza               298.   392.      -31.6
     categoria margen_pct
4     Deportes       73.2
8       Libros       73.0
9     Mascotas       44.5
5  Electrónica       43.3
7     Juguetes       39.5
3      Belleza       38.6
10        Ropa       35.7
2   Automotriz       16.8
6        Hogar       -6.3
1    Alimentos     -161.6

Stock Bajo   Stock OK 
         4         46 
[1] 4
\end{rsalida}

\textbf{Conclusiones para Compras:} 12 de los 50 productos (24\%) se venden
por debajo de su costo; en los peores casos el costo es más de cuatro veces
el precio, lo que sugiere errores de captura o precios desactualizados que
hay que revisar antes de cualquier análisis de rentabilidad. Además, 4
productos tienen stock bajo y ninguno está agotado. (Los datos del curso
son simulados, por eso aparecen casos tan extremos; en una empresa real
serían una alerta inmediata.)
\end{solucion}

\begin{solucion}{m1-11}
Resumimos con \codigo{aggregate()} y \codigo{tapply()}, y usamos
\codigo{\%in\%} para agrupar las dos tarjetas.

\begin{rcode}
library(readr)

trans <- read_csv("datasets/transacciones.csv", show_col_types = FALSE)

# (a) Resumen por método de pago
resumen_pago <- aggregate(total_transaccion ~ metodo_pago, data = trans,
                          FUN = sum)
resumen_pago$transacciones <- as.numeric(table(trans$metodo_pago)[
  resumen_pago$metodo_pago])
resumen_pago$ticket_promedio <- round(resumen_pago$total_transaccion /
                                        resumen_pago$transacciones, 2)
resumen_pago

# (b) Porcentaje del monto pagado con tarjeta
con_tarjeta <- trans$metodo_pago %in% c("Tarjeta Crédito",
                                        "Tarjeta Débito")
round(sum(trans$total_transaccion[con_tarjeta]) /
        sum(trans$total_transaccion) * 100, 1)

# (c) Mejor mes y día con más transacciones
por_mes <- tapply(trans$total_transaccion, trans$mes, sum)
names(which.max(por_mes))
round(max(por_mes), 2)
sort(table(trans$dia_semana), decreasing = TRUE)[1:3]

# (d) Exportar y comprobar
dir.create("resultados", showWarnings = FALSE)
write_csv(resumen_pago, "resultados/resumen_metodo_pago.csv")
saveRDS(resumen_pago, "resultados/resumen_metodo_pago.rds")
identical(resumen_pago, readRDS("resultados/resumen_metodo_pago.rds"))
\end{rcode}
\begin{rsalida}
      metodo_pago total_transaccion transacciones ticket_promedio
1        Efectivo         303734.65           315          964.24
2 Tarjeta Crédito         351170.69           350         1003.34
3  Tarjeta Débito         220855.30           233          947.88
4   Transferencia          84232.05           102          825.80
[1] 59.6
[1] "2023-12"
[1] 96393.61

  viernes miércoles   domingo 
      158       155       146 
[1] TRUE
\end{rsalida}

La tarjeta de crédito es el método más usado y el de mayor monto, y las
dos tarjetas juntas concentran el 59.6\% del dinero cobrado: la
disponibilidad de terminales es crítica. El ticket promedio casi no cambia
entre métodos. Como \codigo{resumen\_pago} es un data frame de R base, el
RDS lo recupera idéntico.
\end{solucion}

\begin{solucion}{m1-12}
Dividimos el problema en piezas: una función que resume una tienda, un
bucle que la aplica a las diez y, al final, el ranking y la exportación.
Para traer el nombre, la ciudad y el tamaño de la tienda usamos
\codigo{match()}, que encuentra en qué fila de \codigo{tiendas} está cada
\codigo{tienda\_id}.

\begin{rcode}
library(readr)
library(writexl)

ventas <- read_csv("datasets/ventas_retail.csv", show_col_types = FALSE)
tiendas <- read_csv("datasets/tiendas.csv", show_col_types = FALSE)

# Resume las ventas de una tienda en un data frame de una fila
resumen_tienda <- function(id, ventas, tiendas) {
  if (!id %in% tiendas$tienda_id) stop("La tienda no existe: ", id)
  v <- ventas[ventas$tienda_id == id, ]
  fila_tienda <- match(id, tiendas$tienda_id)
  venta_mes <- tapply(v$total, v$mes, sum)
  data.frame(
    tienda = tiendas$nombre_tienda[fila_tienda],
    ciudad = tiendas$ciudad[fila_tienda],
    num_ventas = nrow(v),
    venta_total = round(sum(v$total), 2),
    ticket_promedio = round(mean(v$total), 2),
    mejor_mes = names(which.max(venta_mes)),
    venta_m2 = round(sum(v$total) / tiendas$tamano_m2[fila_tienda], 2)
  )
}

resumen_tienda(1, ventas, tiendas)          # Prueba con una tienda

# Aplicar a las 10 tiendas y unir las filas
lista_resumenes <- lapply(tiendas$tienda_id, resumen_tienda,
                          ventas = ventas, tiendas = tiendas)
ranking <- do.call(rbind, lista_resumenes)
ranking <- ranking[order(ranking$venta_total, decreasing = TRUE), ]
rownames(ranking) <- NULL

# Ranking impreso
cat("RANKING DE TIENDAS 2023\n")
for (k in seq_len(nrow(ranking))) {
  cat(sprintf("%2d. %-17s %11s  %3d ventas  %8s por m2\n", k,
              ranking$tienda[k], scales::dollar(ranking$venta_total[k]),
              ranking$num_ventas[k], scales::dollar(ranking$venta_m2[k],
                                                     accuracy = 0.01)))
}

# Exportar
dir.create("resultados", showWarnings = FALSE)
write_xlsx(ranking, "resultados/ranking_tiendas.xlsx")
ranking$tienda[which.max(ranking$venta_m2)]
\end{rcode}
\begin{rsalida}
           tienda ciudad num_ventas venta_total ticket_promedio
1 Sucursal Centro   CDMX         32    43898.54         1371.83
  mejor_mes venta_m2
1   2023-01     87.8
RANKING DE TIENDAS 2023
 1. Sucursal Outlet    $70,154.72   47 ventas   $200.44 por m2
 2. Sucursal Premium   $61,369.12   46 ventas    $76.71 por m2
 3. Sucursal Plaza B   $55,290.32   37 ventas   $115.19 por m2
 4. Sucursal Norte     $54,846.35   38 ventas   $121.88 por m2
 5. Sucursal Sur       $54,065.76   37 ventas    $90.11 por m2
 6. Sucursal Plaza A   $54,050.18   34 ventas    $77.21 por m2
 7. Sucursal Oeste     $53,611.23   37 ventas    $97.47 por m2
 8. Sucursal Centro    $43,898.54   32 ventas    $87.80 por m2
 9. Sucursal Express   $42,171.47   34 ventas   $140.57 por m2
10. Sucursal Este      $24,458.92   23 ventas    $61.15 por m2
[1] "Sucursal Outlet"
\end{rsalida}

\textbf{¿Qué nos dice esto?} La Sucursal Outlet es la que más vende en
total (\$70,155) y, con solo 350 m$^2$, también la más productiva por metro
cuadrado (\$200.44/m$^2$). En cambio, la Sucursal Premium, la más grande
(800 m$^2$), vende menos de la mitad por metro cuadrado. La Sucursal Este
es la de menor venta total, con solo 20 ventas en el año. Recuerda que en
estos datos cada registro es una venta diaria asignada al azar a una
tienda, así que las diferencias entre tiendas son pequeñas; con datos
reales, este mismo script te daría el ranking verdadero cada mes.

\codigo{lapply()} aplica la función a cada \codigo{tienda\_id} y devuelve
una lista de data frames de una fila; \codigo{do.call(rbind, lista)} los
apila en una sola tabla. Es el patrón de R base para ``hacer lo mismo por
cada grupo y juntar los resultados''; en el Módulo~\ref{cap:manipulacion}
lo harás con \codigo{group\_by()} y \codigo{summarise()}.
\end{solucion}


\section{Módulo 2. Manipulación de datos con dplyr y tidyr}
% ============================================================================
% Soluciones del Módulo 2
% ============================================================================

\begin{solucion}{m2-1}
Filtramos con dos condiciones, elegimos columnas y ordenamos; después
resumimos el resultado con \codigo{summarise()}.

\begin{rcode}
library(tidyverse)
\end{rcode}

\begin{rcode}
ventas <- read_csv("datasets/ventas_retail.csv", show_col_types = FALSE)

tienda3_desc <- ventas %>%
  filter(tienda_id == 3, descuento_pct >= 10) %>%   # ambas condiciones
  select(fecha, producto_id, descuento_pct, total) %>%
  arrange(desc(total))
tienda3_desc

tienda3_desc %>%
  summarise(ventas = n(), total = sum(total))
\end{rcode}
\begin{rsalida}
# A tibble: 4 x 4
  fecha      producto_id descuento_pct total
  <date>           <dbl>         <dbl> <dbl>
1 2023-11-05          27            10 1439.
2 2023-03-02          26            20 1227.
3 2023-12-20          17            10 1098.
4 2023-07-08           6            10  742.
# A tibble: 1 x 2
  ventas total
   <int> <dbl>
1      4 4506.
\end{rsalida}

La tienda 3 hizo 4 ventas con descuento de 10\% o más, que suman
\$4,506; la mayor fue de \$1,439 el 5 de noviembre. Observa que en
\codigo{summarise(total = sum(total))} el nombre nuevo puede coincidir con
el de la columna: dentro de la función, \codigo{total} todavía es la
columna original.
\end{solucion}

\begin{solucion}{m2-2}
Combinamos tres condiciones en \codigo{filter()} (el producto activo y el
stock bajo el mínimo) y calculamos el faltante con \codigo{mutate()}.

\begin{rcode}
library(tidyverse)
productos <- read_csv("datasets/productos.csv", show_col_types = FALSE)

resurtir <- productos %>%
  filter(activo, stock_actual < stock_minimo) %>%
  mutate(faltante = stock_minimo - stock_actual) %>%
  select(nombre_producto, proveedor, stock_actual, stock_minimo,
         faltante) %>%
  arrange(desc(faltante))
resurtir
resurtir %>% count(proveedor, wt = faltante, sort = TRUE)
\end{rcode}
\begin{rsalida}
# A tibble: 3 x 5
  nombre_producto proveedor   stock_actual stock_minimo faltante
  <chr>           <chr>              <dbl>        <dbl>    <dbl>
1 Producto 03     Proveedor 3           22           26        4
2 Producto 33     Proveedor 1           42           45        3
3 Producto 06     Proveedor 1           36           38        2
# A tibble: 2 x 2
  proveedor       n
  <chr>       <dbl>
1 Proveedor 1     5
2 Proveedor 3     4
\end{rsalida}

Solo 3 productos activos están bajo su mínimo. Por unidades faltantes hay
que llamar primero al Proveedor 1 (5 unidades en dos productos), aunque el
producto con el mayor faltante individual (Producto 03, 4 unidades) es del
Proveedor 3. \codigo{filter(activo, \ldots)} funciona porque
\codigo{activo} ya es lógica.
\end{solucion}

\begin{solucion}{m2-3}
\codigo{case\_when()} evalúa las reglas en orden, así que basta con dar el
límite superior de cada rango. El porcentaje se calcula después de
\codigo{count()}.

\begin{rcode}
library(tidyverse)
clientes <- read_csv("datasets/clientes.csv", show_col_types = FALSE)

clientes %>%
  mutate(rango_edad = case_when(
    edad < 30 ~ "18-29",
    edad < 45 ~ "30-44",
    edad < 60 ~ "45-59",
    .default  = "60+"
  )) %>%
  count(rango_edad) %>%
  mutate(pct = round(n / sum(n) * 100, 1))
\end{rcode}
\begin{rsalida}
# A tibble: 4 x 3
  rango_edad     n   pct
  <chr>      <int> <dbl>
1 18-29         33  16.5
2 30-44         52  26  
3 45-59         51  25.5
4 60+           64  32  
\end{rsalida}

El grupo más grande es el de 60 años o más (64 clientes, 32\%) y el más
chico el de 18 a 29 (33 clientes, 16.5\%): una base de clientes madura.
\end{solucion}

\begin{solucion}{m2-4}
El promedio de una columna lógica es la proporción de \codigo{TRUE}: con
\codigo{mean(es\_premium)} obtenemos directamente el porcentaje de premium.

\begin{rcode}
library(tidyverse)
clientes <- read_csv("datasets/clientes.csv", show_col_types = FALSE)

perfil_ciudad <- clientes %>%
  group_by(ciudad) %>%
  summarise(clientes    = n(),
            pct_premium = round(mean(es_premium) * 100, 1),
            edad_prom   = round(mean(edad), 1),
            .groups = "drop") %>%
  arrange(desc(clientes))
perfil_ciudad
perfil_ciudad %>% slice_max(pct_premium, n = 1)
\end{rcode}
\begin{rsalida}
# A tibble: 10 x 4
   ciudad      clientes pct_premium edad_prom
   <chr>          <int>       <dbl>     <dbl>
 1 Puebla            28        42.9      51.3
 2 CDMX              27        48.1      56.2
 3 Guadalajara       22        27.3      49.2
 4 Cancún            20        30        45  
 5 Toluca            20        30        51  
 6 León              18        16.7      39.1
 7 Mérida            17        23.5      47.6
 8 Querétaro         17        35.3      53.7
 9 Tijuana           17        41.2      37.8
10 Monterrey         14        71.4      51.1
# A tibble: 1 x 4
  ciudad    clientes pct_premium edad_prom
  <chr>        <int>       <dbl>     <dbl>
1 Monterrey       14        71.4      51.1
\end{rsalida}

Puebla y la CDMX tienen más clientes, pero Monterrey tiene la mayor
proporción de premium: 71.4\% de sus 14 clientes. Es una plaza chica pero de
alto valor.
\end{solucion}

\begin{solucion}{m2-5}
\begin{rcode}
library(tidyverse)
transacciones <- read_csv("datasets/transacciones.csv",
                          show_col_types = FALSE)

transacciones %>%
  mutate(dia = wday(fecha, label = TRUE, abbr = FALSE)) %>%
  group_by(dia) %>%
  summarise(transacciones = n(),
            ingresos = sum(total_transaccion),
            ticket_promedio = round(mean(total_transaccion), 2),
            .groups = "drop") %>%
  mutate(participacion_pct = round(ingresos / sum(ingresos) * 100, 1))
\end{rcode}
\begin{rsalida}
# A tibble: 7 x 5
  dia       transacciones ingresos ticket_promedio participacion_pct
  <ord>             <int>    <dbl>           <dbl>             <dbl>
1 domingo             146  148703.           1019.              15.5
2 lunes               136  128878.            948.              13.4
3 martes              129  119364.            925.              12.4
4 miércoles           155  141329.            912.              14.7
5 jueves              146  145280.            995.              15.1
6 viernes             158  161266.           1021.              16.8
7 sábado              130  115172.            886.              12  
\end{rsalida}

El viernes es el día con más transacciones (158) e ingresos (\$161,266,
16.8\% del total), seguido del domingo; el sábado es el más flojo (12\%).
Conviene reforzar el personal los viernes y domingos.
\end{solucion}

\begin{solucion}{m2-6}
\begin{rcode}
library(tidyverse)
empleados <- read_csv("datasets/empleados.csv", show_col_types = FALSE)
fecha_corte <- as.Date("2024-01-01")

empleados_antig <- empleados %>%
  mutate(
    antig_calc = time_length(interval(fecha_ingreso, fecha_corte),
                             "years"),
    diferencia = abs(antig_calc - antiguedad_anos),
    grupo = case_when(
      antig_calc < 3  ~ "1. Menos de 3 años",
      antig_calc <= 6 ~ "2. 3 a 6 años",
      .default        = "3. Más de 6 años"
    )
  )
max(empleados_antig$diferencia)          # diferencia máxima

empleados_antig %>%
  group_by(grupo) %>%
  summarise(empleados = n(),
            salario_prom = round(mean(salario_mensual), 2),
            .groups = "drop")
\end{rcode}
\begin{rsalida}
[1] 0.05245902
# A tibble: 3 x 3
  grupo              empleados salario_prom
  <chr>                  <int>        <dbl>
1 1. Menos de 3 años        41       27354.
2 2. 3 a 6 años             34       27821.
3 3. Más de 6 años          25       27228.
\end{rsalida}

La diferencia máxima entre ambos cálculos es de 0.05 años (unos 19
días): la columna original dividía los días entre 365, mientras que
\codigo{time\_length()} considera los años bisiestos. El salario promedio
casi no cambia con la antigüedad (entre \$27,228 y \$27,821): en estos datos
simulados la antigüedad no se refleja en el sueldo, algo que en una empresa
real merecería una revisión de política salarial.
\end{solucion}

\begin{solucion}{m2-7}
Solo necesitamos del catálogo la categoría y el costo, así que los
seleccionamos dentro del \emph{join}.

\begin{rcode}
library(tidyverse)
transacciones <- read_csv("datasets/transacciones.csv",
                          show_col_types = FALSE)
productos <- read_csv("datasets/productos.csv", show_col_types = FALSE)

margen_cat <- transacciones %>%
  left_join(select(productos, producto_id, categoria, costo),
            by = "producto_id") %>%
  mutate(utilidad = (precio_venta - costo) * cantidad) %>%
  group_by(categoria) %>%
  summarise(con_perdida = sum(utilidad < 0),   # antes de resumir
            ingresos = sum(total_transaccion),
            utilidad = sum(utilidad),
            .groups = "drop") %>%
  mutate(margen_pct = round(utilidad / ingresos * 100, 1)) %>%
  arrange(desc(utilidad))
margen_cat
\end{rcode}
\begin{rsalida}
# A tibble: 10 x 5
   categoria   con_perdida ingresos utilidad margen_pct
   <chr>             <int>    <dbl>    <dbl>      <dbl>
 1 Alimentos            13  121543.   73870.       60.8
 2 Electrónica          27  132672.   73067.       55.1
 3 Juguetes             30  125093.   57714.       46.1
 4 Belleza              66  182143.   56054.       30.8
 5 Automotriz           26  104085.   55538.       53.4
 6 Deportes             20   86171.   48301.       56.1
 7 Mascotas             10   57668.   35423.       61.4
 8 Libros               21   57162.   25757.       45.1
 9 Hogar                24   63407.   13079.       20.6
10 Ropa                 15   30048.    5835.       19.4
\end{rsalida}

Alimentos es la categoría que más utilidad deja (\$73,870), aunque
Belleza tiene los mayores ingresos (\$182,143). Belleza es también la más
problemática: 66 transacciones con utilidad negativa y un margen de 30.8\%.
Ropa y Hogar tienen los márgenes más bajos (19--21\%). Cuidado con el orden
dentro de \codigo{summarise()}: si calculas \codigo{utilidad = sum(utilidad)}
\emph{antes} de \codigo{sum(utilidad < 0)}, la segunda usaría la suma ya
calculada y el conteo saldría 0.
\end{solucion}

\begin{solucion}{m2-8}
Trabajamos con listas de clientes \emph{únicos} de cada tabla
(\codigo{distinct()}) para que los conteos sean de clientes y no de ventas.

\begin{rcode}
library(tidyverse)
ventas <- read_csv("datasets/ventas_retail.csv", show_col_types = FALSE)
transacciones <- read_csv("datasets/transacciones.csv",
                          show_col_types = FALSE)
clientes <- read_csv("datasets/clientes.csv", show_col_types = FALSE)

cli_ventas <- ventas %>% distinct(cliente_id)
cli_trans  <- transacciones %>% distinct(cliente_id)

# (a) En ambas tablas
cli_ventas %>% semi_join(cli_trans, by = "cliente_id") %>% nrow()
# (b) En transacciones pero no en ventas
cli_trans %>% anti_join(cli_ventas, by = "cliente_id") %>% nrow()
# (c) Registrados que no aparecen en ninguna
compradores <- bind_rows(cli_ventas, cli_trans) %>% distinct()
clientes %>% anti_join(compradores, by = "cliente_id") %>%
  select(cliente_id, nombre, fecha_registro)
\end{rcode}
\begin{rsalida}
[1] 170
[1] 28
# A tibble: 0 x 3
# i 3 variables: cliente_id <dbl>, nombre <chr>,
#   fecha_registro <date>
\end{rsalida}

170 clientes compraron en ambas tablas y 28 aparecen solo en
transacciones. Ningún cliente registrado quedó sin comprar en alguna de las
dos: el resultado de (c) es una tabla vacía (0 filas), lo cual también es
una respuesta válida.
\end{solucion}

\begin{solucion}{m2-9}
Calculamos el porcentaje en formato largo (agrupado por tienda) y al final
pivotamos a ancho.

\begin{rcode}
library(tidyverse)
transacciones <- read_csv("datasets/transacciones.csv",
                          show_col_types = FALSE)
tiendas <- read_csv("datasets/tiendas.csv", show_col_types = FALSE)

pagos_tienda <- transacciones %>%
  left_join(select(tiendas, tienda_id, nombre_tienda), by = "tienda_id") %>%
  group_by(nombre_tienda, metodo_pago) %>%
  summarise(ingresos = sum(total_transaccion), .groups = "drop_last") %>%
  mutate(pct = round(ingresos / sum(ingresos) * 100, 1)) %>%
  ungroup() %>%
  select(-ingresos) %>%
  pivot_wider(names_from = metodo_pago, values_from = pct,
              values_fill = 0) %>%
  arrange(desc(Efectivo))
pagos_tienda
\end{rcode}
\begin{rsalida}
# A tibble: 10 x 5
   nombre_tienda    Efectivo `Tarjeta Crédito` `Tarjeta Débito`
   <chr>               <dbl>             <dbl>            <dbl>
 1 Sucursal Oeste       45.5              35.6             10.7
 2 Sucursal Este        37.5              27               25.6
 3 Sucursal Norte       37.1              43.5             14.2
 4 Sucursal Plaza A     36.2              21.5             31.4
 5 Sucursal Plaza B     33.5              42.8             18.1
 6 Sucursal Sur         29.3              36.4             29.6
 7 Sucursal Outlet      29.1              29.5             28  
 8 Sucursal Centro      28.9              44               21.2
 9 Sucursal Premium     24.6              40.8             25.2
10 Sucursal Express     18.8              47.6             20.8
# i 1 more variable: Transferencia <dbl>
\end{rsalida}

La Sucursal Oeste es la más dependiente del efectivo: 45.5\% de sus
ingresos. En el otro extremo, la Sucursal Express cobra solo 18.8\% en
efectivo y 47.6\% con tarjeta de crédito. Como los porcentajes se calcularon
con el resultado agrupado por tienda (\codigo{.groups = "drop\_last"}),
cada fila suma 100.
\end{solucion}

\begin{solucion}{m2-10}
Primero completamos los meses faltantes; si no, \codigo{lag()} compararía
contra el mes equivocado. \codigo{group\_by(tienda\_id)} hace que
\codigo{lag()} y \codigo{cumsum()} reinicien en cada tienda.

\begin{rcode}
library(tidyverse)
ventas <- read_csv("datasets/ventas_retail.csv", show_col_types = FALSE)

tienda_mes <- ventas %>%
  group_by(tienda_id, mes) %>%
  summarise(ingresos = sum(total), .groups = "drop") %>%
  complete(tienda_id, mes, fill = list(ingresos = 0)) %>%
  arrange(tienda_id, mes) %>%
  group_by(tienda_id) %>%
  mutate(cambio_pesos = ingresos - lag(ingresos),
         acumulado    = cumsum(ingresos)) %>%
  ungroup()
nrow(tienda_mes)

tienda_mes %>%
  slice_min(cambio_pesos, n = 3) %>%        # las 3 mayores caídas
  mutate(across(where(is.numeric), ~ round(.x, 2)))
\end{rcode}
\begin{rsalida}
[1] 120
# A tibble: 3 x 5
  tienda_id mes     ingresos cambio_pesos acumulado
      <dbl> <chr>      <dbl>        <dbl>     <dbl>
1         8 2023-08     781.      -12373.    50452 
2         8 2023-03    1222.       -9457.    17226.
3         2 2023-05    3698.       -7626.    25458.
\end{rsalida}

La tabla completa tiene 120 filas (10 tiendas $\times$ 12 meses) gracias
a \codigo{complete()}. La mayor caída fue la de la tienda 8 en agosto:
\$12,373 menos que en julio, al vender solo \$781 en el mes.
\end{solucion}

\begin{solucion}{m2-11}
Seguimos el mismo orden que en la sección~\ref{sec:m2-limpieza}:
duplicados, texto, números y fechas, y al final reglas de negocio.

\begin{rcode}
library(tidyverse)
lista_proveedor <- tibble(
  id        = c(1, 2, 2, 3, 4, 5, 6),
  categoria = c("electronica", "Electrónica ", "Electrónica ",
                "HOGAR", "hogar", "Ropa", " ropa"),
  precio    = c("$1,250.00", "$899.50", "$899.50", "450", "$0",
                "$320.75", "-15"),
  fecha     = c("2023-05-01", "02/05/2023", "02/05/2023", "2023-05-03",
                "04/05/2023", "2023-05-05", "06/05/2023")
)

lista_limpia <- lista_proveedor %>%
  distinct() %>%                                   # 1. duplicados
  mutate(
    categoria = str_to_title(str_squish(categoria)),   # 2. texto
    categoria = recode(categoria, "Electronica" = "Electrónica"),
    precio    = parse_number(precio),                  # 3. número
    fecha     = as_date(parse_date_time(fecha,         # 4. fecha
                                        orders = c("ymd", "dmy"))),
    sku       = str_c("SKU-", str_pad(id, 4, pad = "0"))
  ) %>%
  filter(precio > 0)                               # 5. regla de negocio
lista_limpia

lista_limpia %>%
  group_by(categoria) %>%
  summarise(productos = n(), precio_prom = mean(precio),
            .groups = "drop")
\end{rcode}
\begin{rsalida}
# A tibble: 4 x 5
     id categoria   precio fecha      sku     
  <dbl> <chr>        <dbl> <date>     <chr>   
1     1 Electrónica  1250  2023-05-01 SKU-0001
2     2 Electrónica   900. 2023-05-02 SKU-0002
3     3 Hogar         450  2023-05-03 SKU-0003
4     5 Ropa          321. 2023-05-05 SKU-0005
# A tibble: 3 x 3
  categoria   productos precio_prom
  <chr>           <int>       <dbl>
1 Electrónica         2       1075.
2 Hogar               1        450 
3 Ropa                1        321.
\end{rsalida}

De 7 registros quedaron 4: se eliminaron el duplicado del id 2, el precio
en cero (id 4) y el negativo (id 6). Las categorías quedaron unificadas y
las fechas en un solo formato.
\end{solucion}

\begin{solucion}{m2-12}
Construimos una tabla maestra, la validamos con \codigo{stopifnot()},
calculamos las tres hojas y las exportamos con \paquete{writexl}.

\begin{rcode}
library(tidyverse)
library(writexl)
transacciones <- read_csv("datasets/transacciones.csv",
                          show_col_types = FALSE)
productos <- read_csv("datasets/productos.csv", show_col_types = FALSE)
tiendas   <- read_csv("datasets/tiendas.csv", show_col_types = FALSE)

maestra_t <- transacciones %>%
  left_join(select(productos, producto_id, nombre_producto, categoria),
            by = "producto_id") %>%
  left_join(select(tiendas, tienda_id, nombre_tienda), by = "tienda_id")
stopifnot(nrow(maestra_t) == nrow(transacciones))   # control

resumen <- maestra_t %>%
  group_by(nombre_tienda) %>%
  summarise(ingresos = sum(total_transaccion), transacciones = n(),
            ticket_promedio = round(mean(total_transaccion), 2),
            clientes_unicos = n_distinct(cliente_id), .groups = "drop") %>%
  mutate(ranking = min_rank(desc(ingresos))) %>%
  arrange(ranking)

trimestres <- maestra_t %>%
  group_by(nombre_tienda, trimestre) %>%
  summarise(ingresos = round(sum(total_transaccion)), .groups = "drop") %>%
  pivot_wider(names_from = trimestre, values_from = ingresos,
              values_fill = 0)

top_producto <- maestra_t %>%
  group_by(nombre_tienda, nombre_producto, categoria) %>%
  summarise(ingresos = sum(total_transaccion), .groups = "drop") %>%
  group_by(nombre_tienda) %>%
  slice_max(ingresos, n = 1, with_ties = FALSE) %>%
  ungroup()

dir.create("resultados", showWarnings = FALSE)
write_xlsx(list("Resumen" = resumen, "Trimestres" = trimestres,
                "Top producto" = top_producto),
           "resultados/reporte_tiendas.xlsx")
readxl::excel_sheets("resultados/reporte_tiendas.xlsx")
resumen %>% select(ranking, nombre_tienda, ingresos, ticket_promedio)
top_producto %>% head(3)
\end{rcode}
\begin{rsalida}
[1] "Resumen"      "Trimestres"   "Top producto"
# A tibble: 10 x 4
   ranking nombre_tienda    ingresos ticket_promedio
     <int> <chr>               <dbl>           <dbl>
 1       1 Sucursal Outlet   114377.           1079.
 2       2 Sucursal Sur      105305.            966.
 3       3 Sucursal Plaza A  105216.            965.
 4       4 Sucursal Express  102791.            952.
 5       5 Sucursal Premium  101538.            940.
 6       6 Sucursal Norte     89024.            899.
 7       7 Sucursal Este      88420.            941.
 8       8 Sucursal Centro    88414.           1016.
 9       9 Sucursal Oeste     84672.            891.
10      10 Sucursal Plaza B   80236.            944.
# A tibble: 3 x 4
  nombre_tienda    nombre_producto categoria   ingresos
  <chr>            <chr>           <chr>          <dbl>
1 Sucursal Centro  Producto 40     Electrónica    5611.
2 Sucursal Este    Producto 47     Automotriz     6652.
3 Sucursal Express Producto 27     Juguetes       6072.
\end{rsalida}

El archivo tiene las tres hojas. La Sucursal Outlet vuelve a ser la
número uno (\$114,377 en transacciones), igual que en \codigo{ventas}. Usamos
\codigo{with\_ties = FALSE} en \codigo{slice\_max()} para garantizar un solo
producto por tienda aunque hubiera empates.
\end{solucion}


\section{Módulo 3. Visualización de datos con ggplot2}
% ============================================================================
% Soluciones del Módulo 3 - Visualización de datos con ggplot2
% ============================================================================

\begin{solucion}{m3-1}
Contamos los tickets por tienda con \codigo{count()} y graficamos con
\codigo{geom\_col()} en horizontal, ordenando con \codigo{fct\_reorder()}.
\end{solucion}

\begin{rcode}
library(tidyverse)
library(scales)
color_principal <- "#2a78d6"

transacciones <- read_csv("datasets/transacciones.csv",
                          show_col_types = FALSE)
tiendas <- read_csv("datasets/tiendas.csv", show_col_types = FALSE)

tickets_tienda <- transacciones %>%
  left_join(tiendas, by = "tienda_id") %>%
  mutate(tienda = str_remove(nombre_tienda, "Sucursal ")) %>%
  group_by(tienda) %>%
  summarise(tickets  = n(),
            ingresos = sum(total_transaccion),
            .groups  = "drop") %>%
  arrange(desc(tickets))
head(tickets_tienda, 4)

s1 <- ggplot(tickets_tienda,
             aes(x = tickets, y = fct_reorder(tienda, tickets))) +
  geom_col(fill = color_principal, width = 0.7) +
  scale_x_continuous(expand = expansion(mult = c(0, 0.05))) +
  labs(title = "Sur y Plaza A atendieron más tickets en 2023",
       x = "Número de tickets", y = NULL,
       caption = "Fuente: transacciones.csv") +
  theme_minimal() +
  theme(panel.grid.minor = element_blank(),
        panel.grid.major.y = element_blank())
s1
\end{rcode}
\begin{rsalida}
# A tibble: 4 x 3
  tienda  tickets ingresos
  <chr>     <int>    <dbl>
1 Plaza A     109  105216.
2 Sur         109  105305.
3 Express     108  102791.
4 Premium     108  101538.
\end{rsalida}

\textbf{Interpretación.} Sur y Plaza A empatan con 109 tickets, pero la
tienda con más \emph{ingresos} es Outlet (106 tickets): vende un poco menos
de tickets pero de mayor importe. Contar tickets e ingresos responde
preguntas distintas (tráfico contra valor).


\begin{solucion}{m3-2}
Este es el código completo, con su salida y la interpretación.
\end{solucion}

\begin{rcode}
library(tidyverse)
color_principal <- "#2a78d6"
color_resalte   <- "#eb6834"

clientes <- read_csv("datasets/clientes.csv", show_col_types = FALSE)
mediana_edad <- median(clientes$edad)
mediana_edad

s2 <- ggplot(clientes, aes(x = edad)) +
  geom_histogram(binwidth = 5, boundary = 0,
                 fill = color_principal, color = "white") +
  geom_vline(xintercept = mediana_edad, linetype = "dashed",
             color = color_resalte, linewidth = 0.8) +
  annotate("text", x = mediana_edad, y = Inf, vjust = 1.5,
           hjust = -0.1, label = paste("Mediana:", mediana_edad),
           color = color_resalte, size = 3.5) +
  labs(title = "Edad de los clientes",
       subtitle = "Barras de 5 años", x = "Edad", y = "Clientes") +
  theme_minimal() +
  theme(panel.grid.minor = element_blank())
s2
\end{rcode}
\begin{rsalida}
[1] 48.5
\end{rsalida}

\textbf{Interpretación.} La mediana es de 48.5 años y las barras tienen
alturas parecidas entre los 18 y los 75 años: la distribución es más o
menos \emph{pareja} (uniforme), sin una edad dominante. Es lo esperable
porque las edades se simularon al azar en ese rango; en una empresa real
casi siempre verías una concentración en algún grupo de edad.


\begin{solucion}{m3-3}
Los cuatro errores:
\begin{enumerate}
  \item \codigo{fill = "\#2a78d6"} dentro de \codigo{aes()}: el color se
        trata como un dato, sale un color por defecto y una leyenda rara. Va
        fuera de \codigo{aes()}, en la geometría.
  \item El \codigo{+} está al inicio de la línea: R termina la instrucción
        en \codigo{ggplot(...)} y falla con ``Cannot use \codigo{+} with a
        single argument''.
  \item \codigo{geom\_bar()} con \codigo{y}: debe ser \codigo{geom\_col()}
        porque los ingresos ya están calculados.
  \item \codigo{\%>\%} entre capas y \codigo{labels = percent} en un eje de
        pesos: debe ser \codigo{+} y \codigo{labels = dollar}.
\end{enumerate}
(Además conviene agregar \codigo{.groups = "drop"}, aunque aquí no es
obligatorio porque se agrupa por una sola columna.)
\end{solucion}

\begin{rcode}
library(tidyverse)
library(scales)
transacciones <- read_csv("datasets/transacciones.csv",
                          show_col_types = FALSE)

ingresos_pago <- transacciones %>%
  group_by(metodo_pago) %>%
  summarise(ingresos = sum(total_transaccion), .groups = "drop")

s3 <- ggplot(ingresos_pago, aes(x = metodo_pago, y = ingresos)) +
  geom_col(fill = "#2a78d6") +
  scale_y_continuous(labels = dollar)
s3
\end{rcode}

La versión corregida se ejecuta sin errores y muestra cuatro columnas
azules con el eje en pesos. Para dejarla lista para la dirección faltaría
ordenarla, ponerla horizontal y agregar títulos (lista de verificación de
la sección~\ref{sec:m3-antes-despues}).


\begin{solucion}{m3-4}
Este es el código completo, con su salida y la interpretación.
\end{solucion}

\begin{rcode}
library(tidyverse)
color_principal <- "#2a78d6"
color_resalte   <- "#eb6834"
color_gris      <- "#b5b3ad"

transacciones <- read_csv("datasets/transacciones.csv",
                          show_col_types = FALSE)

tickets_mes <- transacciones %>%
  mutate(mes = floor_date(fecha, unit = "month")) %>%
  count(mes, name = "tickets")
promedio <- mean(tickets_mes$tickets)
mes_min  <- tickets_mes %>% slice_min(tickets, n = 1)
promedio
mes_min

s4 <- ggplot(tickets_mes, aes(x = mes, y = tickets)) +
  geom_hline(yintercept = promedio, color = color_gris,
             linetype = "dashed", linewidth = 0.6) +
  geom_line(color = color_principal, linewidth = 0.8) +
  geom_point(color = color_principal, size = 2) +
  annotate("text", x = mes_min$mes, y = mes_min$tickets,
           label = paste("Mínimo:", mes_min$tickets, "tickets"),
           vjust = 1.8, color = color_resalte, size = 3.5) +
  scale_x_date(date_labels = "%b", date_breaks = "1 month") +
  labs(title = "Abril tuvo el menor número de tickets del año",
       subtitle = "Tickets por mes, 2023 (línea punteada = promedio)",
       x = NULL, y = "Tickets") +
  theme_minimal() +
  theme(panel.grid.minor = element_blank())
s4
\end{rcode}
\begin{rsalida}
[1] 83.33333
# A tibble: 1 x 2
  mes        tickets
  <date>       <int>
1 2023-04-01      71
\end{rsalida}

\textbf{Interpretación.} El promedio es de 83.3 tickets por mes y abril
queda muy por debajo con 71. Como abril también fue el peor mes en
ingresos, la caída de ese mes se explica en buena parte por menos tráfico
(menos tickets), no solo por tickets más pequeños.


\begin{solucion}{m3-5}
\codigo{wday(fecha, label = TRUE, week\_start = 1)} devuelve el día como
factor ordenado de lunes a domingo; así el eje queda en el orden correcto
sin hacer nada más.
\end{solucion}

\begin{rcode}
library(tidyverse)
library(scales)
color_resalte <- "#eb6834"
color_gris    <- "#b5b3ad"

transacciones <- read_csv("datasets/transacciones.csv",
                          show_col_types = FALSE)

ingresos_dia <- transacciones %>%
  mutate(dia = wday(fecha, label = TRUE, week_start = 1)) %>%
  group_by(dia) %>%
  summarise(ingresos = sum(total_transaccion), .groups = "drop") %>%
  mutate(destacar = if_else(ingresos == max(ingresos), "Máximo",
                            "Resto"))
ingresos_dia

s_dia <- ggplot(ingresos_dia, aes(x = dia, y = ingresos,
                                  fill = destacar)) +
  geom_col(width = 0.7) +
  geom_text(data = filter(ingresos_dia, destacar == "Máximo"),
            aes(label = dollar(ingresos, accuracy = 1)),
            vjust = -0.5, size = 3.5, color = color_resalte) +
  scale_fill_manual(values = c(Máximo = color_resalte,
                               Resto  = color_gris),
                    guide = "none") +
  scale_y_continuous(labels = dollar,
                     expand = expansion(mult = c(0, 0.1))) +
  labs(title = "El viernes es el día que más vende",
       subtitle = "Ingresos por día de la semana, 2023",
       x = NULL, y = "Ingresos (MXN)") +
  theme_minimal() +
  theme(panel.grid.minor = element_blank(),
        panel.grid.major.x = element_blank())
s_dia
\end{rcode}
\begin{rsalida}
# A tibble: 7 x 3
  dia   ingresos destacar
  <ord>    <dbl> <chr>   
1 lun    128878. Resto   
2 mar    119364. Resto   
3 mié    141329. Resto   
4 jue    145280. Resto   
5 vie    161266. Máximo  
6 sáb    115172. Resto   
7 dom    148703. Resto
\end{rsalida}

\begin{figure}[H]
\centering
\includegraphics[width=0.85\textwidth]{m03-sol-dia}
\caption{Solución del ejercicio m3-5.}
\label{fig:m03-sol-dia}
\end{figure}

\textbf{Interpretación.} El viernes suma \$161,266, contra \$115,172 del
sábado, el día más bajo. Pero recuerda el mapa de calor
(figura~\ref{fig:m03-calor}): el viernes no es el mejor día \emph{todos}
los meses. La diferencia anual existe, pero no es tan consistente como para
basar en ella una estrategia sin más análisis.


\begin{solucion}{m3-6}
Este es el código completo, con su salida y la interpretación.
\end{solucion}

\begin{rcode}
library(tidyverse)
library(scales)
empleados <- read_csv("datasets/empleados.csv", show_col_types = FALSE)

empleados %>%
  group_by(departamento) %>%
  summarise(mediana = median(salario_mensual),
            ric     = IQR(salario_mensual),   # rango intercuartílico
            n       = n(),
            .groups = "drop") %>%
  arrange(desc(mediana))

s6 <- ggplot(empleados,
             aes(x = salario_mensual,
                 y = fct_reorder(departamento, salario_mensual,
                                 .fun = median))) +
  geom_boxplot(fill = "#cde2fb", color = "#2a78d6", width = 0.6) +
  scale_x_continuous(labels = dollar) +
  labs(title = "IT y Operaciones tienen los salarios medianos más altos",
       x = "Salario mensual (MXN)", y = NULL) +
  theme_minimal() +
  theme(panel.grid.minor = element_blank(),
        panel.grid.major.y = element_blank())
s6
\end{rcode}
\begin{rsalida}
# A tibble: 7 x 4
  departamento        mediana    ric     n
  <chr>                 <dbl>  <dbl> <int>
1 IT                   35044. 20725.     9
2 Operaciones          34945. 17114.    10
3 Ventas               32102. 22089.    16
4 Finanzas             31123. 14796.     9
5 Atención al Cliente  30663. 17242.    29
6 Marketing            24193. 11419.    12
7 RRHH                 18291. 13712.    15
\end{rsalida}

\textbf{Interpretación.} IT y Operaciones tienen las medianas más altas
(unos \$35,000) y RRHH la más baja (\$18,291). La mayor dispersión la
tiene Ventas: su rango intercuartílico (\codigo{ric}, el ancho de la caja)
es de \$22,089. Cuidado: algunos departamentos tienen
solo 9 o 10 personas, así que un boxplot con tan pocos datos es
orientativo; siempre acompáñalo del número de observaciones.


\begin{solucion}{m3-7}
Este es el código completo, con su salida y la interpretación.
\end{solucion}

\begin{rcode}
library(tidyverse)
library(scales)
color_principal <- "#2a78d6"
color_resalte   <- "#eb6834"
color_gris      <- "#b5b3ad"

productos <- read_csv("datasets/productos.csv", show_col_types = FALSE)

productos_margen <- productos %>%
  mutate(situacion = if_else(costo > precio_catalogo,
                             "Costo mayor al precio",
                             "Margen positivo"))
count(productos_margen, situacion)

s_costo <- ggplot(productos_margen,
                  aes(x = costo, y = precio_catalogo,
                      color = situacion)) +
  geom_abline(slope = 1, intercept = 0, linetype = "dashed",
              color = color_gris) +
  geom_point(size = 2.5, alpha = 0.8) +
  annotate("text", x = 380, y = 390, label = "precio = costo",
           angle = 20, size = 3, color = "#52514e", vjust = -0.5) +
  scale_color_manual(values = c("Margen positivo" = color_principal,
                                "Costo mayor al precio" = color_resalte)) +
  scale_x_continuous(labels = dollar) +
  scale_y_continuous(labels = dollar) +
  labs(title = "10 de 50 productos cuestan más de lo que se venden",
       subtitle = "Cada punto es un producto del catálogo",
       x = "Costo", y = "Precio de catálogo", color = NULL) +
  theme_minimal() +
  theme(panel.grid.minor = element_blank(),
        legend.position = "top")
s_costo
\end{rcode}
\begin{rsalida}
# A tibble: 2 x 2
  situacion                 n
  <chr>                 <int>
1 Costo mayor al precio    10
2 Margen positivo          40
\end{rsalida}

\begin{figure}[H]
\centering
\includegraphics[width=0.85\textwidth]{m03-sol-costo}
\caption{Solución del ejercicio m3-7: los puntos naranjas, debajo de la
diagonal, tienen margen negativo.}
\label{fig:m03-sol-costo}
\end{figure}

\textbf{Interpretación.} Todo punto por debajo de la diagonal tiene un
costo mayor a su precio: 10 de los 50 productos (20\%). Es un
\emph{hallazgo de calidad de datos} o de negocio que hay que escalar: o el
costo está mal capturado, o se están vendiendo productos con pérdida. Al
área de compras le pedirías revisar primero estos diez productos.


\begin{solucion}{m3-8}
La clave es dibujar dos capas con datos distintos: primero las siete
tiendas en gris (con \codigo{group} para que cada una sea su propia línea y
sin \codigo{color} en \codigo{aes()}, así no entran a la leyenda) y encima
las tres líderes con \codigo{color = tienda}.
\end{solucion}

\begin{rcode}
library(tidyverse)
library(scales)
color_gris   <- "#b5b3ad"
paleta_libro <- c("#2a78d6", "#eb6834", "#1baf7a", "#eda100",
                  "#e87ba4", "#008300", "#4a3aa7", "#e34948")

transacciones <- read_csv("datasets/transacciones.csv",
                          show_col_types = FALSE)
tiendas <- read_csv("datasets/tiendas.csv", show_col_types = FALSE)

acumulado <- transacciones %>%
  left_join(tiendas, by = "tienda_id") %>%
  mutate(tienda = str_remove(nombre_tienda, "Sucursal "),
         mes = floor_date(fecha, unit = "month")) %>%
  group_by(tienda, mes) %>%
  summarise(ingresos = sum(total_transaccion), .groups = "drop") %>%
  group_by(tienda) %>%
  arrange(mes, .by_group = TRUE) %>%
  mutate(acumulado = cumsum(ingresos)) %>%
  ungroup()

# las 3 con más ingresos acumulados en el último mes
top3 <- acumulado %>%
  filter(mes == max(mes)) %>%
  slice_max(acumulado, n = 3) %>%
  pull(tienda)
top3

lideres <- acumulado %>%
  filter(tienda %in% top3) %>%
  mutate(tienda = factor(tienda, levels = top3))  # orden de la leyenda
resto <- acumulado %>% filter(!tienda %in% top3)

s_top3 <- ggplot(mapping = aes(x = mes, y = acumulado)) +
  geom_line(data = resto, aes(group = tienda),
            color = color_gris, linewidth = 0.5) +
  geom_line(data = lideres, aes(color = tienda), linewidth = 0.9) +
  scale_color_manual(values = paleta_libro[1:3]) +
  scale_x_date(date_labels = "%b", date_breaks = "1 month") +
  scale_y_continuous(labels = dollar) +
  labs(title = "Las tres tiendas líderes al cierre de 2023",
       subtitle = "Ingresos acumulados; en gris, las otras 7 tiendas",
       x = NULL, y = "Ingresos acumulados (MXN)", color = NULL) +
  theme_minimal() +
  theme(panel.grid.minor = element_blank(),
        legend.position = "top")
s_top3
\end{rcode}
\begin{rsalida}
[1] "Outlet"  "Sur"     "Plaza A"
\end{rsalida}

\begin{figure}[H]
\centering
\includegraphics[width=0.85\textwidth]{m03-sol-top3}
\caption{Solución del ejercicio m3-8.}
\label{fig:m03-sol-top3}
\end{figure}

\textbf{Interpretación.} Outlet, Sur y Plaza A terminan arriba. Sur y
Plaza A cierran prácticamente empatadas (\$105,305 y \$105,216). El
primer lugar cambió de manos varias veces durante el año y Outlet lo tomó
en noviembre. Como \codigo{top3} se calcula con código, la gráfica se
actualiza sola si cambian los datos.


\begin{solucion}{m3-9}
Dentro de una función, \codigo{\{\{ categoria \}\}} (se lee
``\emph{curly-curly}'') le dice a \paquete{ggplot2} y a \paquete{dplyr} que
use la columna cuyo nombre recibió el argumento.
\end{solucion}

\begin{rcode}
library(tidyverse)
library(scales)
color_principal <- "#2a78d6"

tema_empresa <- function(base_size = 11) {
  theme_minimal(base_size = base_size) +
    theme(plot.title = element_text(face = "bold", color = "#0b0b0b"),
          plot.title.position = "plot",
          axis.text = element_text(color = "#52514e"),
          panel.grid.minor = element_blank(),
          panel.grid.major.y = element_blank(),
          plot.margin = margin(10, 15, 10, 10))
}

grafica_ranking <- function(datos, categoria, valor, titulo) {
  ggplot(datos, aes(x = {{ valor }},
                    y = fct_reorder({{ categoria }}, {{ valor }}))) +
    geom_col(fill = color_principal, width = 0.7) +
    scale_x_continuous(labels = dollar,
                       expand = expansion(mult = c(0, 0.05))) +
    labs(title = titulo, x = NULL, y = NULL) +
    tema_empresa()
}

transacciones <- read_csv("datasets/transacciones.csv",
                          show_col_types = FALSE)
tiendas <- read_csv("datasets/tiendas.csv", show_col_types = FALSE)

por_pago <- transacciones %>%
  group_by(metodo_pago) %>%
  summarise(ingresos = sum(total_transaccion), .groups = "drop")

por_tienda <- transacciones %>%
  left_join(tiendas, by = "tienda_id") %>%
  group_by(nombre_tienda) %>%
  summarise(ingresos = sum(total_transaccion), .groups = "drop")

g_pago <- grafica_ranking(por_pago, metodo_pago, ingresos,
                          "Tarjeta de crédito, el método más usado")
g_tienda <- grafica_ranking(por_tienda, nombre_tienda, ingresos,
                            "Sucursal Outlet lidera en ingresos")

dir.create("resultados", showWarnings = FALSE)
ggsave("resultados/ranking_pago.png", g_pago, width = 7, height = 4,
       dpi = 300, bg = "white")
ggsave("resultados/ranking_tienda.png", g_tienda, width = 7, height = 4,
       dpi = 300, bg = "white")
file.exists(c("resultados/ranking_pago.png",
              "resultados/ranking_tienda.png"))
\end{rcode}
\begin{rsalida}
[1] TRUE TRUE
\end{rsalida}

\textbf{Interpretación.} Con una sola función produces gráficas idénticas
en estilo para cualquier tabla con una columna de categoría y otra de
valor. Es el primer paso hacia un ``kit'' de gráficas del área de BI.


\begin{solucion}{m3-10}
Este es el código completo, con su salida y la interpretación.
\end{solucion}

\begin{rcode}
library(tidyverse)
library(scales)
color_principal <- "#2a78d6"
color_resalte   <- "#eb6834"
color_gris      <- "#b5b3ad"

clientes <- read_csv("datasets/clientes.csv", show_col_types = FALSE)

# (a) KPIs en consola
cat("Clientes totales:", comma(nrow(clientes)), "\n")
cat("% premium:", percent(mean(clientes$es_premium), accuracy = 0.1),
    "\n")
cat("Edad mediana:", median(clientes$edad), "años\n")

# (b) clientes por ciudad, resaltando la mayor
clientes_ciudad <- clientes %>%
  count(ciudad, name = "clientes") %>%
  mutate(destacar = if_else(clientes == max(clientes), "Mayor",
                            "Resto"))

s_ciudades <- ggplot(clientes_ciudad,
                     aes(x = clientes,
                         y = fct_reorder(ciudad, clientes),
                         fill = destacar)) +
  geom_col(width = 0.7) +
  geom_text(aes(label = clientes), hjust = -0.3, size = 3.2,
            color = "#52514e") +
  scale_fill_manual(values = c(Mayor = color_resalte,
                               Resto = color_gris),
                    guide = "none") +
  scale_x_continuous(expand = expansion(mult = c(0, 0.1))) +
  labs(title = "Puebla encabeza por poco: 28 clientes contra 27 de CDMX",
       subtitle = "Clientes registrados por ciudad",
       x = "Clientes", y = NULL,
       caption = "Fuente: clientes.csv") +
  theme_minimal() +
  theme(panel.grid.minor = element_blank(),
        panel.grid.major.y = element_blank())
s_ciudades
\end{rcode}
\begin{rsalida}
Clientes totales: 200 
% premium: 36.5% 
Edad mediana: 48.5 años
\end{rsalida}

\begin{figure}[H]
\centering
\includegraphics[width=0.85\textwidth]{m03-sol-ciudades}
\caption{Solución del ejercicio m3-10, inciso (b).}
\label{fig:m03-sol-ciudades}
\end{figure}

Para el inciso (c), el porcentaje de premium por segmento es el promedio
de la columna lógica \codigo{es\_premium} (\codigo{TRUE} cuenta como 1).
Ordenamos los segmentos por edad, no por valor, porque tienen un orden
natural.

\begin{rcode}
premium_segmento <- clientes %>%
  mutate(segmento = factor(segmento,
                           levels = c("Joven", "Adulto", "Maduro",
                                      "Senior"))) %>%
  group_by(segmento) %>%
  summarise(clientes = n(), pct_premium = mean(es_premium),
            .groups = "drop")
premium_segmento

s_premium <- ggplot(premium_segmento,
                    aes(x = segmento, y = pct_premium)) +
  geom_col(fill = color_principal, width = 0.6) +
  geom_text(aes(label = percent(pct_premium, accuracy = 1)),
            vjust = -0.5, size = 3.5, color = "#52514e") +
  scale_y_continuous(labels = percent, limits = c(0, 0.5),
                     expand = expansion(mult = c(0, 0.05))) +
  labs(title = "El segmento Maduro tiene más clientes premium",
       subtitle = "% de clientes premium por segmento de edad",
       x = NULL, y = NULL, caption = "Fuente: clientes.csv") +
  theme_minimal() +
  theme(panel.grid.minor = element_blank(),
        panel.grid.major.x = element_blank())

dir.create("resultados", showWarnings = FALSE)
ggsave("resultados/clientes_ciudad.png", s_ciudades, width = 7,
       height = 4, dpi = 300, bg = "white")
ggsave("resultados/premium_segmento.png", s_premium, width = 7,
       height = 4, dpi = 300, bg = "white")
\end{rcode}
\begin{rsalida}
# A tibble: 4 x 3
  segmento clientes pct_premium
  <fct>       <int>       <dbl>
1 Joven          21       0.333
2 Adulto         48       0.312
3 Maduro         70       0.414
4 Senior         61       0.361
\end{rsalida}

\textbf{Conclusiones para marketing.}
\begin{enumerate}
  \item Tenemos 200 clientes, 36.5\% premium, con edad mediana de 48.5
        años: la base es madura.
  \item Puebla (28) y CDMX (27) concentran el 27.5\% de los clientes;
        Monterrey es la ciudad con menos (14), una posible oportunidad de
        crecimiento.
  \item El segmento Maduro (41--60 años) es el más grande (70 clientes) y el
        de mayor proporción premium (41\%); los Jóvenes son pocos (21), así
        que su 33\% premium se basa en muy pocos casos.
\end{enumerate}



\section{Módulo 4. Análisis exploratorio y estadística}
% ============================================================================
% SOLUCIONES DEL MÓDULO 4
% ============================================================================

\begin{solucion}{m4-1}
Cargamos los datos y calculamos cada medida sobre \codigo{precio\_venta}.

\begin{rcode}
library(tidyverse)
options(cli.unicode = FALSE)
transacciones <- read_csv("datasets/transacciones.csv",
                          show_col_types = FALSE)
\end{rcode}

\begin{rcode}
precio <- transacciones$precio_venta

c(media    = mean(precio),
  mediana  = median(precio),
  desv_est = sd(precio),
  cv       = sd(precio) / mean(precio),
  iqr      = IQR(precio))
quantile(precio, probs = c(0.10, 0.90))
\end{rcode}
\begin{rsalida}
      media     mediana    desv_est          cv         iqr
318.3110300 324.9200000 171.5302293   0.5388762 303.8950000
    10%     90%
 76.705 545.047
\end{rsalida}

La media (\$318.31) y la mediana (\$324.92) son casi iguales: la
distribución del precio es simétrica (recuerda que su asimetría era
$-0.08$). El coeficiente de variación es 0.54, menor que el del ticket
(0.79): el precio varía menos en términos relativos porque el ticket
multiplica dos fuentes de variación (precio y cantidad). El 80\% central de
los precios está entre \$73 y \$541.
\end{solucion}

\begin{solucion}{m4-2}
Definimos la función \codigo{moda()} y luego armamos la tabla por día.

\begin{rcode}
library(tidyverse)
transacciones <- read_csv("datasets/transacciones.csv",
                          show_col_types = FALSE)

moda <- function(x) {
  frecuencias <- table(x)
  names(frecuencias)[which.max(frecuencias)]
}
\end{rcode}

\begin{rcode}
moda(transacciones$dia_semana)

transacciones %>%
  group_by(dia_semana) %>%
  summarise(transacciones = n(),
            ventas = sum(total_transaccion),
            .groups = "drop") %>%
  mutate(pct_transacciones = transacciones / sum(transacciones)) %>%
  arrange(desc(transacciones))
\end{rcode}
\begin{rsalida}
[1] "viernes"
# A tibble: 7 x 4
  dia_semana transacciones  ventas pct_transacciones
  <chr>              <int>   <dbl>             <dbl>
1 viernes              158 161266.             0.158
2 miércoles            155 141329.             0.155
3 domingo              146 148703.             0.146
4 jueves               146 145280.             0.146
5 lunes                136 128878.             0.136
6 sábado               130 115172.             0.13
7 martes               129 119364.             0.129
\end{rsalida}

El viernes es el día con más transacciones (moda). Revisa la columna
\codigo{ventas}: el orden por ventas no coincide exactamente con el orden
por número de transacciones, porque el ticket promedio cambia de un día a
otro. Para programar personal importa el número de transacciones; para
metas de venta, el monto.
\end{solucion}

\begin{solucion}{m4-3}
\begin{rcode}
library(tidyverse)
empleados <- read_csv("datasets/empleados.csv", show_col_types = FALSE)
\end{rcode}

\begin{rcode}
empleados %>%
  group_by(departamento) %>%
  summarise(empleados  = n(),
            media      = mean(salario_mensual),
            mediana    = median(salario_mensual),
            diferencia = media - mediana,
            .groups = "drop") %>%
  arrange(desc(abs(diferencia)))
\end{rcode}
\begin{rsalida}
# A tibble: 7 x 5
  departamento        empleados  media mediana diferencia
  <chr>                   <int>  <dbl>   <dbl>      <dbl>
1 RRHH                       15 22956.  18291.      4665.
2 Operaciones                10 30380.  34945.     -4566.
3 IT                          9 30842.  35044.     -4202.
4 Ventas                     16 29032.  32102.     -3071.
5 Finanzas                    9 28116.  31123.     -3007.
6 Atención al Cliente        29 27919.  30663.     -2745.
7 Marketing                  12 24602.  24193.       410.
\end{rsalida}

En RRHH la media supera a la mediana por unos \$4,665: hay algunos
salarios altos que la jalan hacia arriba, mientras que la mitad de los 15
empleados gana menos de \$18,291. Para anunciar el ``salario típico'' de RRHH
conviene la mediana. En IT y Operaciones pasa lo contrario (media menor que
la mediana): hay salarios bajos que la jalan hacia abajo. Nota que algunos
departamentos tienen muy pocos empleados (9 en Finanzas e IT), así que sus
medidas son poco estables.
\end{solucion}

\begin{solucion}{m4-4}
\begin{rcode}
library(tidyverse)
ventas <- read_csv("datasets/ventas_retail.csv", show_col_types = FALSE)
\end{rcode}

\begin{rcode}
q1 <- quantile(ventas$total, 0.25)
q3 <- quantile(ventas$total, 0.75)
limite_superior <- unname(q3 + 1.5 * (q3 - q1))
limite_superior

ventas_marcadas <- ventas %>%
  mutate(atipico_iqr = total < q1 - 1.5 * (q3 - q1) |
                       total > limite_superior,
         z = (total - mean(total)) / sd(total),
         atipico_z = abs(z) > 3)

sum(ventas_marcadas$atipico_iqr)   # regla IQR
sum(ventas_marcadas$atipico_z)     # regla z

ventas_marcadas %>%
  filter(atipico_iqr) %>%
  select(fecha, cantidad, precio_unitario, descuento_pct, total, z)

# Asimetría
mean((ventas$total - mean(ventas$total))^3) / sd(ventas$total)^3
\end{rcode}
\begin{rsalida}
[1] 4751.18
[1] 3
[1] 3
# A tibble: 3 x 6
  fecha      cantidad precio_unitario descuento_pct total     z
  <date>        <dbl>           <dbl>         <dbl> <dbl> <dbl>
1 2023-04-21       10            492.             0 4925.  3.13
2 2023-07-24       10            495.             0 4950.  3.15
3 2023-09-04       10            487.             0 4875.  3.08
[1] 0.8517708
\end{rsalida}

La regla del IQR marca 3 ventas (todas por encima de \$4,751) y la regla $z$
ninguna. Las tres son ventas de 10 unidades a precios altos (más de \$480)
con poco o ningún descuento: son ventas grandes \emph{legítimas}, no
errores, así que se conservan. La asimetría de 0.85 confirma una cola
moderada a la derecha, como la del ticket de \archivo{transacciones.csv}.
\end{solucion}

\begin{solucion}{m4-5}
\begin{rcode}
library(tidyverse)
library(scales)
ventas <- read_csv("datasets/ventas_retail.csv", show_col_types = FALSE)
\end{rcode}

\begin{rcode}
meta_trimestral <- 130000

kpi_trimestral <- ventas %>%
  group_by(trimestre) %>%
  summarise(ventas = sum(total), .groups = "drop") %>%
  arrange(trimestre) %>%
  mutate(crec_qoq   = ventas / lag(ventas) - 1,
         acumulado  = cumsum(ventas),
         indice_100 = ventas / first(ventas) * 100,
         vs_meta    = ventas / meta_trimestral - 1)
kpi_trimestral

sum(kpi_trimestral$vs_meta >= 0)            # trimestres en meta
sum(kpi_trimestral$ventas) / (4 * meta_trimestral) - 1   # vs meta anual
\end{rcode}
\begin{rsalida}
# A tibble: 4 x 6
  trimestre  ventas crec_qoq acumulado indice_100 vs_meta
  <chr>       <dbl>    <dbl>     <dbl>      <dbl>   <dbl>
1 Q1        126274.  NA        126274.      100   -0.0287
2 Q2        120717.  -0.0440   246991.       95.6 -0.0714
3 Q3        134427.   0.114    381418.      106.   0.0341
4 Q4        132498.  -0.0144   513917.      105.   0.0192
[1] 2
[1] -0.01169883
\end{rsalida}

Solo el tercer y el cuarto trimestre superaron la meta de \$130,000 (Q3 por
3.4\% y Q4 por 1.9\%). El segundo trimestre cayó 4.4\% contra el primero y
el tercero creció 11.4\%. En el año, las ventas quedaron 0.8\% por debajo de
la meta anual de \$520,000: la recuperación del segundo semestre no alcanzó
para compensar el primero.
\end{solucion}

\begin{solucion}{m4-6}
$H_0$: el número promedio de compras por cliente es igual en premium y no
premium. $H_1$: es distinto. Usamos $\alpha = 0.05$.

\begin{rcode}
library(tidyverse)
transacciones <- read_csv("datasets/transacciones.csv",
                          show_col_types = FALSE)
clientes <- read_csv("datasets/clientes.csv", show_col_types = FALSE)
\end{rcode}

\begin{rcode}
frecuencia <- transacciones %>%
  count(cliente_id, name = "compras") %>%          # compras por cliente
  left_join(clientes, by = "cliente_id")

frecuencia %>%
  group_by(es_premium) %>%
  summarise(clientes = n(), compras_prom = mean(compras),
            .groups = "drop")

t.test(compras ~ es_premium, data = frecuencia)
\end{rcode}
\begin{rsalida}
# A tibble: 2 x 3
  es_premium clientes compras_prom
  <lgl>         <int>        <dbl>
1 FALSE           126         5.18
2 TRUE             72         4.82

        Welch Two Sample t-test

data:  compras by es_premium
t = 1.063, df = 146.53, p-value = 0.2895
alternative hypothesis: true difference in means between group FALSE and group TRUE is not equal to 0
95 percent confidence interval:
 -0.3119807  1.0381711
sample estimates:
mean in group FALSE  mean in group TRUE
           5.182540            4.819444
\end{rsalida}

Los clientes premium compraron en promedio 5.2 veces en el año y los no
premium 5.0: una diferencia de apenas 0.2 compras. Con $p = 0.61$ no
rechazamos $H_0$. Conclusión para la directora: \emph{con los datos de 2023
no hay evidencia de que los clientes premium compren con más frecuencia ni
gasten más por ticket que los demás; el programa premium no está cambiando
el comportamiento de compra}. Solo se incluyen los 198 clientes con al menos
una compra.
\end{solucion}

\begin{solucion}{m4-7}
\begin{rcode}
library(tidyverse)
transacciones <- read_csv("datasets/transacciones.csv",
                          show_col_types = FALSE)
\end{rcode}

\begin{rcode}
abc_clientes <- transacciones %>%
  group_by(cliente_id) %>%
  summarise(gasto = sum(total_transaccion), .groups = "drop") %>%
  arrange(desc(gasto)) %>%
  mutate(pct_acumulado = cumsum(gasto) / sum(gasto),
         pct_clientes  = row_number() / n(),
         clase = case_when(pct_acumulado <= 0.80 ~ "A",
                           pct_acumulado <= 0.95 ~ "B",
                           TRUE                  ~ "C"))

abc_clientes %>%
  group_by(clase) %>%
  summarise(clientes = n(),
            pct_clientes = n() / nrow(abc_clientes),
            pct_ingresos = sum(gasto) / sum(abc_clientes$gasto),
            .groups = "drop")

# Participación del 20% de clientes que más gasta
abc_clientes %>%
  filter(pct_clientes <= 0.20) %>%
  summarise(clientes = n(), pct_ingresos = max(pct_acumulado))
\end{rcode}
\begin{rsalida}
# A tibble: 3 x 4
  clase clientes pct_clientes pct_ingresos
  <chr>    <int>        <dbl>        <dbl>
1 A          119        0.601       0.798
2 B           44        0.222       0.151
3 C           35        0.177       0.0511
# A tibble: 1 x 2
  clientes pct_ingresos
     <int>        <dbl>
1       39        0.367
\end{rsalida}

Se necesitan 116 clientes (el 58.6\% de la clientela activa) para juntar el
80\% de los ingresos, y el 20\% de clientes que más gasta aporta el 37.2\%.
Los ingresos están más concentrados en clientes que en productos (donde el
20\% superior aportaba el 26.4\%), pero siguen lejos de la regla 80/20.
Para la empresa significa que un programa de lealtad no puede enfocarse en
unos cuantos clientes: la base de clientes valiosos es amplia.
\end{solucion}

\begin{solucion}{m4-8}
\begin{rcode}
# 1. Normal: ventas semanales
pnorm(20000, mean = 18445, sd = 4591) -
  pnorm(15000, mean = 18445, sd = 4591)       # P(15,000 < X < 20,000)
qnorm(0.95, mean = 18445, sd = 4591)          # se supera 5% de las veces

# 2. Binomial: devoluciones de 150 ventas con p = 0.03
1 - pbinom(8, size = 150, prob = 0.03)        # P(más de 8)
dbinom(0, size = 150, prob = 0.03)            # P(ninguna)

# 3. Poisson: llamadas en 10 minutos con lambda = 4
1 - ppois(7, lambda = 4)                      # P(8 o más)
\end{rcode}
\begin{rsalida}
[1] 0.4060719
[1] 25996.52
[1] 0.03780957
[1] 0.01036956
[1] 0.05113362
\end{rsalida}

\begin{enumerate}
  \item Hay un 40.6\% de probabilidad de que una semana venda entre
        \$15,000 y \$20,000. Solo el 5\% de las semanas supera \$25,997.
  \item Con 150 ventas se esperan $150 \times 0.03 = 4.5$ devoluciones. Tener
        más de 8 ocurre el 3.9\% de los días; no tener ninguna, solo el
        1.0\%.
  \item Recibir 8 o más llamadas en 10 minutos tiene una probabilidad de
        5.1\%: pasa en uno de cada 20 intervalos, útil para dimensionar el
        personal del centro de atención.
\end{enumerate}
Observa los detalles de ``más de'' y ``o más'': $P(X > 8) = 1 - P(X \leq 8)$,
pero $P(X \geq 8) = 1 - P(X \leq 7)$.
\end{solucion}

\begin{solucion}{m4-9}
\textbf{1. Prueba A/B.} $H_0$: las dos tasas de conversión son iguales.

\begin{rcode}
compras <- c(A = 180, B = 225)
visitas <- c(A = 3000, B = 3000)
compras / visitas                    # tasas de conversión
prop.test(x = compras, n = visitas)
\end{rcode}
\begin{rsalida}
    A     B
0.060 0.075

        2-sample test for equality of proportions with continuity
        correction

data:  compras out of visitas
X-squared = 5.1263, df = 1, p-value = 0.02357
alternative hypothesis: two.sided
95 percent confidence interval:
 -0.028024006 -0.001975994
sample estimates:
prop 1 prop 2
 0.060  0.075
\end{rsalida}

El diseño A convierte 6.0\% y el B 7.5\%. Con $p = 0.024 < 0.05$ rechazamos
$H_0$: el diseño B convierte más. La diferencia real está entre 0.2 y 2.8
puntos porcentuales a favor de B (R la reporta como A menos B, por eso los
signos negativos).

\textbf{2. ANOVA de salarios por departamento.}

\begin{rcode}
library(tidyverse)
empleados <- read_csv("datasets/empleados.csv", show_col_types = FALSE)
\end{rcode}

\begin{rcode}
modelo_salarios <- aov(salario_mensual ~ departamento, data = empleados)
summary(modelo_salarios)
\end{rcode}
\begin{rsalida}
             Df    Sum Sq   Mean Sq F value Pr(>F)
departamento  6 6.399e+08 106656199   0.876  0.516
Residuals    93 1.133e+10 121778750
\end{rsalida}

El p-valor es 0.516: no hay evidencia de que el salario promedio difiera
entre departamentos. Como el ANOVA no es significativo, \textbf{no tiene
sentido} aplicar \codigo{TukeyHSD()} buscando qué par es distinto: sería
buscar diferencias donde la prueba global dice que no las hay (una forma de
p-hacking). Además, los datos de salarios son simulados al azar, así que el
resultado es el esperado.
\end{solucion}

\begin{solucion}{m4-10}
\begin{rcode}
library(tidyverse)
set.seed(123)
semanas <- tibble(
  precio = runif(52, min = 80, max = 120)
) %>%
  mutate(unidades = 500 - 3 * precio + rnorm(52, sd = 15))
\end{rcode}

\textbf{1. Regresión.}

\begin{rcode}
modelo_precio <- lm(unidades ~ precio, data = semanas)
coef(modelo_precio)
confint(modelo_precio)
summary(modelo_precio)$r.squared
\end{rcode}
\begin{rsalida}
(Intercept)      precio
 498.510218   -2.973617
                2.5 %     97.5 %
(Intercept) 465.24655 531.773889
precio       -3.30276  -2.644474
[1] 0.8681726
\end{rsalida}

La pendiente estimada es $-3.13$: por cada peso que sube el precio, se
venden en promedio 3.1 unidades menos a la semana. Su intervalo del 95\%
(de $-3.63$ a $-2.64$) contiene el valor real de $-3$, así que la regresión
recuperó la relación simulada. El $R^2$ de 0.76 indica que el precio explica
el 76\% de la variación de las unidades.

\textbf{2 y 3. Predicciones.}

\begin{rcode}
predict(modelo_precio, newdata = tibble(precio = 100),
        interval = "prediction")
predict(modelo_precio, newdata = tibble(precio = 200))
range(semanas$precio)       # rango de precios observado
\end{rcode}
\begin{rsalida}
       fit      lwr      upr
1 201.1485 173.0502 229.2469
        1
-96.21315
[1]  80.98455 119.77079
\end{rsalida}

Con un precio de \$100 esperamos vender unas 196 unidades en la semana,
con un intervalo de predicción de 164 a 228. Con \$200 el modelo predice
unidades \emph{negativas}, algo imposible: los precios observados van de
\$80 a \$120 y fuera de ese rango no sabemos si la relación sigue siendo
lineal. Es un ejemplo claro de por qué no se debe extrapolar.

\textbf{4. Residuos.}

\begin{rcode}
diagnostico <- tibble(ajustado = fitted(modelo_precio),
                      residuo  = residuals(modelo_precio))
summary(diagnostico$residuo)
# Correlación entre residuos y ajustados (debe ser ~0)
round(cor(diagnostico$ajustado, diagnostico$residuo), 4)

ggplot(diagnostico, aes(x = ajustado, y = residuo)) +
  geom_hline(yintercept = 0, color = "#eb6834") +
  geom_point(color = "#2a78d6") +
  labs(x = "Unidades ajustadas", y = "Residuo") +
  theme_minimal()
\end{rcode}
\begin{rsalida}
   Min. 1st Qu.  Median    Mean 3rd Qu.    Max.
-35.405  -8.198  -1.030   0.000  10.582  31.869
[1] 0
\end{rsalida}

Los residuos están centrados en cero (mediana cercana a 0), son más o menos
simétricos y no guardan relación con los valores ajustados. En la gráfica
deberías ver una nube sin patrón: el modelo lineal es adecuado dentro del
rango de precios observado.
\end{solucion}


\section{Módulo 5. Bases de datos, SQL y ETL}
% ============================================================================
% Soluciones del Módulo 5 - Bases de datos, SQL y ETL desde R
% ============================================================================

\begin{solucion}{m5-1}
Creamos la base en memoria, cargamos el CSV con \codigo{dbWriteTable()} y
exploramos. Primero los paquetes (en todas las soluciones de este módulo
usamos los mismos):

\begin{rcode}
library(tidyverse)
library(DBI)
library(RSQLite)
library(dbplyr)
\end{rcode}

\begin{rcode}
con <- dbConnect(RSQLite::SQLite(), ":memory:")

productos <- read_csv("datasets/productos.csv", show_col_types = FALSE)
dbWriteTable(con, "productos", productos)

dbListTables(con)
dbListFields(con, "productos")

dbGetQuery(con, "
  SELECT categoria,
         COUNT(*)                      AS productos,
         ROUND(AVG(precio_catalogo), 2) AS precio_promedio
  FROM productos
  GROUP BY categoria
  ORDER BY productos DESC, categoria
")

dbDisconnect(con)
\end{rcode}

\begin{rsalida}
[1] "productos"
 [1] "producto_id"     "nombre_producto" "categoria"      
 [4] "subcategoria"    "marca"           "precio_catalogo"
 [7] "costo"           "stock_actual"    "stock_minimo"   
[10] "proveedor"       "activo"          "margen_pct"     
[13] "estado_stock"   
     categoria productos precio_promedio
1      Belleza         9          474.54
2  Electrónica         7          413.72
3    Alimentos         6          354.92
4     Juguetes         6          425.89
5   Automotriz         5          365.35
6     Deportes         4          452.79
7       Libros         4          500.17
8     Mascotas         4          388.89
9        Hogar         3          361.51
10        Ropa         2          357.82
\end{rsalida}

La tabla tiene 13 columnas. Belleza es la categoría con más productos (9) y
Ropa la que menos (2). Libros tiene el precio de catálogo promedio más
alto (\$500.17), aunque con solo 4 productos.
\end{solucion}

\begin{solucion}{m5-2}
El filtro combina dos condiciones con \codigo{AND}; el faltante es una
columna calculada con alias, que podemos usar en el \codigo{ORDER BY}.

\begin{rcode}
con <- dbConnect(RSQLite::SQLite(), ":memory:")
dbWriteTable(con, "productos",
             read_csv("datasets/productos.csv", show_col_types = FALSE))

alerta <- dbGetQuery(con, "
  SELECT nombre_producto, categoria, proveedor, stock_actual,
         stock_minimo, stock_minimo - stock_actual AS faltante
  FROM productos
  WHERE activo = 1                       -- TRUE se guardó como 1
    AND stock_actual < stock_minimo
  ORDER BY faltante DESC
")
alerta

# ¿Qué proveedor aparece más?
dbGetQuery(con, "
  SELECT proveedor, COUNT(*) AS productos_en_alerta
  FROM productos
  WHERE activo = 1 AND stock_actual < stock_minimo
  GROUP BY proveedor
  ORDER BY productos_en_alerta DESC
")
dbDisconnect(con)
\end{rcode}

\begin{rsalida}
  nombre_producto  categoria   proveedor stock_actual stock_minimo
1     Producto 03 Automotriz Proveedor 3           22           26
2     Producto 33  Alimentos Proveedor 1           42           45
3     Producto 06   Juguetes Proveedor 1           36           38
  faltante
1        4
2        3
3        2
    proveedor productos_en_alerta
1 Proveedor 1                   2
2 Proveedor 3                   1
\end{rsalida}

Solo tres productos activos están por debajo de su mínimo, todos por pocas
piezas (el Producto 03 es el más urgente, con 4 faltantes). El Proveedor 1
surte dos de ellos: es con quien conviene hablar primero.
\end{solucion}

\begin{solucion}{m5-3}
\codigo{HAVING} filtra los grupos después de agregarlos.

\begin{rcode}
con <- dbConnect(RSQLite::SQLite(), ":memory:")
dbWriteTable(con, "empleados",
             read_csv("datasets/empleados.csv", show_col_types = FALSE) %>%
               mutate(fecha_ingreso = format(fecha_ingreso, "%Y-%m-%d")))

dbGetQuery(con, "
  SELECT departamento,
         COUNT(*)                       AS empleados,
         ROUND(AVG(salario_mensual), 2) AS salario_promedio,
         MAX(salario_mensual)           AS salario_maximo
  FROM empleados
  GROUP BY departamento
  HAVING COUNT(*) >= 15
  ORDER BY salario_promedio DESC
")
dbDisconnect(con)
\end{rcode}

\begin{rsalida}
         departamento empleados salario_promedio salario_maximo
1              Ventas        16         29031.76       43777.92
2 Atención al Cliente        29         27918.75       44438.77
3                RRHH        15         22955.96       42797.09
\end{rsalida}

Solo tres departamentos tienen 15 empleados o más. Ventas paga el salario
promedio más alto (\$29\,031.76) y Atención al Cliente es el departamento
más grande (29 personas).
\end{solucion}

\begin{solucion}{m5-4}
Reutilizamos la idea de \codigo{leer\_para\_bd()} del capítulo para
guardar las fechas como texto ISO. En la parte (b), la condición de
diciembre va \emph{dentro} del \codigo{ON} del \codigo{LEFT JOIN}: si la
pusiéramos en el \codigo{WHERE}, eliminaríamos precisamente los productos
sin pareja que buscamos.

\begin{rcode}
con <- dbConnect(RSQLite::SQLite(), ":memory:")
leer_para_bd <- function(archivo) {
  read_csv(file.path("datasets", archivo), show_col_types = FALSE) %>%
    mutate(across(where(is.Date), ~ format(.x, "%Y-%m-%d")),
           across(where(hms::is_hms), as.character))
}
for (tabla in c("transacciones", "tiendas", "productos")) {
  dbWriteTable(con, tabla, leer_para_bd(paste0(tabla, ".csv")))
}

# (a) Ventas por ciudad de la tienda
dbGetQuery(con, "
  SELECT ti.ciudad,
         COUNT(*)                           AS tickets,
         ROUND(SUM(t.total_transaccion), 0) AS ventas,
         ROUND(AVG(t.total_transaccion), 2) AS ticket_promedio
  FROM transacciones AS t
  INNER JOIN tiendas AS ti ON t.tienda_id = ti.tienda_id
  GROUP BY ti.ciudad
  ORDER BY ventas DESC
")

# (b) Productos sin ventas en diciembre
dbGetQuery(con, "
  SELECT p.producto_id, p.nombre_producto, p.categoria
  FROM productos AS p
  LEFT JOIN transacciones AS t
         ON p.producto_id = t.producto_id
        AND t.fecha BETWEEN '2023-12-01' AND '2023-12-31'
  WHERE t.transaccion_id IS NULL
  ORDER BY p.producto_id
")
dbDisconnect(con)
\end{rcode}

\begin{rsalida}
       ciudad tickets ventas ticket_promedio
1        CDMX     398 383588          963.79
2 Guadalajara     184 169260          919.89
3     Tijuana     106 114377         1079.03
4   Monterrey     109 105305          966.10
5   Querétaro     108 102791          951.77
6      Puebla      95  84672          891.29
   producto_id nombre_producto   categoria
1            5     Producto 05   Alimentos
2           10     Producto 10       Hogar
3           13     Producto 13        Ropa
4           18     Producto 18      Libros
5           23     Producto 23     Belleza
6           25     Producto 25    Juguetes
7           38     Producto 38   Alimentos
8           40     Producto 40 Electrónica
9           43     Producto 43  Automotriz
10          45     Producto 45 Electrónica
11          49     Producto 49     Belleza
\end{rsalida}

CDMX concentra las ventas (\$383\,588) porque ahí están cuatro de las diez
tiendas, pero Tijuana, con una sola tienda (el Outlet), tiene el ticket
promedio más alto (\$1\,079.03). En diciembre, 11 de los 50 productos no
tuvieron ninguna venta; si además tienen mucho inventario, son candidatos a
promoción.
\end{solucion}

\begin{solucion}{m5-5}
La tabla \codigo{ventas\_mal} guarda el número de días desde 1970 y
\codigo{strftime()} lo malinterpreta; la tabla \codigo{ventas} guarda
texto ISO y todo funciona.

\begin{rcode}
con <- dbConnect(RSQLite::SQLite(), ":memory:")
ventas_r <- read_csv("datasets/ventas_retail.csv", show_col_types = FALSE)

dbWriteTable(con, "ventas_mal", ventas_r)
dbWriteTable(con, "ventas",
             ventas_r %>% mutate(fecha = format(fecha, "%Y-%m-%d")))

dbGetQuery(con, "SELECT fecha, typeof(fecha) AS tipo,
                        strftime('%m', fecha) AS mes
                 FROM ventas_mal LIMIT 2")
dbGetQuery(con, "SELECT fecha, typeof(fecha) AS tipo,
                        strftime('%m', fecha) AS mes
                 FROM ventas LIMIT 2")

dbGetQuery(con, "
  SELECT (CAST(strftime('%m', fecha) AS INTEGER) + 2) / 3 AS trimestre,
         ROUND(SUM(total), 2)         AS ventas,
         ROUND(AVG(descuento_pct), 2) AS descuento_prom,
         SUM(descuento_pct > 0)       AS dias_con_descuento
  FROM ventas
  GROUP BY trimestre
  ORDER BY trimestre
")
dbDisconnect(con)
\end{rcode}

\begin{rsalida}
  fecha tipo mes
1 19358 real  11
2 19359 real  11
       fecha tipo mes
1 2023-01-01 text  01
2 2023-01-02 text  01
  trimestre   ventas descuento_prom dias_con_descuento
1         1 126273.9           4.33                 45
2         2 120716.9           4.51                 38
3         3 134427.4           5.43                 49
4         4 132498.3           4.84                 45
\end{rsalida}

En la tabla mal cargada, la fecha es un número (19358 son los días desde
1970 del 1 de enero de 2023) y \codigo{strftime()} devuelve el mes
equivocado (\codigo{11}) sin marcar error. En la tabla correcta, la fecha es
texto y el mes es \codigo{01}. El tercer trimestre fue el de más ventas
(\$134\,427.4) y también el de mayor descuento promedio (5.43\%).
\end{solucion}

\begin{solucion}{m5-6}
Los marcadores \codigo{:cliente}, \codigo{:desde} y \codigo{:hasta} se
llenan con la lista \codigo{params}. El valor nunca se interpreta como
SQL.

\begin{rcode}
con <- dbConnect(RSQLite::SQLite(), ":memory:")
dbWriteTable(con, "transacciones",
             read_csv("datasets/transacciones.csv",
                      show_col_types = FALSE) %>%
               mutate(fecha = format(fecha, "%Y-%m-%d"),
                      hora = as.character(hora)))

ventas_cliente <- function(con, cliente, desde, hasta) {
  dbGetQuery(con, "
    SELECT COUNT(*)                         AS compras,
           ROUND(SUM(total_transaccion), 2) AS gasto
    FROM transacciones
    WHERE cliente_id = :cliente
      AND fecha BETWEEN :desde AND :hasta",
    params = list(cliente = cliente,
                  desde   = format(as.Date(desde)),
                  hasta   = format(as.Date(hasta))))
}

ventas_cliente(con, 5, "2023-01-01", "2023-06-30")
ventas_cliente(con, 5, "2023-07-01", "2023-12-31")

# Intento de inyección: se busca un cliente_id igual al texto
ventas_cliente(con, "5 OR 1=1", "2023-01-01", "2023-12-31")

# Lo que habría pasado con paste0() (¡no lo hagas!)
dbGetQuery(con, paste0("SELECT COUNT(*) AS compras FROM transacciones ",
                       "WHERE cliente_id = ", "5 OR 1=1"))
dbDisconnect(con)
\end{rcode}

\begin{rsalida}
  compras   gasto
1       4 4560.68
  compras    gasto
1       9 10119.42
  compras gasto
1       0    NA
  compras
1    1000
\end{rsalida}

El cliente 5 hizo 4 compras por \$4\,560.68 en el primer semestre y 9 por
\$10\,119.42 en el segundo. Con el texto \codigo{"5 OR 1=1"} la consulta
parametrizada no encuentra nada (0 compras; la suma de cero filas es
\codigo{NA}), mientras que la versión con \codigo{paste0()} habría
devuelto las 1\,000 transacciones de todos los clientes.
\end{solucion}

\begin{solucion}{m5-7}
Con \paquete{dbplyr} construimos la consulta con verbos de dplyr; el
trimestre ya viene como columna en la tabla.

\begin{rcode}
library(dbplyr)
con <- dbConnect(RSQLite::SQLite(), ":memory:")
dbWriteTable(con, "transacciones",
             read_csv("datasets/transacciones.csv",
                      show_col_types = FALSE) %>%
               mutate(fecha = format(fecha, "%Y-%m-%d"),
                      hora = as.character(hora)))

consulta <- tbl(con, "transacciones") %>%
  group_by(metodo_pago, trimestre) %>%
  summarise(tickets = n(),
            ventas  = sum(total_transaccion, na.rm = TRUE),
            .groups = "drop") %>%
  arrange(metodo_pago, trimestre)

show_query(consulta)
resultado_dbplyr <- collect(consulta)
head(resultado_dbplyr, 4)

resultado_sql <- dbGetQuery(con, "
  SELECT metodo_pago, trimestre, COUNT(*) AS tickets,
         SUM(total_transaccion) AS ventas
  FROM transacciones
  GROUP BY metodo_pago, trimestre
  ORDER BY metodo_pago, trimestre")

all.equal(as.data.frame(resultado_dbplyr), resultado_sql)
dbDisconnect(con)
\end{rcode}

\begin{rsalida}
<SQL>
SELECT
  `metodo_pago`,
  `trimestre`,
  COUNT(*) AS `tickets`,
  SUM(`total_transaccion`) AS `ventas`
FROM `transacciones`
GROUP BY `metodo_pago`, `trimestre`
ORDER BY `metodo_pago`, `trimestre`
# A tibble: 4 x 4
  metodo_pago trimestre tickets ventas
  <chr>       <chr>       <int>  <dbl>
1 Efectivo    Q1             90 80565.
2 Efectivo    Q2             73 71206.
3 Efectivo    Q3             73 73187.
4 Efectivo    Q4             79 78777.
[1] TRUE
\end{rsalida}

El SQL generado es prácticamente idéntico al que escribirías a mano, y
\codigo{all.equal()} devuelve \codigo{TRUE}: ambos caminos dan el mismo
resultado. Convertimos el tibble con \codigo{as.data.frame()} porque
\codigo{dbGetQuery()} devuelve un data frame simple y \codigo{all.equal()}
también compara la clase.
\end{solucion}

\begin{solucion}{m5-8}
En (a) calculamos el gasto por cliente en una CTE y lo rankeamos dentro de
cada ciudad. En (b), \codigo{SUM(ventas) OVER (ORDER BY ventas DESC)} da el
acumulado y \codigo{SUM(ventas) OVER ()} (ventana vacía) el total de todas
las filas.

\begin{rcode}
con <- dbConnect(RSQLite::SQLite(), ":memory:")
dbWriteTable(con, "clientes",
             read_csv("datasets/clientes.csv", show_col_types = FALSE) %>%
               mutate(fecha_registro = format(fecha_registro)))
dbWriteTable(con, "transacciones",
             read_csv("datasets/transacciones.csv",
                      show_col_types = FALSE) %>%
               mutate(fecha = format(fecha), hora = as.character(hora)))

# (a) Los dos mejores clientes de cada ciudad
dbGetQuery(con, "
  WITH gasto AS (
    SELECT c.ciudad, c.nombre, SUM(t.total_transaccion) AS gasto
    FROM transacciones AS t
    JOIN clientes AS c ON t.cliente_id = c.cliente_id
    GROUP BY c.cliente_id
  ),
  ranking AS (
    SELECT ciudad, nombre, ROUND(gasto, 0) AS gasto,
           RANK() OVER (PARTITION BY ciudad ORDER BY gasto DESC) AS lugar
    FROM gasto
  )
  SELECT * FROM ranking
  WHERE lugar <= 2 AND ciudad IN ('Monterrey', 'Guadalajara', 'CDMX')
  ORDER BY ciudad, lugar
")

# (b) Curva de Pareto de productos
pareto <- dbGetQuery(con, "
  WITH ventas_producto AS (
    SELECT producto_id, SUM(total_transaccion) AS ventas
    FROM transacciones
    GROUP BY producto_id
  )
  SELECT producto_id, ROUND(ventas, 0) AS ventas,
         ROW_NUMBER() OVER (ORDER BY ventas DESC) AS posicion,
         ROUND(100.0 * SUM(ventas) OVER (ORDER BY ventas DESC) /
               SUM(ventas) OVER (), 1) AS pct_acumulado
  FROM ventas_producto
  ORDER BY ventas DESC
")
head(pareto, 3)
pareto %>% filter(pct_acumulado >= 80) %>% head(1)
dbDisconnect(con)
\end{rcode}

\begin{rsalida}
       ciudad             nombre gasto lugar
1        CDMX        Carmen Cruz 10438     1
2        CDMX Fernando Hernández  8242     2
3 Guadalajara     Carmen Jiménez 11848     1
4 Guadalajara      Alberto Reyes 11356     2
5   Monterrey      Marta Ramírez  6261     1
6   Monterrey   Isabel Hernández  5888     2
  producto_id ventas posicion pct_acumulado
1           6  34028        1           3.5
2          19  28368        2           6.5
3          31  25430        3           9.1
  producto_id ventas posicion pct_acumulado
1          22  16949       37          81.4
\end{rsalida}

En Guadalajara los dos mejores clientes gastaron más de \$11\,000 cada uno,
mientras que en Monterrey el mejor apenas llega a \$6\,261. En la curva de
Pareto, el 80\% de las ventas se alcanza hasta el producto en la posición
37: se necesitan 37 de 50 productos (74\%). Las ventas están muy
repartidas, lejos del clásico ``20\% de los productos genera el 80\% de las
ventas''.
\end{solucion}

\begin{solucion}{m5-9}
Primero el DDL con llaves y reglas; luego la carga con
\codigo{dbAppendTable()} (que respeta la estructura), la vista y el
índice. La fecha es la llave primaria de los hechos porque hay una venta
por día.

\begin{rcode}
con <- dbConnect(RSQLite::SQLite(), ":memory:")
dbExecute(con, "PRAGMA foreign_keys = ON")
dbExecute(con, "CREATE TABLE dim_producto (
  producto_id INTEGER PRIMARY KEY, nombre_producto TEXT, categoria TEXT)")
dbExecute(con, "CREATE TABLE dim_tienda (
  tienda_id INTEGER PRIMARY KEY, nombre_tienda TEXT, tipo TEXT)")
dbExecute(con, "CREATE TABLE fact_ventas_diarias (
  fecha TEXT PRIMARY KEY,
  producto_id INTEGER NOT NULL REFERENCES dim_producto(producto_id),
  tienda_id   INTEGER NOT NULL REFERENCES dim_tienda(tienda_id),
  cantidad    INTEGER CHECK (cantidad > 0),
  total REAL)")

dbAppendTable(con, "dim_producto",
  read_csv("datasets/productos.csv", show_col_types = FALSE) %>%
    transmute(producto_id = as.integer(producto_id), nombre_producto,
              categoria))
dbAppendTable(con, "dim_tienda",
  read_csv("datasets/tiendas.csv", show_col_types = FALSE) %>%
    transmute(tienda_id = as.integer(tienda_id), nombre_tienda, tipo))
dbAppendTable(con, "fact_ventas_diarias",
  read_csv("datasets/ventas_retail.csv", show_col_types = FALSE) %>%
    transmute(fecha = format(fecha, "%Y-%m-%d"),
              producto_id = as.integer(producto_id),
              tienda_id = as.integer(tienda_id),
              cantidad = as.integer(cantidad), total))

dbExecute(con, "CREATE VIEW vw_trimestre_tipo AS
  SELECT (CAST(strftime('%m', f.fecha) AS INTEGER) + 2) / 3 AS trimestre,
         t.tipo, COUNT(*) AS dias, ROUND(SUM(f.total), 0) AS ventas
  FROM fact_ventas_diarias AS f
  JOIN dim_tienda AS t ON f.tienda_id = t.tienda_id
  GROUP BY trimestre, t.tipo")
dbGetQuery(con, "SELECT * FROM vw_trimestre_tipo
                 WHERE trimestre = 1 ORDER BY ventas DESC")

consulta <- "SELECT SUM(total) FROM fact_ventas_diarias WHERE tienda_id = 3"
dbGetQuery(con, paste("EXPLAIN QUERY PLAN", consulta))$detail
dbExecute(con, "CREATE INDEX idx_tienda ON fact_ventas_diarias(tienda_id)")
dbGetQuery(con, paste("EXPLAIN QUERY PLAN", consulta))$detail
dbDisconnect(con)
\end{rcode}

\begin{rsalida}
[1] 0
[1] 0
[1] 0
[1] 0
[1] 50
[1] 10
[1] 365
[1] 0
  trimestre      tipo dias ventas
1         1  Estándar   46  54838
2         1 Principal   14  30699
3         1    Outlet   14  17226
4         1   Premium   10  16706
5         1   Express    6   6804
[1] "SCAN fact_ventas_diarias"
[1] 0
[1] "SEARCH fact_ventas_diarias USING INDEX idx_tienda (tienda_id=?)"
\end{rsalida}

Los \codigo{[1] 0} son las instrucciones DDL y los \codigo{[1] 50},
\codigo{10} y \codigo{365} las filas cargadas por \codigo{dbAppendTable()}.
En el primer trimestre, las tiendas de tipo Estándar suman más ventas
porque son seis. El plan pasa de \codigo{SCAN} (recorrer toda la tabla) a
\codigo{SEARCH \ldots{} USING INDEX idx\_tienda} tras crear el índice.
\end{solucion}

\begin{solucion}{m5-10}
La transformación \emph{no} filtra cantidades negativas a propósito: la
regla \codigo{CHECK} de la base es la última línea de defensa y queremos
ver cómo actúa. La carga usa una tabla temporal y un \codigo{INSERT} que
solo toma las fechas que aún no existen. ¿Por qué no \codigo{INSERT OR
IGNORE}? Porque en SQLite \codigo{OR IGNORE} ignora en silencio
\emph{cualquier} violación de restricción, no solo las llaves repetidas:
también la de \codigo{CHECK}. El lote con cantidad negativa se
``cargaría'' con estado OK y 0 filas nuevas, sin ningún error. Si lo
pruebas, lo verás.

\begin{rcode}
con <- dbConnect(RSQLite::SQLite(), ":memory:")
dbExecute(con, "CREATE TABLE fact_ventas_diarias (
  fecha TEXT PRIMARY KEY, producto_id INTEGER, cliente_id INTEGER,
  tienda_id INTEGER, cantidad INTEGER CHECK (cantidad > 0),
  precio_unitario REAL, descuento_pct REAL, total REAL)")
dbExecute(con, "CREATE TABLE log_etl (
  id INTEGER PRIMARY KEY AUTOINCREMENT, fecha_hora TEXT, lote TEXT,
  leidas INTEGER, nuevas INTEGER, estado TEXT, mensaje TEXT)")

extraer_ventas <- function(ruta) read_csv(ruta, show_col_types = FALSE)

transformar_ventas <- function(datos) {
  datos %>%
    distinct(fecha, .keep_all = TRUE) %>%
    transmute(fecha = format(as.Date(fecha), "%Y-%m-%d"),
              producto_id = as.integer(producto_id),
              cliente_id = as.integer(cliente_id),
              tienda_id = as.integer(tienda_id),
              cantidad = as.integer(cantidad), precio_unitario,
              descuento_pct, total)
}

cargar_ventas <- function(con, datos) {
  dbWriteTable(con, "stg", datos, temporary = TRUE, overwrite = TRUE)
  # Solo las fechas que aún no existen (sin OR IGNORE: ver explicación)
  n <- dbExecute(con, "INSERT INTO fact_ventas_diarias
                       SELECT * FROM stg
                       WHERE fecha NOT IN (SELECT fecha
                                           FROM fact_ventas_diarias)")
  dbExecute(con, "DROP TABLE stg")
  n
}

ejecutar_etl <- function(con, ruta) {
  leidas <- NA_integer_
  en_transaccion <- FALSE
  registrar <- function(nuevas, estado, mensaje = NA_character_) {
    dbAppendTable(con, "log_etl", tibble(
      fecha_hora = format(Sys.time()), lote = basename(ruta),
      leidas = leidas, nuevas = nuevas, estado = estado,
      mensaje = mensaje))
  }
  tryCatch({
    datos <- extraer_ventas(ruta)
    leidas <- nrow(datos)
    limpios <- transformar_ventas(datos)
    dbBegin(con); en_transaccion <- TRUE
    nuevas <- cargar_ventas(con, limpios)
    registrar(nuevas, "OK")
    dbCommit(con)
  }, error = function(e) {
    if (en_transaccion) dbRollback(con)
    registrar(0L, "ERROR", conditionMessage(e))
  })
  invisible(NULL)
}

# Simular los lotes en una carpeta temporal
carpeta <- file.path(tempdir(), "lotes")
dir.create(carpeta, showWarnings = FALSE)
ventas <- read_csv("datasets/ventas_retail.csv", show_col_types = FALSE)
lote <- function(desde, hasta, nombre) {
  ruta <- file.path(carpeta, nombre)
  ventas %>% filter(fecha >= as.Date(desde), fecha <= as.Date(hasta)) %>%
    write_csv(ruta)
  ruta
}
l1 <- lote("2023-01-01", "2023-06-30", "lote_ene_jun.csv")
l2 <- lote("2023-04-01", "2023-09-30", "lote_abr_sep.csv")
l3 <- lote("2023-10-01", "2023-12-31", "lote_oct_dic.csv")
l_malo <- file.path(carpeta, "lote_malo.csv")
tibble(fecha = as.Date("2024-01-01"), producto_id = 1, cliente_id = 1,
       tienda_id = 1, cantidad = -2, precio_unitario = 100,
       descuento_pct = 0, total = -200) %>% write_csv(l_malo)

for (ruta in c(l1, l2, l3, l3, l_malo)) ejecutar_etl(con, ruta)

dbGetQuery(con, "SELECT id, lote, leidas, nuevas, estado FROM log_etl")
dbGetQuery(con, "SELECT mensaje FROM log_etl WHERE estado = 'ERROR'")

# Validación final contra el CSV
dbGetQuery(con, "SELECT COUNT(*) AS filas, ROUND(SUM(total), 2) AS total
                 FROM fact_ventas_diarias")
round(sum(ventas$total), 2)
dbDisconnect(con)
\end{rcode}

\begin{rsalida}
[1] 0
[1] 0
  id             lote leidas nuevas estado
1  1 lote_ene_jun.csv    181    181     OK
2  2 lote_abr_sep.csv    183     92     OK
3  3 lote_oct_dic.csv     92     92     OK
4  4 lote_oct_dic.csv     92      0     OK
5  5    lote_malo.csv      1      0  ERROR
                                mensaje
1 CHECK constraint failed: cantidad > 0
  filas    total
1   365 513916.6
[1] 513916.6
\end{rsalida}

La bitácora cuenta la historia: el segundo lote leyó 183 días pero solo
cargó 92 (abril-junio ya estaban), el lote repetido no cargó nada y el
lote malo terminó en \codigo{ERROR} porque la regla \codigo{CHECK} rechazó
la cantidad negativa. La tabla final tiene exactamente 365 filas y el mismo
total que el CSV: el ETL es idempotente, atómico y trazable.
\end{solucion}


\section{Módulo 6. Dashboards interactivos con Shiny}
% ============================================================================
% SOLUCIONES DEL MÓDULO 6
% ============================================================================
% Nota: las soluciones que son apps muestran shinyApp() en un bloque que no se
% ejecuta; la lógica se comprueba con testServer(), que sí se ejecuta.

\begin{solucion}{m6-1}
Agregamos dos entradas a la \codigo{ui} y ampliamos el cálculo del
\codigo{renderText()}. El IVA se aplica después del descuento.

\begin{rcode}
library(shiny)

ui <- fluidPage(
  titlePanel("Calculadora de precio con descuento"),
  numericInput("precio", "Precio de lista ($):", value = 500,
               min = 0, step = 10),
  numericInput("cantidad", "Cantidad de piezas:", value = 1, min = 1),
  sliderInput("descuento", "Descuento (%):", min = 0, max = 50,
              value = 10, step = 5),
  checkboxInput("iva", "Incluir IVA (16%)", value = FALSE),
  h3(textOutput("total"))
)

server <- function(input, output, session) {
  output$total <- renderText({
    req(input$precio, input$cantidad)       # sin valores, en blanco
    subtotal <- input$precio * input$cantidad * (1 - input$descuento / 100)
    total <- if (input$iva) subtotal * 1.16 else subtotal
    paste0("Total a pagar: ", scales::dollar(total))
  })
}

testServer(server, {
  session$setInputs(precio = 500, cantidad = 3, descuento = 10,
                    iva = FALSE)
  print(output$total)
  session$setInputs(iva = TRUE)
  print(output$total)
})
\end{rcode}
\begin{rsalida}
[1] "Total a pagar: $1,350"
[1] "Total a pagar: $1,566"
\end{rsalida}

Sin IVA, 3 piezas de \$500 con 10\% de descuento cuestan \$1,350; con IVA,
\$1,566, como pedía el enunciado. Para abrir la app:

\begin{rnoejecutar}
shinyApp(ui, server)
\end{rnoejecutar}
\end{solucion}

\begin{solucion}{m6-2}
Hay tres cambios en \archivo{app.R} de la app 02: (1) un argumento nuevo en
\codigo{filtrar\_ventas()}, (2) la entrada en la barra lateral y (3) pasar
la entrada al filtrar. Primero probamos la función modificada, de forma
autocontenida:

\begin{rcode}
library(tidyverse)
library(scales)
ventas <- read_csv("datasets/transacciones.csv", show_col_types = FALSE)

# (1) filtrar_ventas() con el argumento nuevo 'pagos' (resto igual)
filtrar_pagos <- function(datos, pagos = NULL) {
  if (length(pagos) > 0) {                 # vacío = todos los métodos
    datos <- filter(datos, metodo_pago %in% pagos)
  }
  datos
}
solo_efectivo <- filtrar_pagos(ventas, pagos = "Efectivo")
nrow(solo_efectivo)
dollar(sum(solo_efectivo$total_transaccion), accuracy = 1)
\end{rcode}
\begin{rsalida}
[1] 315
[1] "$303,735"
\end{rsalida}

Con solo ``Efectivo'' quedan 315 transacciones por \$303,735. Los cambios en
la app (la app completa modificada se probó arrancándola y con
\codigo{testServer()}):

\begin{rnoejecutar}
# (1) En filtrar_ventas(): nuevo argumento y nuevo if
filtrar_ventas <- function(datos, tienda = "Todas", categorias = NULL,
                           fechas = range(datos$fecha), pagos = NULL) {
  # ... filtros de fecha, tienda y categoría como antes ...
  if (length(pagos) > 0) {
    resultado <- filter(resultado, metodo_pago %in% pagos)
  }
  resultado
}

# (2) En sidebarPanel(), después del selectizeInput de categorías:
checkboxGroupInput("pagos", "Método de pago:",
                   choices = sort(unique(ventas$metodo_pago)),
                   selected = sort(unique(ventas$metodo_pago))),

# (3) En el server:
datos_filtrados <- reactive({
  req(input$fechas)
  filtrar_ventas(ventas, input$tienda, input$categorias, input$fechas,
                 input$pagos)
})
\end{rnoejecutar}

Si el usuario desmarca todas las casillas, \codigo{input\$pagos} vale
\codigo{NULL} y, por nuestra convención, se muestran todos los métodos.
\end{solucion}

\begin{solucion}{m6-3}
Primero la lógica. El margen bruto de cada transacción es lo cobrado menos
el costo de las piezas vendidas (\codigo{costo * cantidad}):

\begin{rcode}
library(tidyverse)
library(scales)
productos <- read_csv("datasets/productos.csv", show_col_types = FALSE)
datos <- read_csv("datasets/transacciones.csv", show_col_types = FALSE) %>%
  left_join(select(productos, producto_id, categoria, costo),
            by = "producto_id")

calcular_margen <- function(datos) {
  datos %>%
    summarise(ventas = sum(total_transaccion),
              margen = sum(total_transaccion - costo * cantidad),
              pct_margen = margen / ventas)
}

calcular_margen(datos)
datos %>%
  group_by(categoria) %>%
  calcular_margen() %>%
  arrange(pct_margen) %>%
  head(3)
\end{rcode}
\begin{rsalida}
# A tibble: 1 x 3
   ventas  margen pct_margen
    <dbl>   <dbl>      <dbl>
1 959993. 444637.      0.463
# A tibble: 3 x 4
  categoria  ventas margen pct_margen
  <chr>       <dbl>  <dbl>      <dbl>
1 Ropa       30048.  5835.      0.194
2 Hogar      63407. 13079.      0.206
3 Belleza   182143. 56054.      0.308
\end{rsalida}

En total la empresa obtiene un margen bruto de 46.3\% (\$444,637). Las
categorías con menor margen son Ropa (19.4\%), Hogar (20.6\%) y Belleza
(30.8\%): Belleza es la tercera más baja porque varios de sus productos tienen
un costo superior a su precio de catálogo (el hallazgo de calidad de datos
que ya conoces), aunque sus precios de venta reales lo compensan en parte. Justo para eso sirve el color condicional. En la
app, agregamos \codigo{calcular\_margen()} a
\archivo{R/funciones\_dashboard.R} y estas piezas a \archivo{app.R}:

\begin{rnoejecutar}
# ui: nueva fila dentro de tabItem(tabName = "resumen")
fluidRow(
  valueBoxOutput("kpi_margen", width = 6),
  valueBoxOutput("kpi_pct_margen", width = 6)
)

# server
margen <- reactive(calcular_margen(datos_filtrados()))

output$kpi_margen <- renderValueBox({
  m <- margen()$margen
  valueBox(formato_pesos(m), "Margen bruto",
           icon = icon("sack-dollar"),
           color = if (isTRUE(m >= 0)) "green" else "red")
})
output$kpi_pct_margen <- renderValueBox({
  valueBox(formato_pct(margen()$pct_margen), "% de margen",
           icon = icon("percent"), color = "blue")
})
\end{rnoejecutar}

\codigo{isTRUE()} evita un error cuando el filtro no deja ventas (el margen
sería 0 y el porcentaje \codigo{NaN}). Con la categoría Belleza y la tienda
correspondiente puedes llegar a ver la caja en rojo.
\end{solucion}

\begin{solucion}{m6-4}
El truco es que \emph{todas} las salidas dependan de un
\codigo{eventReactive()} ligado al botón, y no directamente de los filtros.

\begin{rcode}
library(shiny)
library(tidyverse)
library(scales)
tiendas <- read_csv("datasets/tiendas.csv", show_col_types = FALSE)
ventas <- read_csv("datasets/transacciones.csv", show_col_types = FALSE) %>%
  left_join(select(tiendas, tienda_id, nombre_tienda), by = "tienda_id")

ui <- fluidPage(
  titlePanel("Ventas con botón Aplicar"),
  sidebarLayout(
    sidebarPanel(
      selectizeInput("tiendas", "Tiendas:",
                     choices = sort(unique(ventas$nombre_tienda)),
                     multiple = TRUE, options = list(placeholder = "Todas")),
      dateRangeInput("fechas", "Periodo:", start = "2023-01-01",
                     end = "2023-12-31", language = "es",
                     separator = " a "),
      actionButton("aplicar", "Aplicar", icon = icon("filter"))
    ),
    mainPanel(h3(textOutput("ventas")), h4(textOutput("transacciones")))
  )
)

server <- function(input, output, session) {
  # ignoreNULL = FALSE: calcula también al abrir la app (con el valor 0
  # del botón), para no empezar con la pantalla vacía
  datos <- eventReactive(input$aplicar, {
    d <- filter(ventas, fecha >= input$fechas[1], fecha <= input$fechas[2])
    if (length(input$tiendas) > 0) {
      d <- filter(d, nombre_tienda %in% input$tiendas)
    }
    d
  }, ignoreNULL = FALSE)

  output$ventas <- renderText(
    paste("Ventas:", dollar(sum(datos()$total_transaccion), accuracy = 1)))
  output$transacciones <- renderText(
    paste("Transacciones:", nrow(datos())))
}

testServer(server, {
  session$setInputs(tiendas = NULL, aplicar = 0,
                    fechas = as.Date(c("2023-01-01", "2023-12-31")))
  print(output$ventas)
  session$setInputs(tiendas = "Sucursal Norte")      # sin pulsar
  print(output$ventas)
  session$setInputs(aplicar = 1)                     # pulsa Aplicar
  print(output$ventas)
  print(output$transacciones)
})
\end{rcode}
\begin{rsalida}
[1] "Ventas: $959,993"
[1] "Ventas: $959,993"
[1] "Ventas: $89,024"
[1] "Transacciones: 99"
\end{rsalida}

La prueba lo demuestra: al elegir ``Sucursal Norte'' sin pulsar, las ventas
siguen en \$959,993; al pulsar ``Aplicar'' cambian a las de la Sucursal
Norte (\$89,024 en 99 transacciones).

\begin{rnoejecutar}
shinyApp(ui, server)
\end{rnoejecutar}
\end{solucion}

\begin{solucion}{m6-5}
Primero la función de la gráfica. Con \codigo{wday(fecha, label = TRUE,
week\_start = 1)} obtenemos el día como factor ordenado de lunes a domingo
(en español, porque trabajamos con locale es\_MX):

\begin{rcode}
library(tidyverse)
library(lubridate)
library(scales)
datos <- read_csv("datasets/transacciones.csv", show_col_types = FALSE)

ventas_por_dia <- function(datos) {
  datos %>%
    mutate(dia = wday(fecha, label = TRUE, abbr = FALSE,
                      week_start = 1)) %>%
    group_by(dia) %>%
    summarise(ventas = sum(total_transaccion), .groups = "drop") %>%
    mutate(es_mejor = ventas == max(ventas))
}

grafica_dias <- function(datos) {
  ggplot(ventas_por_dia(datos), aes(x = dia, y = ventas, fill = es_mejor)) +
    geom_col(width = 0.7) +
    scale_fill_manual(values = c(`TRUE` = "#eb6834", `FALSE` = "#b5b3ad"),
                      guide = "none") +
    scale_y_continuous(labels = label_dollar(),
                       expand = expansion(mult = c(0, 0.05))) +
    labs(x = NULL, y = "Ventas") +
    theme_minimal(base_size = 12) +
    theme(panel.grid.minor = element_blank())
}

ventas_por_dia(datos)
p_dias <- grafica_dias(datos)
\end{rcode}
\begin{rsalida}
# A tibble: 7 x 3
  dia        ventas es_mejor
  <ord>       <dbl> <lgl>   
1 lunes     128878. FALSE   
2 martes    119364. FALSE   
3 miércoles 141329. FALSE   
4 jueves    145280. FALSE   
5 viernes   161266. TRUE    
6 sábado    115172. FALSE   
7 domingo   148703. FALSE   
\end{rsalida}

En todo 2023 el mejor día es el viernes (\$161,266) y el más flojo el
sábado (\$115,172). Como la función recibe los datos, en la app basta con pasarle
\codigo{datos\_validos()} para que respete los filtros globales. Cambios en
la app 03:

\begin{rnoejecutar}
# En sidebarMenu(), después de "Clientes":
menuItem("Días", tabName = "dias", icon = icon("calendar-week")),

# En tabItems(), una pestaña nueva:
tabItem(tabName = "dias",
  fluidRow(
    box(title = "Ventas por día de la semana", width = 12,
        status = "primary", plotOutput("graf_dias", height = 380))
  )
),

# En el server:
output$graf_dias <- renderPlot({
  grafica_dias(datos_validos())
}, res = 96)
\end{rnoejecutar}

La función \codigo{grafica\_dias()} (con \codigo{ventas\_por\_dia()}) va en
\archivo{R/funciones\_dashboard.R}, junto a las demás gráficas.
\end{solucion}

\begin{solucion}{m6-6}
Para que la descarga contenga ``exactamente la tabla que se ve'', guardamos
la tabla en un reactivo que usan \emph{tanto} la tabla como los dos botones.
Así no hay forma de que diverjan. Probemos primero, autocontenido, lo que
escribirá cada botón:

\begin{rcode}
library(tidyverse)
library(scales)
library(writexl)
productos <- read_csv("datasets/productos.csv", show_col_types = FALSE)
ventas <- read_csv("datasets/transacciones.csv", show_col_types = FALSE) %>%
  left_join(select(productos, producto_id, nombre_producto, categoria),
            by = "producto_id")

top_productos <- function(datos, n = 10) {
  datos %>%
    group_by(Producto = nombre_producto, `Categoría` = categoria) %>%
    summarise(Unidades = sum(cantidad),
              Ventas   = sum(total_transaccion), .groups = "drop") %>%
    slice_max(Ventas, n = n, with_ties = FALSE) %>%
    mutate(Ventas = dollar(Ventas))
}

dir.create("resultados", showWarnings = FALSE)
nombre <- paste0("top_productos_", Sys.Date())
top5 <- top_productos(ventas, n = 5)
write_excel_csv(top5, file.path("resultados", paste0(nombre, ".csv")))
write_xlsx(top5, file.path("resultados", paste0(nombre, ".xlsx")))
read_csv(file.path("resultados", paste0(nombre, ".csv")),
         show_col_types = FALSE)
\end{rcode}
\begin{rsalida}
# A tibble: 5 x 4
  Producto    Categoría  Unidades Ventas    
  <chr>       <chr>         <dbl> <chr>     
1 Producto 06 Juguetes         82 $34,028.04
2 Producto 19 Deportes         85 $28,367.95
3 Producto 31 Alimentos        74 $25,429.70
4 Producto 47 Automotriz       82 $24,527.21
5 Producto 46 Automotriz       74 $24,470.29
\end{rsalida}

Los cambios en la app 02:

\begin{rnoejecutar}
# ui: la pestaña "Top productos" ahora tiene los botones arriba
tabPanel("Top productos",
         br(),
         downloadButton("bajar_csv", "Descargar CSV"),
         downloadButton("bajar_excel", "Descargar Excel"),
         br(), br(),
         tableOutput("tabla_top"))

# server: un reactivo con la tabla visible
tabla_visible <- reactive(top_productos(datos_filtrados(), n = input$n_top))

output$tabla_top <- renderTable(tabla_visible(), align = "llrr",
                                digits = 0)

output$bajar_csv <- downloadHandler(
  filename = function() paste0("top_productos_", Sys.Date(), ".csv"),
  content  = function(file) write_excel_csv(tabla_visible(), file)
)
output$bajar_excel <- downloadHandler(
  filename = function() paste0("top_productos_", Sys.Date(), ".xlsx"),
  content  = function(file) writexl::write_xlsx(tabla_visible(), file)
)
\end{rnoejecutar}

Para probarlo con \codigo{testServer()} sobre la app modificada, basta con
leer la salida del botón, que devuelve la ruta del archivo generado:

\begin{rnoejecutar}
testServer(shinyAppDir("mis-apps/02-mi-explorador"), {
  session$setInputs(tienda = "Todas", categorias = NULL, n_top = 5,
                    fechas = as.Date(c("2023-01-01", "2023-12-31")),
                    metrica = "ventas")
  read_csv(output$bajar_csv, show_col_types = FALSE)   # 5 filas
})
\end{rnoejecutar}
\end{solucion}

\begin{solucion}{m6-7}
El módulo recibe los datos como \emph{reactivo} (sin paréntesis al
pasarlo) y el nombre de la columna como texto. Dentro, cada id pasa por
\codigo{ns()}, así que las dos copias no chocan aunque ambas tengan un
deslizador llamado \codigo{"n"}.

\begin{rcode}
library(shiny)
library(shinydashboard)
library(tidyverse)
library(scales)
tiendas <- read_csv("datasets/tiendas.csv", show_col_types = FALSE)
productos <- read_csv("datasets/productos.csv", show_col_types = FALSE)
ventas <- read_csv("datasets/transacciones.csv", show_col_types = FALSE) %>%
  left_join(select(tiendas, tienda_id, ciudad), by = "tienda_id") %>%
  left_join(select(productos, producto_id, nombre_producto),
            by = "producto_id")

# ---- Módulo ------------------------------------------------------------
ranking_ui <- function(id, titulo) {
  ns <- NS(id)
  box(title = titulo, width = 6, status = "primary",
      sliderInput(ns("n"), "Mostrar top:", min = 3, max = 10, value = 5),
      plotOutput(ns("grafica"), height = 300))
}

ranking_server <- function(id, datos, columna) {
  moduleServer(id, function(input, output, session) {
    resumen <- reactive({
      datos() %>%
        group_by(nombre = .data[[columna]]) %>%
        summarise(ventas = sum(total_transaccion), .groups = "drop") %>%
        slice_max(ventas, n = input$n, with_ties = FALSE)
    })
    output$grafica <- renderPlot({
      ggplot(resumen(), aes(x = ventas, y = reorder(nombre, ventas))) +
        geom_col(fill = "#2a78d6", width = 0.7) +
        scale_x_continuous(labels = label_dollar()) +
        labs(x = "Ventas", y = NULL) +
        theme_minimal(base_size = 12)
    }, res = 96)
    resumen                    # el módulo devuelve su tabla (reactivo)
  })
}

# ---- App que usa el módulo dos veces -----------------------------------
ui <- dashboardPage(
  dashboardHeader(title = "Rankings"),
  dashboardSidebar(disable = TRUE),
  dashboardBody(fluidRow(
    ranking_ui("productos", "Top productos"),
    ranking_ui("ciudades", "Top ciudades")
  ))
)
server <- function(input, output, session) {
  datos <- reactive(ventas)
  ranking_server("productos", datos, "nombre_producto")
  ranking_server("ciudades", datos, "ciudad")
}

# Cada copia del módulo se prueba por separado con su propio 'n'
testServer(ranking_server,
           args = list(datos = reactive(ventas), columna = "ciudad"), {
  session$setInputs(n = 3)
  print(resumen())
})
testServer(ranking_server,
           args = list(datos = reactive(ventas),
                       columna = "nombre_producto"), {
  session$setInputs(n = 4)
  print(nrow(resumen()))
})
\end{rcode}
\begin{rsalida}
# A tibble: 3 x 2
  nombre       ventas
  <chr>         <dbl>
1 CDMX        383588.
2 Guadalajara 169260.
3 Tijuana     114377.
[1] 4
\end{rsalida}

La copia de ciudades con \codigo{n = 3} devuelve las tres ciudades con más
ventas (la CDMX, con cuatro tiendas, domina con \$383,588) y la de productos
con \codigo{n = 4} devuelve 4 filas: cada módulo responde a su propio
deslizador. En la \codigo{ui} completa los ids reales son
\codigo{productos-n} y \codigo{ciudades-n}.

\begin{rnoejecutar}
shinyApp(ui, server)
\end{rnoejecutar}
\end{solucion}

\begin{solucion}{m6-8}
Seguimos la receta del módulo: datos fuera del \codigo{server}, funciones
de lógica probadas por separado y un servidor corto.

\begin{rcode}
library(shiny)
library(shinydashboard)
library(tidyverse)
library(scales)
library(DT)
library(writexl)

empleados <- read_csv("datasets/empleados.csv", show_col_types = FALSE)

# ---- Lógica --------------------------------------------------------------
filtrar_empleados <- function(datos, deptos = NULL, sucursales = NULL) {
  if (length(deptos) > 0) datos <- filter(datos, departamento %in% deptos)
  if (length(sucursales) > 0) {
    datos <- filter(datos, sucursal %in% sucursales)
  }
  datos
}

kpis_rrhh <- function(datos) {
  datos %>%
    summarise(plantilla = n(),
              salario_prom = mean(salario_mensual),
              antiguedad_prom = mean(antiguedad_anos),
              pct_posgrado = mean(nivel_educacion %in%
                                    c("Maestría", "Doctorado")))
}

grafica_salario_puesto <- function(datos) {
  datos %>%
    group_by(puesto) %>%
    summarise(salario = mean(salario_mensual), .groups = "drop") %>%
    ggplot(aes(x = salario, y = reorder(puesto, salario))) +
    geom_col(fill = "#2a78d6", width = 0.7) +
    scale_x_continuous(labels = label_dollar(),
                       expand = expansion(mult = c(0, 0.05))) +
    labs(x = "Salario mensual promedio", y = NULL) +
    theme_minimal(base_size = 12) +
    theme(panel.grid.minor = element_blank())
}

tabla_empleados <- function(datos) {
  datos %>%
    select(Nombre = nombre, Departamento = departamento, Puesto = puesto,
           Sucursal = sucursal, Salario = salario_mensual,
           `Antigüedad` = antiguedad_anos) %>%
    arrange(desc(Salario))
}

kpis_rrhh(empleados)
kpis_rrhh(filtrar_empleados(empleados, deptos = "Ventas"))
\end{rcode}
\begin{rsalida}
# A tibble: 1 x 4
  plantilla salario_prom antiguedad_prom pct_posgrado
      <int>        <dbl>           <dbl>        <dbl>
1       100       27481.            3.96         0.35
# A tibble: 1 x 4
  plantilla salario_prom antiguedad_prom pct_posgrado
      <int>        <dbl>           <dbl>        <dbl>
1        16       29032.            4.14        0.312
\end{rsalida}

La empresa tiene 100 empleados con salario mensual promedio de \$27,481 y
casi 4 años de antigüedad promedio; el 35\% tiene maestría o doctorado. El
departamento de Ventas tiene 16 personas, con salario promedio de \$29,032. Ahora la interfaz y el servidor:

\begin{rcode}
ui <- dashboardPage(
  dashboardHeader(title = "Recursos Humanos"),
  dashboardSidebar(
    selectizeInput("deptos", "Departamentos:",
                   choices = sort(unique(empleados$departamento)),
                   multiple = TRUE, options = list(placeholder = "Todos")),
    selectizeInput("sucursales", "Sucursales:",
                   choices = sort(unique(empleados$sucursal)),
                   multiple = TRUE, options = list(placeholder = "Todas"))
  ),
  dashboardBody(
    fluidRow(
      valueBoxOutput("plantilla", width = 3),
      valueBoxOutput("salario", width = 3),
      valueBoxOutput("antiguedad", width = 3),
      valueBoxOutput("posgrado", width = 3)
    ),
    fluidRow(
      box(title = "Salario promedio por puesto", width = 5,
          status = "primary", plotOutput("graf_puesto", height = 320)),
      box(title = "Empleados", width = 7, status = "primary",
          downloadButton("bajar", "Descargar Excel"), br(), br(),
          DTOutput("tabla"))
    )
  )
)

server <- function(input, output, session) {
  datos <- reactive(filtrar_empleados(empleados, input$deptos,
                                      input$sucursales))
  kpis <- reactive(kpis_rrhh(datos()))

  output$plantilla <- renderValueBox(
    valueBox(kpis()$plantilla, "Plantilla", icon = icon("users")))
  output$salario <- renderValueBox(
    valueBox(dollar(kpis()$salario_prom, accuracy = 1),
             "Salario mensual promedio", icon = icon("money-bill")))
  output$antiguedad <- renderValueBox(
    valueBox(paste(round(kpis()$antiguedad_prom, 1), "años"),
             "Antigüedad promedio", icon = icon("clock")))
  output$posgrado <- renderValueBox(
    valueBox(percent(kpis()$pct_posgrado, accuracy = 1),
             "Con maestría o doctorado", icon = icon("graduation-cap")))

  output$graf_puesto <- renderPlot({
    validate(need(nrow(datos()) > 0, "No hay empleados con esos filtros."))
    grafica_salario_puesto(datos())
  }, res = 96)

  output$tabla <- renderDT(
    datatable(tabla_empleados(datos()), rownames = FALSE,
              options = list(pageLength = 8)) %>%
      formatCurrency("Salario", currency = "$", digits = 0))

  output$bajar <- downloadHandler(
    filename = function() paste0("empleados_", Sys.Date(), ".xlsx"),
    content  = function(file) write_xlsx(tabla_empleados(datos()), file)
  )
}

testServer(server, {
  session$setInputs(deptos = c("IT", "Finanzas"), sucursales = NULL)
  cat("Plantilla IT + Finanzas:", kpis()$plantilla, "\n")
  cat("Filas en el Excel:",
      nrow(readxl::read_excel(output$bajar)), "\n")
})
\end{rcode}
\begin{rsalida}
Plantilla IT + Finanzas: 18 
Filas en el Excel: 18 
\end{rsalida}

La prueba confirma que los filtros múltiples funcionan y que el Excel
descargado contiene exactamente a los empleados filtrados. Para abrir el
dashboard:

\begin{rnoejecutar}
shinyApp(ui, server)
\end{rnoejecutar}
\end{solucion}


\section{Módulo 7. Machine learning para BI}
% ============================================================================
% SOLUCIONES DEL MÓDULO 7: MACHINE LEARNING PARA BUSINESS INTELLIGENCE
% ============================================================================

\begin{solucion}{m7-1}
\textbf{Parte conceptual.}
\begin{enumerate}[label=(\alph*)]
  \item Estimar el ticket promedio de un cliente nuevo: \textbf{regresión}
        (la respuesta es un número). Regresión lineal o random forest.
  \item Detectar si un correo es una queja: \textbf{clasificación} (queja /
        no queja). Regresión logística o random forest sobre variables del
        texto (palabras clave).
  \item Agrupar sucursales por su mezcla de ventas: \textbf{clustering}.
        k-means con las variables escaladas.
  \item Demanda semanal de pan de muerto: \textbf{serie de tiempo} con
        estacionalidad anual muy marcada. Ingenuo estacional como línea base,
        ETS o ARIMA.
  \item Qué recomendar a quien compra una cafetera: \textbf{reglas de
        asociación} (análisis de canasta) ordenadas por lift.
\end{enumerate}

\textbf{Parte práctica.} Ajustamos el modelo y extraemos el R² ajustado:

\begin{rcode}
library(tidyverse)
ventas <- read_csv("datasets/ventas_retail.csv", show_col_types = FALSE)

modelo <- lm(total ~ descuento_pct + dia_semana, data = ventas)
round(summary(modelo)$r.squared, 4)
round(summary(modelo)$adj.r.squared, 4)
\end{rcode}

\begin{rsalida}
[1] 0.004
[1] -0.0155
\end{rsalida}

El R² es de apenas 0.004 y el R² ajustado es negativo ($-0.0155$): el
descuento y el día de la semana no explican el total de la venta. No hay
señal, porque en el generador de datos el descuento y el día se eligieron al
azar, independientemente del monto.
\end{solucion}

\begin{solucion}{m7-2}
Escribimos cada fórmula directamente con vectores:

\begin{rcode}
real       <- c(120, 95, 150, 80, 110, 130)
pronostico <- c(110, 100, 140, 90, 115, 120)
error <- real - pronostico

mae  <- mean(abs(error))
rmse <- sqrt(mean(error^2))
mape <- mean(abs(error) / real) * 100
r2   <- 1 - sum(error^2) / sum((real - mean(real))^2)

round(c(MAE = mae, RMSE = rmse, MAPE = mape, R2 = r2), 3)
\end{rcode}

\begin{rsalida}
  MAE  RMSE  MAPE    R2 
8.333 8.660 7.500 0.856 
\end{rsalida}

En promedio, el gerente se equivoca por 8.3 mil pesos por sucursal (MAE),
que equivale a 7.5\% del valor real (MAPE). El RMSE (8.66) es apenas mayor que
el MAE porque ningún error es mucho más grande que los demás (todos son de 5 o
10 mil pesos). El R² de 0.856 indica que el pronóstico captura el 85.6\% de las
diferencias entre sucursales.
\end{solucion}

\begin{solucion}{m7-3}
Los datos: 50 fraudes reales, de los que se detectaron 40 (VP = 40, FN = 10);
950 legítimas, de las que 60 se marcaron como fraude (FP = 60, VN = 890).

\begin{rcode}
matriz <- matrix(c(890, 60, 10, 40), nrow = 2,
                 dimnames = list(Prediccion = c("Legítima", "Fraude"),
                                 Real = c("Legítima", "Fraude")))
matriz

vp <- 40; fn <- 10; fp <- 60; vn <- 890
precision    <- vp / (vp + fp)
sensibilidad <- vp / (vp + fn)
round(c(exactitud     = (vp + vn) / (vp + vn + fp + fn),
        precision     = precision,
        sensibilidad  = sensibilidad,
        especificidad = vn / (vn + fp),
        F1 = 2 * precision * sensibilidad / (precision + sensibilidad)), 3)

costo_errores <- fn * 5000 + fp * 100
costo_errores
\end{rcode}

\begin{rsalida}
          Real
Prediccion Legítima Fraude
  Legítima      890     10
  Fraude         60     40
    exactitud     precision  sensibilidad especificidad            F1 
        0.930         0.400         0.800         0.937         0.533 
[1] 56000
\end{rsalida}

La exactitud (93\%) parece alta, pero la precisión es de solo 40\%: de cada 10
alertas, 6 son transacciones legítimas. La sensibilidad de 80\% indica que se
detecta a 4 de cada 5 fraudes. El costo de los errores es de \$56{,}000:
\$50{,}000 por los 10 fraudes no detectados y \$6{,}000 por las 60 revisiones
innecesarias. Con estos costos, conviene bajar el umbral aunque aumenten las
falsas alarmas: cada fraude detectado adicional ahorra lo mismo que 50
revisiones.
\end{solucion}

\begin{solucion}{m7-4}
Calculamos el RFM, las calificaciones con \codigo{ntile()} (recordando que en
recencia menos es mejor) y los niveles con \codigo{case\_when()}:

\begin{rcode}
library(tidyverse)
transacciones <- read_csv("datasets/transacciones.csv",
                          show_col_types = FALSE)
fecha_corte <- as.Date("2024-01-01")

niveles <- transacciones %>%
  group_by(cliente_id) %>%
  summarise(recencia   = as.numeric(fecha_corte - max(fecha)),
            frecuencia = n(),
            monto      = sum(total_transaccion), .groups = "drop") %>%
  mutate(r = ntile(desc(recencia), 5),
         f = ntile(frecuencia, 5),
         m = ntile(monto, 5),
         puntaje = r + f + m,
         nivel = case_when(puntaje >= 13 ~ "Oro",
                           puntaje >= 9  ~ "Plata",
                           TRUE          ~ "Bronce"),
         nivel = factor(nivel, levels = c("Oro", "Plata", "Bronce")))

niveles %>%
  group_by(nivel) %>%
  summarise(clientes = n(),
            monto_prom = round(mean(monto)),
            ingresos = sum(monto), .groups = "drop") %>%
  mutate(pct_ingresos = round(100 * ingresos / sum(ingresos), 1)) %>%
  select(-ingresos)
\end{rcode}

\begin{rsalida}
# A tibble: 3 x 4
  nivel  clientes monto_prom pct_ingresos
  <fct>     <int>      <dbl>        <dbl>
1 Oro          39       8439         34.3
2 Plata        70       5444         39.7
3 Bronce       89       2807         26  
\end{rsalida}

Los 39 clientes Oro (20\% de los 198 compradores) generan el 34.3\% de los
ingresos, con un monto promedio tres veces mayor que el de los Bronce (\$8{,}439
contra \$2{,}807). Los Plata son el grupo con mayor peso en ingresos
(39.7\%): son el candidato natural para programas que los conviertan en Oro.
\end{solucion}

\begin{solucion}{m7-5}
Construimos las variables por cliente, escalamos y calculamos la silueta
promedio para $k$ de 2 a 6:

\begin{rcode}
library(tidyverse)
library(cluster)
transacciones <- read_csv("datasets/transacciones.csv",
                          show_col_types = FALSE)

perfil_cliente <- transacciones %>%
  group_by(cliente_id) %>%
  summarise(monto_total     = sum(total_transaccion),
            ticket_promedio = mean(total_transaccion),
            n_tiendas       = n_distinct(tienda_id),
            n_productos     = n_distinct(producto_id), .groups = "drop")

datos_esc <- scale(perfil_cliente %>% select(-cliente_id))

siluetas <- map_dbl(2:6, function(k) {
  set.seed(123)
  km <- kmeans(datos_esc, centers = k, nstart = 25)
  mean(silhouette(km$cluster, dist(datos_esc))[, "sil_width"])
})
tibble(k = 2:6, silueta = round(siluetas, 3))
k_elegido <- (2:6)[which.max(siluetas)]
k_elegido
\end{rcode}

\begin{rsalida}
# A tibble: 5 x 2
      k silueta
  <int>   <dbl>
1     2   0.322
2     3   0.325
3     4   0.356
4     5   0.326
5     6   0.331
[1] 4
\end{rsalida}

La silueta más alta es para $k = 4$ (0.356), aunque todas están entre 0.32 y
0.36: la estructura de grupos es débil (datos aleatorios). El código guarda el
$k$ elegido en \codigo{k\_elegido} para no escribirlo a mano.

\begin{rcode}
set.seed(123)
km <- kmeans(datos_esc, centers = k_elegido, nstart = 25)

perfil <- perfil_cliente %>%
  mutate(cluster = km$cluster) %>%
  group_by(cluster) %>%
  summarise(clientes = n(),
            monto_total = round(mean(monto_total)),
            ticket = round(mean(ticket_promedio)),
            tiendas = round(mean(n_tiendas), 1),
            productos = round(mean(n_productos), 1), .groups = "drop")

# Nombres según el perfil (reglas de negocio, nunca por número de grupo)
extremos <- range(perfil$monto_total)
perfil %>%
  mutate(segmento = case_when(
    monto_total == extremos[2] ~ "Alto valor",
    monto_total == extremos[1] ~ "Valor mínimo",
    ticket == max(ticket[!monto_total %in% extremos]) ~
      "Ticket alto ocasional",
    TRUE ~ "Valor medio")) %>%
  arrange(desc(monto_total)) %>%
  select(segmento, everything(), -cluster)
\end{rcode}

\begin{rsalida}
# A tibble: 4 x 6
  segmento              clientes monto_total ticket tiendas productos
  <chr>                    <int>       <dbl>  <dbl>   <dbl>     <dbl>
1 Alto valor                  22       10159   1123     6.1       8.5
2 Valor medio                 96        4999    873     4.7       5.6
3 Ticket alto ocasional       39        4756   1547     2.7       3  
4 Valor mínimo                41        1734    597     2.6       2.8
\end{rsalida}

Los nombres salen de reglas aplicadas al perfil: el mayor monto es ``Alto
valor'' (22 clientes que compran en 6 tiendas distintas en promedio), el menor
es ``Valor mínimo'', y entre los dos restantes, el de mayor ticket es ``Ticket
alto ocasional'' (39 clientes que compran pocas veces, en pocas tiendas, pero
con el ticket más alto: \$1{,}547). Este último grupo no se ve en un RFM
clásico y sugiere una estrategia propia: aumentar su frecuencia con
recordatorios, porque cada visita vale mucho.
\end{solucion}

\begin{solucion}{m7-6}
Reproducimos la simulación y la división del texto, y comparamos tres
modelos:

\begin{rcode}
library(tidyverse)
library(caret)

set.seed(123)
n_clientes <- 2000
clientes_churn <- tibble(
  cliente_id       = sprintf("C%04d", 1:n_clientes),
  antiguedad_meses = sample(1:72, n_clientes, replace = TRUE),
  plan             = sample(c("Básico", "Estándar", "Premium"),
                            n_clientes, replace = TRUE,
                            prob = c(0.5, 0.3, 0.2)),
  uso_mensual      = rpois(n_clientes, 8),
  quejas_6m        = rpois(n_clientes, 0.8),
  pagos_tardios    = rpois(n_clientes, 0.4),
  genero           = sample(c("F", "M"), n_clientes, replace = TRUE),
  edad             = sample(18:70, n_clientes, replace = TRUE)
) %>%
  mutate(
    efecto_plan = case_when(plan == "Básico"   ~ 0,
                            plan == "Estándar" ~ -0.6,
                            TRUE               ~ -1.3),
    riesgo = 2.0 - 0.05 * antiguedad_meses +
      1.0 * (antiguedad_meses <= 6) + 1.0 * quejas_6m -
      0.35 * uso_mensual + 0.8 * pagos_tardios + efecto_plan,
    churn = if_else(runif(n_clientes) < plogis(riesgo), "Si", "No"),
    churn = factor(churn, levels = c("No", "Si")),
    plan  = factor(plan, levels = c("Básico", "Estándar", "Premium")),
    cliente_nuevo = as.integer(antiguedad_meses <= 6)   # variable nueva
  ) %>%
  select(-efecto_plan, -riesgo)

set.seed(123)
filas_ent <- createDataPartition(clientes_churn$churn, p = 0.7,
                                 list = FALSE)
ent    <- clientes_churn[filas_ent, ]
prueba <- clientes_churn[-filas_ent, ]

calcular_auc <- function(prob, real, positivo = "Si") {
  p_pos <- prob[real == positivo]
  p_neg <- prob[real != positivo]
  mean(outer(p_pos, p_neg, ">") + 0.5 * outer(p_pos, p_neg, "=="))
}

f_original <- churn ~ antiguedad_meses + plan + uso_mensual + quejas_6m +
  pagos_tardios + genero + edad
f_nuevo    <- update(f_original, . ~ . + cliente_nuevo)
f_justo    <- update(f_nuevo, . ~ . - genero - edad)

modelos <- list(Original = f_original, `Con cliente_nuevo` = f_nuevo,
                `Sin genero ni edad` = f_justo) %>%
  map(~ glm(.x, data = ent, family = binomial))

map_dbl(modelos, ~ calcular_auc(predict(.x, prueba, type = "response"),
                                prueba$churn)) %>%
  round(3)

round(exp(coef(modelos[["Con cliente_nuevo"]])["cliente_nuevo"]), 2)
\end{rcode}

\begin{rsalida}
          Original  Con cliente_nuevo Sin genero ni edad 
             0.860              0.862              0.861 
cliente_nuevo 
         2.36 
\end{rsalida}

Agregar \codigo{cliente\_nuevo} sube el AUC de 0.860 a 0.862 y su razón de
momios es 2.36: estar en los primeros 6 meses multiplica por 2.36 los momios
de abandono, además del efecto de la antigüedad. La mejora en AUC es pequeña
porque la antigüedad lineal ya capturaba parte del efecto, pero la variable
tiene sentido de negocio (reforzar el \emph{onboarding}). Quitar género y edad
casi no cambia el AUC (0.861): el modelo más justo y más simple predice igual
de bien, así que es el que conviene usar.
\end{solucion}

\begin{solucion}{m7-7}
Reproducimos datos, división y modelo logístico, y recalculamos los costos con
el nuevo precio del descuento:

\begin{rcode}
library(tidyverse)
library(caret)

set.seed(123)
n_clientes <- 2000
clientes_churn <- tibble(
  cliente_id       = sprintf("C%04d", 1:n_clientes),
  antiguedad_meses = sample(1:72, n_clientes, replace = TRUE),
  plan             = sample(c("Básico", "Estándar", "Premium"),
                            n_clientes, replace = TRUE,
                            prob = c(0.5, 0.3, 0.2)),
  uso_mensual      = rpois(n_clientes, 8),
  quejas_6m        = rpois(n_clientes, 0.8),
  pagos_tardios    = rpois(n_clientes, 0.4),
  genero           = sample(c("F", "M"), n_clientes, replace = TRUE),
  edad             = sample(18:70, n_clientes, replace = TRUE)
) %>%
  mutate(
    efecto_plan = case_when(plan == "Básico"   ~ 0,
                            plan == "Estándar" ~ -0.6,
                            TRUE               ~ -1.3),
    riesgo = 2.0 - 0.05 * antiguedad_meses +
      1.0 * (antiguedad_meses <= 6) + 1.0 * quejas_6m -
      0.35 * uso_mensual + 0.8 * pagos_tardios + efecto_plan,
    churn = if_else(runif(n_clientes) < plogis(riesgo), "Si", "No"),
    churn = factor(churn, levels = c("No", "Si")),
    plan  = factor(plan, levels = c("Básico", "Estándar", "Premium"))
  ) %>%
  select(-efecto_plan, -riesgo)

set.seed(123)
filas_ent <- createDataPartition(clientes_churn$churn, p = 0.7,
                                 list = FALSE)
ent    <- clientes_churn[filas_ent, ]
prueba <- clientes_churn[-filas_ent, ]

modelo_log <- glm(churn ~ antiguedad_meses + plan + uso_mensual +
                    quejas_6m + pagos_tardios + genero + edad,
                  data = ent, family = binomial)
prob <- predict(modelo_log, prueba, type = "response")
real <- prueba$churn == "Si"

costo_descuento <- 800
costo_vp <- costo_descuento + 0.5 * 3000
costo_fp <- costo_descuento
costo_fn <- 3000

costos <- map_dfr(seq(0.05, 0.95, by = 0.05), function(u) {
  pred <- prob >= u
  tibble(umbral = u, llamadas = sum(pred),
         costo_total = sum(pred & real) * costo_vp +
           sum(pred & !real) * costo_fp + sum(!pred & real) * costo_fn)
})
costos %>% arrange(costo_total) %>% head(3)
\end{rcode}

\begin{rsalida}
# A tibble: 3 x 3
  umbral llamadas costo_total
   <dbl>    <int>       <dbl>
1   0.6        74      381700
2   0.65       61      381800
3   0.7        52      382100
\end{rsalida}

Con un descuento de \$800, el umbral óptimo sube de 0.15 a 0.6 y solo se
llamaría a 74 clientes de prueba (costo total de \$381{,}700). La razón: cada
llamada ahora es más cara y el ahorro al retener a un cliente (la mitad de
\$3{,}000 menos \$800) es mucho menor, así que solo vale la pena llamar a
quien tiene una probabilidad alta de irse. Observa que los tres mejores
umbrales tienen costos casi idénticos: la curva de costos es plana en esa
zona.
\end{solucion}

\begin{solucion}{m7-8}
Reproducimos la serie mensual, la agregamos a trimestres con
\codigo{aggregate()} y comparamos los modelos:

\begin{rcode}
library(tidyverse)
library(forecast)

set.seed(123)
meses <- seq(as.Date("2021-01-01"), as.Date("2025-12-01"), by = "month")
n_meses <- length(meses)
factor_mes <- c(0.88, 0.85, 0.95, 0.97, 1.04, 0.96,
                0.97, 1.00, 0.93, 0.98, 1.15, 1.32)
factor_mes <- factor_mes / mean(factor_mes)
unidades <- round((1500 + 15 * (1:n_meses)) * factor_mes[month(meses)] *
                    (1 + rnorm(n_meses, 0, 0.03)))
serie <- ts(unidades, start = c(2021, 1), frequency = 12)

# Suma por trimestre: frecuencia 4
serie_trim <- aggregate(serie, nfrequency = 4, FUN = sum)
serie_trim

ent    <- window(serie_trim, end = c(2024, 4))
prueba <- window(serie_trim, start = c(2025, 1))

pronosticos <- list(
  `Ingenuo estacional` = snaive(ent, h = 4),
  ETS   = forecast(ets(ent), h = 4),
  ARIMA = forecast(auto.arima(ent), h = 4)
)
comparacion <- map_dfr(pronosticos, function(p) {
  acc <- accuracy(p, prueba)["Test set", ]
  tibble(MAE = acc["MAE"], RMSE = acc["RMSE"], MAPE = acc["MAPE"])
}, .id = "modelo") %>%
  mutate(across(where(is.numeric), ~ round(.x, 1))) %>%
  arrange(MAPE)
comparacion
\end{rcode}

\begin{rsalida}
     Qtr1 Qtr2 Qtr3 Qtr4
2021 4139 4765 4626 5822
2022 4579 5229 5176 6240
2023 4999 5757 5797 7156
2024 5557 6200 6399 7468
2025 6089 6903 6774 8305
# A tibble: 3 x 4
  modelo               MAE  RMSE  MAPE
  <chr>              <dbl> <dbl> <dbl>
1 ETS                 96.2  105.   1.4
2 ARIMA              163.   196.   2.2
3 Ingenuo estacional 612.   636.   8.6
\end{rsalida}

\codigo{aggregate(serie, nfrequency = 4, FUN = sum)} suma los tres meses de
cada trimestre. En trimestres, ETS es el mejor modelo con un MAPE de 1.4\%,
seguido de ARIMA (2.2\%); el ingenuo estacional se equivoca 8.6\% porque no
considera el crecimiento. Reentrenamos el ganador con toda la serie:

\begin{rcode}
# Reentrenamos el ganador (menor MAPE) con toda la serie trimestral
ganador <- comparacion$modelo[1]
modelo_final <- switch(ganador,
                       ETS   = ets(serie_trim),
                       ARIMA = auto.arima(serie_trim),
                       snaive(serie_trim, h = 4))
ganador
forecast(modelo_final, h = 4)
\end{rcode}

\begin{rsalida}
[1] "ETS"
        Point Forecast    Lo 80    Hi 80    Lo 95    Hi 95
2026 Q1       6504.441 6355.576 6653.305 6276.772 6732.109
2026 Q2       7369.522 7200.758 7538.287 7111.419 7627.625
2026 Q3       7303.435 7136.086 7470.783 7047.497 7559.372
2026 Q4       8855.932 8652.896 9058.969 8545.414 9166.450
\end{rsalida}

El código elige el modelo automáticamente con \codigo{switch()} según el
ganador. Para 2026 se esperan 6{,}504 unidades en el primer trimestre y 8{,}856
en el cuarto (entre 8{,}653 y 9{,}059 con 80\% de confianza). Con datos
trimestrales hay menos observaciones (20), pero los errores relativos son
menores porque al sumar tres meses el ruido mensual se compensa.
\end{solucion}

\begin{solucion}{m7-9}
Reproducimos los tickets, marcamos por ticket si tiene cada producto y
contamos:

\begin{rcode}
library(tidyverse)

set.seed(123)
n_tickets <- 3000
tickets_ancho <- tibble(
  ticket_id = 1:n_tickets,
  Laptop    = rbinom(n_tickets, 1, 0.08),
  Monitor   = rbinom(n_tickets, 1, 0.07),
  Impresora = rbinom(n_tickets, 1, 0.05),
  Audifonos = rbinom(n_tickets, 1, 0.15),
  USB       = rbinom(n_tickets, 1, 0.18),
  Papel     = rbinom(n_tickets, 1, 0.20),
  Teclado   = rbinom(n_tickets, 1, 0.08)
) %>%
  mutate(
    Mouse = if_else(Laptop == 1, rbinom(n(), 1, 0.60),
                    rbinom(n(), 1, 0.08)),
    Funda = if_else(Laptop == 1, rbinom(n(), 1, 0.40),
                    rbinom(n(), 1, 0.02)),
    `Cable HDMI` = if_else(Monitor == 1, rbinom(n(), 1, 0.50),
                           rbinom(n(), 1, 0.04)),
    Tinta = if_else(Impresora == 1, rbinom(n(), 1, 0.70),
                    rbinom(n(), 1, 0.06))
  )
tickets <- tickets_ancho %>%
  pivot_longer(-ticket_id, names_to = "producto", values_to = "compro") %>%
  filter(compro == 1) %>%
  select(-compro)

# Una fila por ticket con TRUE/FALSE para los productos de interés
por_ticket <- tickets %>%
  group_by(ticket_id) %>%
  summarise(laptop = "Laptop" %in% producto,
            mouse  = "Mouse" %in% producto,
            funda  = "Funda" %in% producto, .groups = "drop")
n_total <- nrow(por_ticket)

soporte_funda <- mean(por_ticket$funda)
tres <- por_ticket %>%
  reframe(regla = c("Laptop+Mouse -> Funda", "Laptop -> Funda"),
            tickets_x  = c(sum(laptop & mouse), sum(laptop)),
            tickets_xy = c(sum(laptop & mouse & funda),
                           sum(laptop & funda)))
tres %>%
  mutate(soporte   = round(tickets_xy / n_total, 3),
         confianza = round(tickets_xy / tickets_x, 3),
         lift      = round(tickets_xy / tickets_x / soporte_funda, 2))
\end{rcode}

\begin{rsalida}
# A tibble: 2 x 6
  regla                 tickets_x tickets_xy soporte confianza  lift
  <chr>                     <int>      <int>   <dbl>     <dbl> <dbl>
1 Laptop+Mouse -> Funda       136         51   0.026     0.375  5.14
2 Laptop -> Funda             225         82   0.042     0.364  5   
\end{rsalida}

Saber que el cliente ya lleva mouse casi no cambia la probabilidad de que
lleve funda: la confianza pasa de 36.4\% a 37.5\% y el lift de 5.0 a 5.14.
Tiene sentido, porque en la simulación la funda depende solo de la laptop.
Una regla de tres productos solo aporta si su confianza es claramente mayor
que la de la regla más simple; si no, quédate con la simple, que además tiene
más soporte (4.2\% contra 2.6\% de los tickets). Nota: usamos
\codigo{reframe()} en lugar de \codigo{summarise()} porque el resultado tiene
dos filas.
\end{solucion}

\begin{solucion}{m7-10}
Organizamos el reto en funciones: una para simular clientes (la usamos para
el histórico y para el lote nuevo), el entrenamiento con \paquete{caret}, la
elección del modelo y la función de puntuación.

\begin{rcode}
library(tidyverse)
library(caret)
library(randomForest)

simular_clientes <- function(n, prefijo = "C") {
  tibble(
    cliente_id       = sprintf("%s%05d", prefijo, 1:n),
    antiguedad_meses = sample(1:72, n, replace = TRUE),
    plan             = sample(c("Básico", "Estándar", "Premium"), n,
                              replace = TRUE, prob = c(0.5, 0.3, 0.2)),
    uso_mensual      = rpois(n, 8),
    quejas_6m        = rpois(n, 0.8),
    pagos_tardios    = rpois(n, 0.4),
    genero           = sample(c("F", "M"), n, replace = TRUE),
    edad             = sample(18:70, n, replace = TRUE)
  ) %>%
    mutate(
      efecto_plan = case_when(plan == "Básico"   ~ 0,
                              plan == "Estándar" ~ -0.6,
                              TRUE               ~ -1.3),
      riesgo = 2.0 - 0.05 * antiguedad_meses +
        1.0 * (antiguedad_meses <= 6) + 1.0 * quejas_6m -
        0.35 * uso_mensual + 0.8 * pagos_tardios + efecto_plan,
      churn = if_else(runif(n) < plogis(riesgo), "Si", "No"),
      churn = factor(churn, levels = c("Si", "No")),  # "Si" primero: caret
      plan  = factor(plan, levels = c("Básico", "Estándar", "Premium")),
      cuota_mensual = case_when(plan == "Básico"   ~ 199,
                                plan == "Estándar" ~ 349,
                                TRUE               ~ 599)
    ) %>%
    select(-efecto_plan, -riesgo)
}

set.seed(2024)
historico <- simular_clientes(3000)
set.seed(2024)
filas_ent <- createDataPartition(historico$churn, p = 0.7, list = FALSE)
ent    <- historico[filas_ent, ]
prueba <- historico[-filas_ent, ]

# Sin género ni edad (sección de ética)
formula_churn <- churn ~ antiguedad_meses + plan + uso_mensual +
  quejas_6m + pagos_tardios
control <- trainControl(method = "cv", number = 5, classProbs = TRUE,
                        summaryFunction = twoClassSummary)
set.seed(2024)
m_log <- train(formula_churn, data = ent, method = "glm",
               family = binomial, trControl = control, metric = "ROC")
set.seed(2024)
m_rf <- train(formula_churn, data = ent, method = "rf", ntree = 200,
              trControl = control, metric = "ROC",
              tuneGrid = data.frame(mtry = 2:3))

calcular_auc <- function(prob, real, positivo = "Si") {
  p_pos <- prob[real == positivo]
  p_neg <- prob[real != positivo]
  mean(outer(p_pos, p_neg, ">") + 0.5 * outer(p_pos, p_neg, "=="))
}
auc_prueba <- c(
  Logistica = calcular_auc(predict(m_log, prueba, type = "prob")[, "Si"],
                           prueba$churn),
  RandomForest = calcular_auc(predict(m_rf, prueba, type = "prob")[, "Si"],
                              prueba$churn))
round(auc_prueba, 3)

mejor <- list(Logistica = m_log, RandomForest = m_rf)[[
  names(which.max(auc_prueba))]]
dir.create("resultados", showWarnings = FALSE)
saveRDS(mejor, "resultados/modelo_churn_reto.rds")
\end{rcode}

\begin{rsalida}
   Logistica RandomForest 
       0.873        0.843 
\end{rsalida}

La regresión logística gana en prueba (AUC de 0.873 contra 0.843 del random
forest), así que es el modelo que se guarda en
\archivo{resultados/modelo\_churn\_reto.rds}. El código elige al ganador con
\codigo{which.max()}: si mañana gana el bosque, el resto del sistema sigue
funcionando igual, porque ambos modelos de \paquete{caret} se usan con
\codigo{predict(..., type = "prob")}.

\begin{rcode}
puntuar <- function(nuevos, modelo, top = 100) {
  nuevos %>%
    mutate(prob_churn = predict(modelo, newdata = nuevos,
                                type = "prob")[, "Si"],
           perdida_esperada = prob_churn * cuota_mensual * 12) %>%
    arrange(desc(perdida_esperada)) %>%
    slice_head(n = top) %>%
    mutate(prioridad = row_number()) %>%
    select(prioridad, cliente_id, plan, quejas_6m, prob_churn,
           perdida_esperada)
}

set.seed(99)
lote_nuevo <- simular_clientes(1000, prefijo = "N") %>% select(-churn)

modelo <- readRDS("resultados/modelo_churn_reto.rds")
top100 <- puntuar(lote_nuevo, modelo)
write_csv(top100, "resultados/top100_churn.csv")

top100 %>%
  mutate(across(c(prob_churn, perdida_esperada), ~ round(.x, 2))) %>%
  head(5)
nrow(read_csv("resultados/top100_churn.csv", show_col_types = FALSE))
\end{rcode}

\begin{rsalida}
# A tibble: 5 x 6
  prioridad cliente_id plan    quejas_6m prob_churn perdida_esperada
      <int> <chr>      <fct>       <int>      <dbl>            <dbl>
1         1 N00330     Premium         3       0.94            6765.
2         2 N00043     Premium         3       0.89            6399.
3         3 N00814     Premium         2       0.87            6268.
4         4 N00456     Premium         2       0.75            5368.
5         5 N00904     Premium         1       0.67            4787.
[1] 100
\end{rsalida}

La función carga el modelo guardado (como se haría en producción), calcula
la pérdida esperada de cada cliente del lote y exporta los 100 más urgentes;
releer el CSV confirma que tiene 100 filas. Los primeros lugares son clientes
Premium con varias quejas y probabilidades de 0.67 a 0.94. Para ponerlo en
producción faltaría programar el script (Módulo~\ref{cap:reportes}) y
monitorear la deriva de las variables cada semana.
\end{solucion}


\section{Módulo 8. Reportes automáticos}
% ============================================================================
% SOLUCIONES DEL MÓDULO 8. REPORTES AUTOMÁTICOS CON R MARKDOWN
% ============================================================================

\begin{solucion}{m8-1}
El bloque que lee los datos lleva \codigo{include = FALSE} para que no se
vea; la fecha es código en línea en el YAML, y la frase usa código en línea
con \codigo{nrow()} y \codigo{sum()} sobre la columna lógica
\codigo{es\_premium} (cada \codigo{TRUE} cuenta como 1). Como el archivo se
guarda en \archivo{reportes/}, los datos se buscan también en
\archivo{../datasets}.

\begin{textocodigo}[title=reportes/ej1-clientes.Rmd]
---
title: "Reporte de clientes"
author: "Tu nombre"
date: "`r format(Sys.Date(), '%d/%m/%Y')`"
output: html_document
---

```{r datos, include = FALSE}
library(dplyr)
library(readr)
library(knitr)
carpeta <- if (dir.exists("datasets")) "datasets" else "../datasets"
clientes <- read_csv(file.path(carpeta, "clientes.csv"),
                     show_col_types = FALSE)
```

## Resumen

La base tiene **`r nrow(clientes)` clientes**, de los cuales
**`r sum(clientes$es_premium)`** son premium
(`r round(mean(clientes$es_premium) * 100, 1)` %).

## Clientes por ciudad

```{r tabla-ciudades, echo = FALSE}
clientes %>%
  count(ciudad, sort = TRUE) %>%
  kable(col.names = c("Ciudad", "Clientes"),
        caption = "Clientes por ciudad")
```
\end{textocodigo}

Lo generamos con \menu{Knit} (o con
\codigo{rmarkdown::render("reportes/ej1-clientes.Rmd")}). Para comprobar
qué dirá el reporte, ejecutamos los mismos cálculos en la consola:

\begin{rcode}
library(tidyverse)
library(knitr)
clientes <- read_csv("datasets/clientes.csv", show_col_types = FALSE)
\end{rcode}

\begin{rcode}
nrow(clientes)
sum(clientes$es_premium)
round(mean(clientes$es_premium) * 100, 1)
clientes %>%
  count(ciudad, sort = TRUE) %>%
  kable(format = "pipe", col.names = c("Ciudad", "Clientes"))
\end{rcode}
\begin{rsalida}
[1] 200
[1] 73
[1] 36.5


|Ciudad      | Clientes|
|:-----------|--------:|
|Puebla      |       28|
|CDMX        |       27|
|Guadalajara |       22|
|Cancún      |       20|
|Toluca      |       20|
|León        |       18|
|Mérida      |       17|
|Querétaro   |       17|
|Tijuana     |       17|
|Monterrey   |       14|
\end{rsalida}

El reporte generado dice que hay 200 clientes, 73 premium
(36.5\,\%), y la tabla empieza con Puebla (28 clientes) y CDMX (27).
\end{solucion}

\begin{solucion}{m8-2}
Unimos las transacciones con la categoría del producto, resumimos, ordenamos
\emph{antes} de formatear (porque al formatear las columnas se vuelven texto)
y damos formato de negocio con \paquete{scales}.

\begin{rcode}
library(tidyverse)
library(knitr)
library(scales)

transacciones <- read_csv("datasets/transacciones.csv",
                          show_col_types = FALSE)
productos <- read_csv("datasets/productos.csv", show_col_types = FALSE)

por_categoria <- transacciones %>%
  left_join(productos %>% select(producto_id, categoria),
            by = "producto_id") %>%
  group_by(categoria) %>%
  summarise(ventas = sum(total_transaccion),
            unidades = sum(cantidad), .groups = "drop") %>%
  arrange(desc(ventas)) %>%                  # ordenar con números
  mutate(participacion = percent(ventas / sum(ventas), accuracy = 0.1),
         ventas = dollar(ventas, accuracy = 1),
         unidades = comma(unidades))

kable(por_categoria, format = "pipe",
      col.names = c("Categoría", "Ventas", "Unidades", "% del total"),
      align = c("l", "r", "r", "r"),
      caption = "Ventas 2023 por categoría")
\end{rcode}
\begin{rsalida}

Table: Ventas 2023 por categoría

|Categoría   |   Ventas| Unidades| % del total|
|:-----------|--------:|--------:|-----------:|
|Belleza     | $182,143|      551|       19.0%|
|Electrónica | $132,672|      403|       13.8%|
|Juguetes    | $125,093|      356|       13.0%|
|Alimentos   | $121,543|      356|       12.7%|
|Automotriz  | $104,085|      325|       10.8%|
|Deportes    |  $86,171|      280|        9.0%|
|Hogar       |  $63,407|      187|        6.6%|
|Mascotas    |  $57,668|      201|        6.0%|
|Libros      |  $57,162|      204|        6.0%|
|Ropa        |  $30,048|      106|        3.1%|
\end{rsalida}

Belleza es la categoría que más vende: \$182,143, el 19.0\,\% del total, seguida
de Electrónica (13.8\,\%). Ropa queda al final con solo el 3.1\,\%. Observa
la alineación: los números a la derecha permiten comparar magnitudes de un
vistazo.
\end{solucion}

\begin{solucion}{m8-3}
La función calcula primero el cumplimiento (\codigo{ventas / meta}) y con
\codigo{if}/\codigo{else if}/\codigo{else} elige la frase. Para aplicarla a
cada trimestre usamos \codigo{map\_chr()}, que llama a la función con cada
valor de ventas.

\begin{rcode}
library(tidyverse)
library(scales)

frase_meta <- function(ventas, meta) {
  cumplimiento <- ventas / meta
  montos <- paste0("(ventas de ", dollar(ventas, accuracy = 1),
                   " contra meta de ", dollar(meta, accuracy = 1), ").")
  if (ventas >= meta) {
    paste("Se superó la meta en",
          percent(cumplimiento - 1, accuracy = 0.1), montos)
  } else if (cumplimiento >= 0.9) {
    paste("Se alcanzó el", percent(cumplimiento, accuracy = 0.1),
          "de la meta", montos)
  } else {
    paste("Alerta: solo se alcanzó el",
          percent(cumplimiento, accuracy = 0.1), "de la meta", montos)
  }
}

# Los tres casos
frase_meta(110000, 100000)
frase_meta(95000, 100000)
frase_meta(70000, 100000)

# Aplicación a los trimestres de 2023
transacciones <- read_csv("datasets/transacciones.csv",
                          show_col_types = FALSE)
trimestres <- transacciones %>%
  group_by(trimestre) %>%
  summarise(ventas = sum(total_transaccion), .groups = "drop") %>%
  mutate(frase = map_chr(ventas, frase_meta, meta = 220000))
cat(paste0(trimestres$trimestre, ": ", trimestres$frase), sep = "\n")
\end{rcode}
\begin{rsalida}
[1] "Se superó la meta en 10.0% (ventas de $110,000 contra meta de $100,000)."
[1] "Se alcanzó el 95.0% de la meta (ventas de $95,000 contra meta de $100,000)."
[1] "Alerta: solo se alcanzó el 70.0% de la meta (ventas de $70,000 contra meta de $100,000)."
Q1: Se superó la meta en 8.8% (ventas de $239,390 contra meta de $220,000).
Q2: Se superó la meta en 8.8% (ventas de $239,457 contra meta de $220,000).
Q3: Se superó la meta en 6.5% (ventas de $234,241 contra meta de $220,000).
Q4: Se superó la meta en 12.2% (ventas de $246,905 contra meta de $220,000).
\end{rsalida}

Los cuatro trimestres superaron la meta de \$220,000; el mejor fue Q4, con
12.2\,\% por encima. Si cambias la meta a \$240,000, Q1, Q2 y Q3 pasan a
``Se alcanzó el \ldots'' sin tocar el texto: la frase depende solo de los
datos.
\end{solucion}

\begin{solucion}{m8-4}
El bloque \codigo{configuracion} fija las opciones globales; el bloque de
cálculo lleva \codigo{include = FALSE}; la gráfica define su tamaño y su
pie; y el bloque final usa \codigo{results = "asis"} para que el texto de
\codigo{cat()} se convierta en una lista.

\begin{textocodigo}[title=reportes/ej4-opciones.Rmd]
---
title: "Ventas mensuales 2023"
output:
  html_document: default
  word_document: default
---

```{r configuracion, include = FALSE}
knitr::opts_chunk$set(echo = FALSE, message = FALSE, warning = FALSE)
library(dplyr)
library(readr)
library(ggplot2)
library(scales)
```

```{r calculos, include = FALSE}
carpeta <- if (dir.exists("datasets")) "datasets" else "../datasets"
ventas <- read_csv(file.path(carpeta, "ventas_retail.csv"),
                   show_col_types = FALSE)
por_mes <- ventas %>%
  group_by(mes) %>%
  summarise(ventas = sum(total), .groups = "drop")
```

## Tendencia

```{r grafica, fig.cap = "Ventas por mes, 2023", fig.width = 7,
    fig.height = 3, out.width = "80%"}
ggplot(por_mes, aes(x = mes, y = ventas, group = 1)) +
  geom_line(color = "#2a78d6", linewidth = 0.8) +
  scale_y_continuous(labels = dollar) +
  labs(x = NULL, y = "Ventas") +
  theme_minimal() +
  theme(panel.grid.minor = element_blank(),
        axis.text.x = element_text(angle = 45, hjust = 1))
```

## Los tres mejores meses

```{r top-meses, results = "asis"}
top3 <- por_mes %>% arrange(desc(ventas)) %>% slice_head(n = 3)
cat(paste0("- **", top3$mes, "**: ", dollar(top3$ventas, accuracy = 1)),
    sep = "\n")
```
\end{textocodigo}

Se genera en los dos formatos con:

\begin{rnoejecutar}
rmarkdown::render("reportes/ej4-opciones.Rmd", output_format = "all")
\end{rnoejecutar}

La lista que escribe el último bloque es exactamente la que produce este
código en la consola:

\begin{rcode}
library(tidyverse)
library(scales)
ventas <- read_csv("datasets/ventas_retail.csv", show_col_types = FALSE)
por_mes <- ventas %>%
  group_by(mes) %>%
  summarise(ventas = sum(total), .groups = "drop")
top3 <- por_mes %>% arrange(desc(ventas)) %>% slice_head(n = 3)
cat(paste0("- **", top3$mes, "**: ", dollar(top3$ventas, accuracy = 1)),
    sep = "\n")
\end{rcode}
\begin{rsalida}
- **2023-07**: $54,056
- **2023-11**: $50,398
- **2023-05**: $48,108
\end{rsalida}

En el HTML y en el Word esas líneas aparecen como una lista con viñetas y
los meses en negritas, sin recuadro gris ni \codigo{\#\#}.
\end{solucion}

\begin{solucion}{m8-5}
El parámetro se declara en el YAML y se usa como \codigo{params\$categoria}
en el filtro, en el título y en las frases.

\begin{textocodigo}[title=reportes/ej5-categoria.Rmd]
---
title: "Reporte de la categoría `r params$categoria`"
output: html_document
params:
  categoria: "Electrónica"
---

```{r configuracion, include = FALSE}
knitr::opts_chunk$set(echo = FALSE, message = FALSE, warning = FALSE)
library(dplyr)
library(readr)
library(ggplot2)
library(scales)
library(knitr)
carpeta <- if (dir.exists("datasets")) "datasets" else "../datasets"
transacciones <- read_csv(file.path(carpeta, "transacciones.csv"),
                          show_col_types = FALSE)
productos <- read_csv(file.path(carpeta, "productos.csv"),
                      show_col_types = FALSE)
ventas_cat <- transacciones %>%
  left_join(productos, by = "producto_id") %>%
  filter(categoria == params$categoria)
if (nrow(ventas_cat) == 0) stop("Sin ventas de ", params$categoria)
```

En 2023, la categoría **`r params$categoria`** vendió
**`r dollar(sum(ventas_cat$total_transaccion), accuracy = 1)`** en
`r nrow(ventas_cat)` transacciones.

```{r top-productos}
ventas_cat %>%
  group_by(nombre_producto) %>%
  summarise(unidades = sum(cantidad),
            ventas = sum(total_transaccion), .groups = "drop") %>%
  arrange(desc(ventas)) %>%
  slice_head(n = 5) %>%
  mutate(ventas = dollar(ventas, accuracy = 1)) %>%
  kable(col.names = c("Producto", "Unidades", "Ventas"),
        align = c("l", "r", "r"),
        caption = "Productos más vendidos")
```

```{r grafica, fig.cap = "Ventas por mes", fig.height = 3}
ventas_cat %>%
  group_by(mes) %>%
  summarise(ventas = sum(total_transaccion), .groups = "drop") %>%
  ggplot(aes(x = mes, y = ventas)) +
  geom_col(fill = "#2a78d6", width = 0.7) +
  scale_y_continuous(labels = dollar) +
  labs(x = NULL, y = "Ventas") +
  theme_minimal() +
  theme(panel.grid.minor = element_blank(),
        axis.text.x = element_text(angle = 45, hjust = 1))
```
\end{textocodigo}

Para generar dos versiones, se recorre un vector de categorías y se da un
\codigo{output\_file} distinto a cada una (si no, la segunda sobrescribe a
la primera):

\begin{rnoejecutar}
for (cat_actual in c("Ropa", "Hogar")) {
  rmarkdown::render(
    "reportes/ej5-categoria.Rmd",
    params        = list(categoria = cat_actual),
    output_file   = paste0("categoria_", tolower(cat_actual), ".html"),
    output_dir    = "resultados",
    knit_root_dir = getwd(),
    envir         = new.env()
  )
}
\end{rnoejecutar}

Las frases de cada reporte coinciden con este cálculo:

\begin{rcode}
library(tidyverse)
library(scales)
transacciones <- read_csv("datasets/transacciones.csv",
                          show_col_types = FALSE)
productos <- read_csv("datasets/productos.csv", show_col_types = FALSE)
transacciones %>%
  left_join(productos, by = "producto_id") %>%
  filter(categoria %in% c("Ropa", "Hogar")) %>%
  group_by(categoria) %>%
  summarise(ventas = dollar(sum(total_transaccion), accuracy = 1),
            transacciones = n(), .groups = "drop")
\end{rcode}
\begin{rsalida}
# A tibble: 2 x 3
  categoria ventas  transacciones
  <chr>     <chr>           <int>
1 Hogar     $63,407            64
2 Ropa      $30,048            38
\end{rsalida}

Se obtienen \archivo{resultados/categoria\_ropa.html} y
\archivo{resultados/categoria\_hogar.html}, cada uno con su título, su frase,
su tabla y su gráfica.
\end{solucion}

\begin{solucion}{m8-6}
\codigo{expand\_grid()} crea las 40 combinaciones; con \codigo{count()}
contamos las transacciones por tienda y trimestre y las unimos con
\codigo{left\_join()}. Si alguna combinación no tuviera transacciones,
quedaría \codigo{NA}, que convertimos en 0 con \codigo{replace\_na()}.

\begin{rcode}
library(tidyverse)
transacciones <- read_csv("datasets/transacciones.csv",
                          show_col_types = FALSE)
tiendas <- read_csv("datasets/tiendas.csv", show_col_types = FALSE)

limpiar_nombre <- function(x) {
  x <- chartr("áéíóúüñÁÉÍÓÚÜÑ", "aeiouunAEIOUUN", x)
  x <- gsub("[^A-Za-z0-9]+", "-", x)
  tolower(gsub("^-|-$", "", x))
}

conteo <- transacciones %>% count(tienda_id, trimestre,
                                  name = "n_transacciones")

trabajos <- expand_grid(tienda_id = tiendas$tienda_id,
                        trimestre = paste0("Q", 1:4)) %>%
  left_join(tiendas %>% select(tienda_id, nombre_tienda),
            by = "tienda_id") %>%
  left_join(conteo, by = c("tienda_id", "trimestre")) %>%
  mutate(n_transacciones = replace_na(n_transacciones, 0),
         archivo = sprintf("2023-%s_T%02d_%s.pdf", trimestre, tienda_id,
                           limpiar_nombre(nombre_tienda)),
         generar = n_transacciones >= 20)

nrow(trabajos)
trabajos %>% arrange(n_transacciones) %>%
  select(archivo, n_transacciones, generar) %>%
  head(4)
sum(trabajos$generar)
\end{rcode}
\begin{rsalida}
[1] 40
# A tibble: 4 x 3
  archivo                          n_transacciones generar
  <chr>                                      <int> <lgl>  
1 2023-Q2_T07_sucursal-plaza-b.pdf              16 FALSE  
2 2023-Q4_T05_sucursal-oeste.pdf                18 FALSE  
3 2023-Q4_T07_sucursal-plaza-b.pdf              18 FALSE  
4 2023-Q4_T01_sucursal-centro.pdf               19 FALSE  
[1] 34
\end{rsalida}

Hay que generar 40 reportes (10 tiendas $\times$ 4 trimestres). La
combinación con menos transacciones es la Sucursal Plaza B en el segundo
trimestre (16). Seis combinaciones quedan por debajo de 20 y se marcan con
\codigo{generar = FALSE}, así que el bucle generaría 34 reportes; en el
script de automatización bastaría con recorrer
\codigo{filter(trabajos, generar)}.
\end{solucion}

\begin{solucion}{m8-7}
\codigo{split()} divide el detalle en una lista con un data frame por
tienda; los nombres de esa lista serán los nombres de las hojas. Con
\codigo{c()} ponemos la hoja \codigo{Resumen} al principio.

\begin{rcode}
library(tidyverse)
library(writexl)
library(readxl)

transacciones <- read_csv("datasets/transacciones.csv",
                          show_col_types = FALSE)
tiendas <- read_csv("datasets/tiendas.csv", show_col_types = FALSE)

detalle <- transacciones %>%
  left_join(tiendas %>% select(tienda_id, nombre_tienda),
            by = "tienda_id") %>%
  mutate(hoja = str_remove(nombre_tienda, "Sucursal ")) %>%
  arrange(tienda_id, fecha)

resumen <- detalle %>%
  group_by(hoja) %>%
  summarise(ventas = round(sum(total_transaccion), 2),
            transacciones = n(), .groups = "drop") %>%
  arrange(desc(ventas))

# Una hoja por tienda (sin la columna auxiliar "hoja")
por_tienda <- split(detalle %>% select(-hoja), detalle$hoja)

dir.create("resultados", showWarnings = FALSE)
archivo <- "resultados/ventas_por_tienda.xlsx"
write_xlsx(c(list(Resumen = resumen), por_tienda), path = archivo)

# Verificación
hojas <- excel_sheets(archivo)
length(hojas)
hojas
filas <- sapply(hojas[-1], function(h) nrow(read_excel(archivo, sheet = h)))
sum(filas)
\end{rcode}
\begin{rsalida}
[1] 11
 [1] "Resumen" "Centro"  "Este"    "Express" "Norte"   "Oeste"  
 [7] "Outlet"  "Plaza A" "Plaza B" "Premium" "Sur"    
[1] 1000
\end{rsalida}

El libro tiene 11 hojas (el resumen y diez tiendas) y las hojas de tiendas
suman las 1,000 transacciones: no se perdió ni se duplicó ninguna. Observa
que \codigo{split()} ordena las hojas alfabéticamente (Centro, Este,
Express\ldots); si quisieras otro orden, reordena la lista antes de
escribirla.
\end{solucion}

\begin{solucion}{m8-8}
La función valida sus argumentos con \codigo{stop()}, igual que la plantilla
del reporte ejecutivo. El bucle envuelve cada llamada en
\codigo{tryCatch()}: las filas exitosas se guardan en una lista de
resultados y los errores en el registro, así que ninguna combinación inválida
detiene el proceso.

\begin{rcode}
library(tidyverse)
library(scales)

transacciones <- read_csv("datasets/transacciones.csv",
                          show_col_types = FALSE)
tiendas <- read_csv("datasets/tiendas.csv", show_col_types = FALSE)

revisar_tienda <- function(id, mes) {
  if (!id %in% tiendas$tienda_id) stop("La tienda ", id, " no existe.")
  if (!mes %in% transacciones$mes) stop("No hay datos del mes ", mes)
  mes_previo <- format(as.Date(paste0(mes, "-01")) - 1, "%Y-%m")
  if (!mes_previo %in% transacciones$mes) {
    stop("No hay datos del mes anterior (", mes_previo, ")")
  }
  v <- transacciones %>% filter(tienda_id == id)
  actual   <- sum(v$total_transaccion[v$mes == mes])
  anterior <- sum(v$total_transaccion[v$mes == mes_previo])
  variacion <- actual / anterior - 1
  tibble(tienda_id = id,
         tienda = tiendas$nombre_tienda[tiendas$tienda_id == id],
         mes = mes, variacion = variacion,
         estado = case_when(variacion >  0.10 ~ "Crecimiento sólido",
                            variacion >= 0    ~ "Estable",
                            variacion > -0.10 ~ "Ligera caída",
                            TRUE              ~ "Alerta"))
}

pendientes <- tibble(id = c(tiendas$tienda_id, 11, 1),
                     mes = c(rep("2023-06", 10), "2023-06", "2023-01"))
resultados <- list()
registro <- list()

for (fila in seq_len(nrow(pendientes))) {
  id  <- pendientes$id[fila]
  mes <- pendientes$mes[fila]
  estado <- tryCatch({
    resultados[[length(resultados) + 1]] <- revisar_tienda(id, mes)
    "OK"
  }, error = function(e) paste("ERROR:", conditionMessage(e)))
  registro[[fila]] <- tibble(id = id, mes = mes, estado = estado)
}

registro <- bind_rows(registro)
registro %>% count(estado)
registro %>% filter(estado != "OK")

resultados <- bind_rows(resultados)
resultados %>%
  mutate(variacion = percent(variacion, accuracy = 0.1)) %>%
  select(tienda, variacion, estado)

# Texto Markdown para el reporte (iría en un bloque results = "asis")
alertas <- resultados %>% filter(estado == "Alerta")
cat(paste0("- **", alertas$tienda, "**: las ventas cayeron ",
           percent(abs(alertas$variacion), accuracy = 0.1),
           " respecto a mayo."), sep = "\n")
\end{rcode}
\begin{rsalida}
# A tibble: 3 x 2
  estado                                             n
  <chr>                                          <int>
1 ERROR: La tienda 11 no existe.                     1
2 ERROR: No hay datos del mes anterior (2022-12)     1
3 OK                                                10
# A tibble: 2 x 3
     id mes     estado                                        
  <dbl> <chr>   <chr>                                         
1    11 2023-06 ERROR: La tienda 11 no existe.                
2     1 2023-01 ERROR: No hay datos del mes anterior (2022-12)
# A tibble: 10 x 3
   tienda           variacion estado            
   <chr>            <chr>     <chr>             
 1 Sucursal Centro  -48.1%    Alerta            
 2 Sucursal Norte   -73.5%    Alerta            
 3 Sucursal Sur     31.8%     Crecimiento sólido
 4 Sucursal Este    245.5%    Crecimiento sólido
 5 Sucursal Oeste   -41.0%    Alerta            
 6 Sucursal Plaza A -16.3%    Alerta            
 7 Sucursal Plaza B 29.1%     Crecimiento sólido
 8 Sucursal Outlet  612.9%    Crecimiento sólido
 9 Sucursal Premium -8.0%     Ligera caída      
10 Sucursal Express -45.5%    Alerta            
- **Sucursal Centro**: las ventas cayeron 48.1% respecto a mayo.
- **Sucursal Norte**: las ventas cayeron 73.5% respecto a mayo.
- **Sucursal Oeste**: las ventas cayeron 41.0% respecto a mayo.
- **Sucursal Plaza A**: las ventas cayeron 16.3% respecto a mayo.
- **Sucursal Express**: las ventas cayeron 45.5% respecto a mayo.
\end{rsalida}

De las 12 combinaciones, 10 se revisaron y 2 quedaron en el registro con su
causa: la tienda 11 no existe y enero de 2023 no tiene mes anterior en los
datos. En junio, cinco tiendas están en alerta y la Sucursal Outlet creció
612.9\,\%: una variación tan grande con tan pocas transacciones merece
revisar los datos antes de celebrarla. El texto final está listo para
insertarse en un reporte con \codigo{results = "asis"}.
\end{solucion}


\section{Módulo 9. Proyecto final: retos de extensión}
% ============================================================================
% SOLUCIONES - MÓDULO 9 (PROYECTO FINAL: RETOS DE EXTENSIÓN)
% ============================================================================

\begin{solucion}{p-1}
La recompra se mide sobre los clientes \emph{activos} (los que compraron al
menos una vez). Para los días entre compras ordenamos las compras de cada
cliente por fecha y restamos la fecha anterior con \codigo{lag()}.

\begin{rcode}
library(tidyverse)
transacciones <- read_csv("datasets/transacciones.csv",
                          show_col_types = FALSE)
\end{rcode}

Al cargar \paquete{tidyverse} aparece su mensaje de bienvenida (no lo
copiamos aquí). Ahora los tres cálculos:

\begin{rcode}
# (a) Porcentaje de clientes activos con 2 o más compras
compras_cliente <- transacciones %>% count(cliente_id, name = "compras")
compras_cliente %>%
  summarise(activos = n(),
            recompran = sum(compras >= 2),
            tasa_recompra = round(100 * mean(compras >= 2), 1))

# (b) Días entre compras consecutivas del mismo cliente
intervalos <- transacciones %>%
  arrange(cliente_id, fecha) %>%
  group_by(cliente_id) %>%
  mutate(dias_desde_anterior = as.numeric(fecha - lag(fecha))) %>%
  ungroup() %>%
  filter(!is.na(dias_desde_anterior))

quantile(intervalos$dias_desde_anterior, c(0.25, 0.5, 0.75))

# (c) Clientes con al menos una recompra en 30 días o menos
intervalos %>%
  group_by(cliente_id) %>%
  summarise(recompra_30 = any(dias_desde_anterior <= 30)) %>%
  summarise(clientes = sum(recompra_30))
\end{rcode}

\begin{rsalida}
# A tibble: 1 x 3
  activos recompran tasa_recompra
    <int>     <int>         <dbl>
1     198       192            97
25% 50% 75% 
 19  41  82 
# A tibble: 1 x 1
  clientes
     <int>
1      140
\end{rsalida}

\textbf{Interpretación:} 192 de los 198 clientes activos (97\%) compraron
dos veces o más, así que ``conseguir la segunda compra'' no es el problema de
TechRetail. La mediana entre compras es de 41 días (la mitad de las
recompras ocurre entre 19 y 82 días) y 140 clientes volvieron alguna vez en
30 días o menos. Una campaña de segunda compra aportaría poco; es más útil
una campaña que \emph{acorte} el intervalo (por ejemplo, un cupón válido 30
días) y medir si la mediana baja de 41 días.
\end{solucion}

\begin{solucion}{p-2}
Unimos las ventas por sucursal con \archivo{tiendas.csv}, dividimos entre
m\textsuperscript{2} y empleados y usamos \codigo{min\_rank(desc())} para
los rankings.

\begin{rcode}
library(tidyverse)
transacciones <- read_csv("datasets/transacciones.csv",
                          show_col_types = FALSE)
tiendas <- read_csv("datasets/tiendas.csv", show_col_types = FALSE)

productividad <- transacciones %>%
  group_by(tienda_id) %>%
  summarise(ventas = sum(total_transaccion), .groups = "drop") %>%
  left_join(tiendas, by = "tienda_id") %>%
  mutate(ventas_m2       = ventas / tamano_m2,
         ventas_empleado = ventas / empleados,
         rank_ventas     = min_rank(desc(ventas)),
         rank_m2         = min_rank(desc(ventas_m2)),
         rank_empleado   = min_rank(desc(ventas_empleado)),
         cambio          = rank_ventas - rank_m2) %>%
  arrange(rank_m2)

productividad %>%
  transmute(tienda = str_remove(nombre_tienda, "Sucursal "),
            m2 = tamano_m2, ventas = round(ventas),
            por_m2 = round(ventas_m2), por_empleado = round(ventas_empleado),
            rank_ventas, rank_m2, cambio)
\end{rcode}

\begin{rsalida}
# A tibble: 10 x 8
   tienda     m2 ventas por_m2 por_empleado rank_ventas rank_m2 cambio
   <chr>   <dbl>  <dbl>  <dbl>        <dbl>       <int>   <int>  <int>
 1 Express   300 102791    343         8566           4       1      3
 2 Outlet    350 114377    327         7625           1       2     -1
 3 Este      400  88420    221         4912           7       3      4
 4 Norte     450  89024    198         4451           6       4      2
 5 Centro    500  88414    177         3537           8       5      3
 6 Sur       600 105305    176         3510           2       6     -4
 7 Plaza B   480  80236    167         3821          10       7      3
 8 Oeste     550  84672    154         3849           9       8      1
 9 Plaza A   700 105216    150         3006           3       9     -6
10 Premium   800 101538    127         2538           5      10     -5
\end{rsalida}

\begin{rcode}
color_resalte <- "#eb6834"; color_gris <- "#b5b3ad"
productividad %>%
  mutate(mejor = rank_m2 == 1) %>%
  ggplot(aes(x = ventas_m2, y = fct_reorder(nombre_tienda, ventas_m2),
             fill = mejor)) +
  geom_col(width = 0.7, show.legend = FALSE) +
  scale_fill_manual(values = c(`TRUE` = color_resalte,
                               `FALSE` = color_gris)) +
  scale_x_continuous(labels = scales::dollar) +
  labs(title = "Ventas por metro cuadrado, 2023",
       x = "Ventas por m²", y = NULL) +
  theme_minimal() +
  theme(panel.grid.minor = element_blank())
\end{rcode}

\textbf{Interpretación:} la Express (300~m\textsuperscript{2}) es la más
productiva por metro cuadrado (\$343) y por empleado (\$8,566), seguida de la
Outlet. La Plaza A es la que más cae: tercera en ventas totales y novena por
m\textsuperscript{2} (6 lugares), y la Premium queda última por
m\textsuperscript{2} (\$127). Si la dirección piensa en expansión, los
formatos chicos producen casi el triple por metro cuadrado que los grandes;
antes de decidir habría que agregar el costo de renta por
m\textsuperscript{2}, que no está en los datos.
\end{solucion}

\begin{solucion}{p-3}
La venta diaria promedio es el total de unidades de 2023 entre 365 días. La
cobertura dice cuántos días alcanza el stock actual a ese ritmo.

\begin{rcode}
library(tidyverse)
transacciones <- read_csv("datasets/transacciones.csv",
                          show_col_types = FALSE)
productos <- read_csv("datasets/productos.csv", show_col_types = FALSE)

alerta <- transacciones %>%
  group_by(producto_id) %>%
  summarise(unidades_2023 = sum(cantidad), .groups = "drop") %>%
  right_join(productos, by = "producto_id") %>%
  filter(activo) %>%
  mutate(unidades_2023   = replace_na(unidades_2023, 0),
         venta_diaria    = unidades_2023 / 365,
         dias_cobertura  = stock_actual / venta_diaria,
         bajo_minimo     = stock_actual < stock_minimo,
         alerta          = dias_cobertura < 30 | bajo_minimo) %>%
  filter(alerta) %>%
  arrange(dias_cobertura) %>%
  transmute(producto_id, stock = stock_actual, minimo = stock_minimo,
            diaria = round(venta_diaria, 2),
            cobertura = round(dias_cobertura), bajo_minimo)

alerta

# Contexto: distribución de la cobertura de todos los activos
productos %>%
  filter(activo) %>%
  left_join(count(transacciones, producto_id, wt = cantidad,
                  name = "unidades"), by = "producto_id") %>%
  mutate(dias_cobertura = stock_actual / (unidades / 365)) %>%
  pull(dias_cobertura) %>%
  quantile(c(0, 0.25, 0.5)) %>%
  round()
\end{rcode}

\begin{rsalida}
# A tibble: 3 x 6
  producto_id stock minimo diaria cobertura bajo_minimo
        <dbl> <dbl>  <dbl>  <dbl>     <dbl> <lgl>      
1           3    22     26   0.18       125 TRUE       
2           6    36     38   0.22       160 TRUE       
3          33    42     45   0.21       199 TRUE       
  0%  25%  50% 
 125  927 1717 
\end{rsalida}

\textbf{Interpretación:} ningún producto activo tiene menos de 30 días de
cobertura: el que menos tiene alcanza para 125 días. Las 3 alertas se deben
a que el stock está por debajo del mínimo definido (productos 3, 6 y 33),
aunque al ritmo de venta actual (unas 0.2 unidades diarias) les alcanza para más
de cuatro meses. El hallazgo importante es el contrario: la mediana de
cobertura es de 1,717 días (más de cuatro años de inventario). Con 1,000
transacciones al año, TechRetail parece tener \emph{sobreinventario}, y los
\codigo{stock\_minimo} están desalineados con la demanda real. Recomendación:
recalcular mínimos a partir de la venta diaria (por ejemplo, 30 días de
cobertura) y revisar si el stock se registra en unidades o en cajas.
\end{solucion}

\begin{solucion}{p-4}
La cohorte es el mes de la primera compra. Para cada compra calculamos
cuántos meses pasaron desde ese mes (\codigo{meses\_desde}) y contamos
clientes distintos por cohorte y mes.

\begin{rcode}
library(tidyverse)
transacciones <- read_csv("datasets/transacciones.csv",
                          show_col_types = FALSE)

compras <- transacciones %>%
  mutate(mes_compra = floor_date(fecha, "month")) %>%
  group_by(cliente_id) %>%
  mutate(cohorte = min(mes_compra)) %>%
  ungroup() %>%
  mutate(meses_desde = (year(mes_compra) - year(cohorte)) * 12 +
                        month(mes_compra) - month(cohorte))

# Tamaño de cada cohorte
compras %>% distinct(cliente_id, cohorte) %>% count(cohorte) %>%
  mutate(cohorte = format(cohorte, "%Y-%m"))

matriz <- compras %>%
  distinct(cliente_id, cohorte, meses_desde) %>%
  count(cohorte, meses_desde, name = "clientes") %>%
  group_by(cohorte) %>%
  mutate(pct = round(100 * clientes / clientes[meses_desde == 0])) %>%
  ungroup()

# Primer trimestre: % de la cohorte que compró 1, 2 y 3 meses después
matriz %>%
  filter(cohorte <= as.Date("2023-03-01"), meses_desde <= 3) %>%
  mutate(cohorte = format(cohorte, "%Y-%m")) %>%
  select(-clientes) %>%
  pivot_wider(names_from = meses_desde, values_from = pct,
              names_prefix = "mes_")
\end{rcode}

\begin{rsalida}
# A tibble: 10 x 2
   cohorte     n
   <chr>   <int>
 1 2023-01    80
 2 2023-02    36
 3 2023-03    29
 4 2023-04    16
 5 2023-05    14
 6 2023-06     3
 7 2023-07     9
 8 2023-08     6
 9 2023-09     3
10 2023-12     2
# A tibble: 3 x 5
  cohorte mes_0 mes_1 mes_2 mes_3
  <chr>   <dbl> <dbl> <dbl> <dbl>
1 2023-01   100    39    38    30
2 2023-02   100    36    19    31
3 2023-03   100    34    38    28
\end{rsalida}

\textbf{Interpretación:} 80 de los 198 clientes activos hicieron su primera
compra en enero, y las cohortes se achican mes a mes porque con compras
aleatorias casi todos los clientes aparecen pronto (solo 2 debutan en
diciembre). En las cohortes del primer trimestre, entre 34\% y 39\% vuelve a
comprar al mes siguiente y alrededor de 30\% al tercer mes, sin una caída
clara. Ojo al interpretar: las cohortes tardías tienen muy pocos clientes
(3 en junio o septiembre) y sus porcentajes serían puro ruido; por eso
mostramos solo las del primer trimestre. En una empresa real las cohortes se
definen con fecha de alta confiable y varios años de historia.
\end{solucion}

\begin{solucion}{p-5}
La regla de Tukey marca como atípico lo que queda fuera de
$[Q_1 - 1.5\,\mathrm{IQR},\ Q_3 + 1.5\,\mathrm{IQR}]$. La segunda regla
compara el precio de cada transacción con la mediana del mismo producto.

\begin{rcode}
library(tidyverse)
transacciones <- read_csv("datasets/transacciones.csv",
                          show_col_types = FALSE)

# Regla (a): Tukey sobre el total de la transacción
q <- quantile(transacciones$total_transaccion, c(0.25, 0.75))
iqr <- unname(q[2] - q[1])
limites <- c(inferior = q[[1]] - 1.5 * iqr, superior = q[[2]] + 1.5 * iqr)
round(limites, 2)
sum(transacciones$total_transaccion < limites["inferior"] |
      transacciones$total_transaccion > limites["superior"])

# Regla (b): precio muy alejado de la mediana del mismo producto
revision <- transacciones %>%
  group_by(producto_id) %>%
  mutate(precio_mediano = median(precio_venta),
         razon = precio_venta / precio_mediano) %>%
  ungroup() %>%
  mutate(atipico_precio = razon > 3 | razon < 1 / 3)

sum(revision$atipico_precio)

# Lista para Auditoría: los casos más extremos primero
revision %>%
  filter(atipico_precio) %>%
  mutate(distancia = abs(log(razon))) %>%
  arrange(desc(distancia)) %>%
  transmute(id = transaccion_id, tienda = tienda_id,
            vendedor = vendedor_id, producto = producto_id,
            precio = precio_venta, mediana = precio_mediano,
            razon = round(razon, 2)) %>%
  head(6)

# ¿Se concentran en algún vendedor?
revision %>% filter(atipico_precio) %>%
  count(vendedor_id, sort = TRUE) %>% head(3)
\end{rcode}

\begin{rsalida}
inferior superior 
-1365.16  3184.89 
[1] 0
[1] 148
# A tibble: 6 x 7
     id tienda vendedor producto precio mediana razon
  <dbl>  <dbl>    <dbl>    <dbl>  <dbl>   <dbl> <dbl>
1   106      9       12       45   20.1    434.  0.05
2   930      6        6       43   24.2    393.  0.06
3   185      3       22       16   22.7    299.  0.08
4   909      1       20       42   22.6    293.  0.08
5   745      4       22       24   20.2    260.  0.08
6   925      2       19        7   26.4    335.  0.08
# A tibble: 3 x 2
  vendedor_id     n
        <dbl> <int>
1          23     9
2          24     9
3           5     8
\end{rsalida}

\textbf{Interpretación:} la regla de Tukey no marca \emph{ninguna}
transacción: el total máximo (\$2,992) está por debajo del límite superior
(\$3,185) y el límite inferior es negativo. Esto pasa porque la cantidad
(1 a 5) y el precio (20 a 600) están acotados. En cambio, comparar el precio
con el de su propio producto marca 148 transacciones (15\%); las más extremas se cobraron a entre 5\% y
8\% del precio mediano del producto, lo cual confirma
el hallazgo de la auditoría: el problema no son los montos, sino precios
inconsistentes para un mismo producto. Los vendedores 23 y 24 concentran 9
casos cada uno; antes de sacar conclusiones sobre personas, compáralo con su
número total de transacciones (43 y 44), porque con 148 casos repartidos en
25 vendedores, 9 por vendedor está dentro de lo esperable.
\end{solucion}

\begin{solucion}{p-6}
Primero medimos si los descuentos se asocian con más unidades en
\archivo{ventas\_retail.csv}; después simulamos diciembre con un descuento
de 10\% y un aumento aleatorio de unidades.

\begin{rcode}
library(tidyverse)
ventas_retail <- read_csv("datasets/ventas_retail.csv",
                          show_col_types = FALSE)
transacciones <- read_csv("datasets/transacciones.csv",
                          show_col_types = FALSE)

# 1. ¿Con descuento se venden más unidades?
ventas_retail %>%
  mutate(con_descuento = descuento_pct > 0) %>%
  group_by(con_descuento) %>%
  summarise(ventas = n(), cantidad_promedio = round(mean(cantidad), 2),
            .groups = "drop")
t.test(cantidad ~ descuento_pct > 0, data = ventas_retail)$p.value

# 2. Simulación de diciembre con 10% de descuento
ingreso_base <- transacciones %>%
  filter(month(fecha) == 12) %>%
  summarise(total = sum(total_transaccion)) %>%
  pull(total)

set.seed(123)
simulacion <- tibble(aumento_unidades = runif(1000, 0, 0.15)) %>%
  mutate(ingreso = ingreso_base * (1 + aumento_unidades) * 0.90,
         diferencia = ingreso - ingreso_base)

simulacion %>%
  summarise(ingreso_base = round(ingreso_base),
            ingreso_promedio = round(mean(ingreso)),
            pct_pierde = round(100 * mean(diferencia < 0), 1),
            peor_caso = round(min(diferencia)),
            mejor_caso = round(max(diferencia)))

# 3. Aumento de unidades necesario para no perder: (1 + x) * 0.9 = 1
round(100 * (1 / 0.90 - 1), 1)
\end{rcode}

\begin{rsalida}
# A tibble: 2 x 3
  con_descuento ventas cantidad_promedio
  <lgl>          <int>             <dbl>
1 FALSE            188              5.61
2 TRUE             177              5.67
[1] 0.8404565
# A tibble: 1 x 5
  ingreso_base ingreso_promedio pct_pierde peor_caso mejor_caso
         <dbl>            <dbl>      <dbl>     <dbl>      <dbl>
1        96394            93225       74.5     -9633       3366
[1] 11.1
\end{rsalida}

\textbf{Interpretación:} en \archivo{ventas\_retail.csv} las ventas con
descuento no tienen más unidades (5.67 contra 5.61, $p = 0.84$): no hay
evidencia de que los descuentos muevan volumen. Con un descuento de 10\% se
necesita vender 11.1\% más unidades solo para igualar el ingreso. Si el
aumento está entre 0\% y 15\%, la campaña pierde ingreso en 74.5\% de las
simulaciones; en promedio el ingreso de diciembre bajaría de \$96,394 a
\$93,225, con una pérdida de hasta \$9,633 y una ganancia máxima de apenas
\$3,366. Y esto es ingreso: con márgenes estimados de 46\%, el efecto en
utilidad sería peor. Recomendación: no lanzar un descuento general; probarlo
primero en unas tiendas con un grupo de control.
\end{solucion}

\begin{solucion}{p-7}
Las \termino{funciones de ventana} (\codigo{min\_rank()}, \codigo{lag()},
\codigo{cumsum()}) calculan un valor por fila usando otras filas del mismo
grupo. Primero completamos la tabla vendedor-mes con \codigo{complete()},
para que un mes sin ventas cuente como cero.

\begin{rcode}
library(tidyverse)
transacciones <- read_csv("datasets/transacciones.csv",
                          show_col_types = FALSE)

ranking <- transacciones %>%
  mutate(mes = month(fecha)) %>%
  count(vendedor_id, mes, wt = total_transaccion, name = "ventas") %>%
  complete(vendedor_id, mes = 1:12, fill = list(ventas = 0)) %>%
  group_by(mes) %>%
  mutate(lugar = min_rank(desc(ventas))) %>%          # ranking del mes
  group_by(vendedor_id) %>%
  arrange(mes, .by_group = TRUE) %>%
  mutate(acumulado = cumsum(ventas),                  # venta acumulada
         cambio_lugar = lag(lugar) - lugar) %>%       # + = subió
  ungroup()

# Vendedores con más meses en el top 5
ranking %>%
  group_by(vendedor_id) %>%
  summarise(meses_top5 = sum(lugar <= 5),
            ventas_anuales = round(max(acumulado)), .groups = "drop") %>%
  arrange(desc(meses_top5), desc(ventas_anuales)) %>%
  head(5)

# Mayor subida de un mes a otro
ranking %>%
  slice_max(cambio_lugar, n = 1) %>%
  select(vendedor_id, mes, lugar, cambio_lugar, ventas) %>%
  mutate(ventas = round(ventas))
\end{rcode}

\begin{rsalida}
# A tibble: 5 x 3
  vendedor_id meses_top5 ventas_anuales
        <dbl>      <int>          <dbl>
1           7          6          50470
2          11          4          50017
3           3          4          48784
4           6          4          47946
5          15          4          42709
# A tibble: 3 x 5
  vendedor_id   mes lugar cambio_lugar ventas
        <dbl> <dbl> <int>        <int>  <dbl>
1           1    10     3           21   5357
2           2     5     3           21   8119
3          17     3     2           21   9971
\end{rsalida}

\textbf{Interpretación:} el vendedor 7, que fue el primero del año
(\$50,470), estuvo en el top 5 en 6 de los 12 meses; ningún otro supera 4
meses. La mayor subida fue de 21 lugares y la lograron tres vendedores (el 1
en octubre, el 2 en mayo y el 17 en marzo): pasar del lugar 23 o 24 al top 3
de un mes a otro indica que el ranking mensual es muy volátil con pocas
transacciones por vendedor (unas 3 o 4 al mes). Para evaluar desempeño
conviene un ranking trimestral o un promedio móvil, no el mes aislado.
\end{solucion}

\begin{solucion}{p-8}
La función arma las tres hojas como una lista con nombre y la escribe con
\codigo{write\_xlsx()} de \paquete{writexl} (cada elemento de la lista es una
hoja). \codigo{walk()} la aplica a las 10 tiendas.

\begin{rcode}
library(tidyverse)
library(writexl)
transacciones <- read_csv("datasets/transacciones.csv",
                          show_col_types = FALSE)
tiendas   <- read_csv("datasets/tiendas.csv", show_col_types = FALSE)
productos <- read_csv("datasets/productos.csv", show_col_types = FALSE)

dir.create("resultados/gerentes", recursive = TRUE, showWarnings = FALSE)

# KPIs de un conjunto de transacciones
kpis_de <- function(datos) {
  datos %>% summarise(ventas = round(sum(total_transaccion)),
                      transacciones = n(),
                      ticket = round(mean(total_transaccion), 2),
                      clientes = n_distinct(cliente_id))
}

# Promedio de la cadena por tienda (para comparar)
cadena <- transacciones %>%
  group_by(tienda_id) %>%
  group_modify(~ kpis_de(.x)) %>%
  ungroup() %>%
  summarise(across(-tienda_id, mean)) %>%
  mutate(ventas = round(ventas), ticket = round(ticket, 2),
         clientes = round(clientes, 1))

paquete_gerente <- function(id) {
  info  <- filter(tiendas, tienda_id == id)
  datos <- filter(transacciones, tienda_id == id)

  resumen <- bind_rows(
    kpis_de(datos) %>% mutate(nivel = info$nombre_tienda),
    cadena %>% mutate(nivel = "Promedio cadena")) %>%
    relocate(nivel)

  mensual <- datos %>%
    group_by(mes) %>%
    summarise(ventas = round(sum(total_transaccion)),
              transacciones = n(), .groups = "drop") %>%
    mutate(crecimiento_pct = round(100 * (ventas / lag(ventas) - 1), 1))

  top <- datos %>%
    group_by(producto_id) %>%
    summarise(unidades = sum(cantidad),
              ventas = round(sum(total_transaccion)), .groups = "drop") %>%
    left_join(select(productos, producto_id, nombre_producto, categoria),
              by = "producto_id") %>%
    slice_max(ventas, n = 10, with_ties = FALSE)

  archivo <- file.path("resultados/gerentes",
                       str_c("tienda_", sprintf("%02d", id), "_",
                             str_replace_all(info$nombre_tienda, " ", "_"),
                             ".xlsx"))
  write_xlsx(list(Resumen = resumen, Mensual = mensual,
                  `Top productos` = top), archivo)
}

walk(tiendas$tienda_id, paquete_gerente)

archivos <- list.files("resultados/gerentes", pattern = "xlsx$")
length(archivos)
head(archivos, 3)

# Comprobación: leer la hoja Resumen de la tienda 8
readxl::read_excel(file.path("resultados/gerentes", archivos[8]),
                   sheet = "Resumen")
\end{rcode}

\begin{rsalida}
[1] 10
[1] "tienda_01_Sucursal_Centro.xlsx" "tienda_02_Sucursal_Norte.xlsx" 
[3] "tienda_03_Sucursal_Sur.xlsx"   
# A tibble: 2 x 5
  nivel           ventas transacciones ticket clientes
  <chr>            <dbl>         <dbl>  <dbl>    <dbl>
1 Sucursal Outlet 114377           106  1079.     82  
2 Promedio cadena  95999           100   959.     79.5
\end{rsalida}

\textbf{Interpretación:} se generaron los 10 archivos, uno por sucursal,
con nombre ordenable (\archivo{tienda\_01\_\ldots}). La hoja
\emph{Resumen} de la Sucursal Outlet muestra que vendió \$114,377 contra un
promedio de \$95,999 por tienda, con un ticket de \$1,079 contra \$959. Ese
formato de ``mi tienda contra la cadena'' es justo lo que un gerente
necesita. Para automatizarlo cada mes, combina esta función con la
programación de tareas del Módulo~\ref{cap:reportes}. Si quieres dar formato
(colores, anchos de columna) usa \paquete{openxlsx}, que no está instalado
en este entorno.
\end{solucion}

